% Force des acides et des bases, couples acide/base — Chimie Tle TI — Cours signé OnBuch+
% Programme officiel MINESEC, Terminale TI. Compiler avec : tectonic L07-force-des-acides-et-bases-couples.tex
\newcommand{\DOCMATIERE}{Chimie}
\newcommand{\DOCNIVEAU}{Tle TI}
\newcommand{\DOCMODULE}{Module 2 · Acides et bases}
\newcommand{\DOCLECON}{07}
\newcommand{\DOCTITRE}{Force des acides et des bases, couples acide/base}
% OnBuch+ — préambule « cours signés » ENRICHI (compilé avec tectonic / XeLaTeX)
% Système visuel « Pop » d'OnBuch+ : fond crème (jamais blanc), bordures encre
% épaisses, ombres « dures » décalées, titres Archivo Black en capitales.
% Version enrichie pour les cours de chimie : chimie (mhchem, chemfig),
% graphes (pgfplots), boîtes pédagogiques supplémentaires (définition,
% propriété, méthode, attention, le savais-tu, exercice résolu, exercice,
% TP, bilan), page de garde refaite et signature OnBuch+ ancrée en bas.
%
% Macros attendues dans main.tex AVANT \input{preamble} :
%   \DOCMATIERE  \DOCNIVEAU  \DOCMODULE  \DOCLECON (numéro)  \DOCTITRE
\documentclass[11pt]{article}

\usepackage{fontspec}
\usepackage[french]{babel}
\usepackage[a4paper,margin=2cm,top=3.4cm,bottom=2.8cm,headheight=1.4cm,footskip=1.3cm]{geometry}
\usepackage{amsmath,amssymb,amsfonts}
\usepackage[table]{xcolor} % [table] : fournit aussi \rowcolors (alternance de lignes)
\usepackage{pagecolor}
\usepackage{tikz}
\usetikzlibrary{arrows.meta,positioning,calc,shapes.geometric,shapes.misc,shapes.arrows,decorations.pathmorphing,decorations.markings,patterns,shadows,mindmap,trees,fit,backgrounds,matrix,intersections}
\usepackage{pgfplots}
\pgfplotsset{compat=1.18}
\usepackage[version=4]{mhchem}
\usepackage{chemfig}
\usepackage{enumitem}
\usepackage{fancyhdr}
% [most] charge toutes les bibliothèques tcolorbox (listings, minted, raster,
% documentation...) et gonfle inutilement la mémoire TeX sur un document long
% (~300 boîtes) : on ne charge que ce qui est réellement utilisé.
\usepackage[skins,breakable]{tcolorbox}
\usepackage{graphicx}
\usepackage{colortbl}
\usepackage{booktabs}
\usepackage{tabularx}
\usepackage{array}
\usepackage{multicol}
\usepackage{siunitx}
% parse-numbers=false : accepte aussi \SI{2,5 \times 10^{-3}}{...}
\sisetup{locale=FR,output-decimal-marker={,},group-minimum-digits=4,parse-numbers=false}
\DeclareSIUnit\ohm{\ensuremath{\symup{\Omega}}} % U+2126 absent de la police
\usepackage{qrcode}
% \degree : commande fréquemment utilisée par les modèles pour un angle ou une
% température en dehors de \SI{}{} (non fournie par les paquets chargés).
\providecommand{\degree}{^{\circ}}
\usepackage{lastpage}
\usepackage{hyperref}
\hypersetup{hidelinks}

% ============================================================
% Palette « Pop » OnBuch+ — mêmes hex que la classe P (plus_ui.dart)
% ============================================================
\definecolor{poporange}{HTML}{F59321}
\definecolor{popdark}{HTML}{E07A0C}
\definecolor{popcream}{HTML}{FFF6E8}
\definecolor{popink}{HTML}{1C1714}
\definecolor{poppurple}{HTML}{7A5AE0}
\definecolor{popgreen}{HTML}{2E9E6A}
\definecolor{poppink}{HTML}{E0457B}
\definecolor{popblue}{HTML}{2F7DE1}
\definecolor{popgold}{HTML}{C98A12}
\definecolor{popmuted}{HTML}{8A7F76}
\definecolor{popline}{HTML}{EADFD2}
\definecolor{popgray}{HTML}{8A7F76}
\colorlet{popgrey}{popgray}
% Alias tolérants (le modèle nomme parfois les couleurs différemment)
\colorlet{popwhite}{white}
\colorlet{popblack}{popink}
\colorlet{popred}{poppink}
\colorlet{popyellow}{popgold}
% Teintes claires pour fonds de tableaux / zones
\colorlet{popblueL}{popblue!10!white}
\colorlet{poporangeL}{poporange!14!white}
\colorlet{popgreenL}{popgreen!12!white}
\colorlet{poppurpleL}{poppurple!12!white}
\colorlet{poppinkL}{poppink!12!white}
\colorlet{popgoldL}{popgold!14!white}

\pagecolor{popcream}

% ============================================================
% Typographie
% ============================================================
\defaultfontfeatures{Ligatures=TeX}
\setmainfont{PlusJakartaSans-Medium.ttf}[Path=../../../fonts/,BoldFont=PlusJakartaSans-Bold.ttf]
% Maths en sans-serif (Fira Math) avec chiffres et lettres droites de Plus
% Jakarta, pour que les formules se fondent dans le texte.
\usepackage[math-style=ISO,bold-style=ISO]{unicode-math}
\setmathfont{FiraMath-Regular.otf}[Path=../../../fonts/]
\setmathfont{PlusJakartaSans-Medium.ttf}[Path=../../../fonts/,range={up/num,up/{Latin,latin},\mathup}]
\usepackage{icomma} % virgule décimale sans espace en mode math
% Caractères absents de la police de texte (sinon carré vide) : rendus par
% leur équivalent mathématique, en texte comme en formule.
\usepackage{newunicodechar}
\newunicodechar{α}{\ensuremath{\alpha}}
\newunicodechar{Α}{\ensuremath{\alpha}}
\newunicodechar{β}{\ensuremath{\beta}}
\newunicodechar{δ}{\ensuremath{\delta}}
\newunicodechar{✓}{\popcheck{popgreen}}
\newfontfamily\popdisplay{ArchivoBlack-Regular.ttf}[Path=../../../fonts/]
\newfontfamily\poplogo{SpaceGrotesk-Bold.ttf}[Path=../../../fonts/]
\newfontfamily\popsemi{PlusJakartaSans-SemiBold.ttf}[Path=../../../fonts/]
\newfontfamily\popxbold{PlusJakartaSans-ExtraBold.ttf}[Path=../../../fonts/]
\newfontfamily\popxboldls{PlusJakartaSans-ExtraBold.ttf}[Path=../../../fonts/,LetterSpace=3.0]

\setlist[enumerate]{leftmargin=1.7em,itemsep=5pt,topsep=4pt}
\setlist[itemize]{leftmargin=1.7em,itemsep=4pt,topsep=4pt,label={\color{poporange}\textbullet}}
\setlength{\parindent}{0pt}
\setlength{\parskip}{5pt}
\renewcommand{\arraystretch}{1.25}

% chemfig : liaisons un peu plus épaisses, encre Pop
\setchemfig{atom sep=2em,bond style={line width=0.9pt,popink},cram width=3pt}

% pgfplots : style Pop par défaut
\pgfplotsset{
  every axis/.append style={axis line style={popink,line width=1pt},tick style={popink},
    grid style={popline,line width=0.5pt},label style={font=\small},tick label style={font=\footnotesize}},
  popcurve/.style={poporange,line width=1.8pt,no markers,smooth},
}

% ============================================================
% Pill / squiggle
% ============================================================
\newcommand{\poppill}[3]{%
  \tikz[baseline=-0.55ex]\node[draw=popink,line width=1.2pt,rounded corners=2.6pt,%
    fill=#2,inner xsep=6.5pt,inner ysep=3pt]{%
    \popxboldls\fontsize{7}{7}\selectfont\color{#3}\MakeUppercase{#1}};%
}
\newcommand{\popsquiggle}[1]{%
  \tikz[baseline=-0.4ex]{
    \draw[#1,line width=2.6pt,line cap=round]
      plot [smooth] coordinates {(0,0.09) (0.34,0.01) (0.68,0.21) (1.02,0.01) (1.36,0.10)};
  }%
}
% Mot-clé surligné (marqueur orange sous le texte)
\newcommand{\cle}[1]{{\popxbold\color{popdark}#1}}

% ============================================================
% En-tête / pied
% ============================================================
\pagestyle{fancy}
\fancyhf{}
\renewcommand{\headrulewidth}{0pt}
\renewcommand{\footrulewidth}{0pt}
\lhead{\raisebox{-0.15ex}{\poplogo\fontsize{13}{13}\selectfont\color{popink}OnBuch\color{poporange}+}}
\rhead{\poppill{\DOCMATIERE\ \textbullet\ \DOCNIVEAU\ \textbullet\ Leçon \DOCLECON}{white}{popink}}
\lfoot{\popxbold\fontsize{7.6}{7.6}\selectfont\color{popmuted}\MakeUppercase{Cours signé OnBuch+}\ \textbullet\ {\popsemi\DOCTITRE}}
\rfoot{\tikz[baseline=-0.6ex]\node[rounded corners=8.5pt,fill=popink,minimum height=17pt,minimum width=17pt,inner xsep=5pt,inner sep=0]{\popxbold\fontsize{8}{8}\selectfont\color{popcream}\thepage\,/\,\pageref*{LastPage}};}

% ============================================================
% Titres
% ============================================================
\newcommand{\chaptitre}[2]{%
  \begin{tcolorbox}[enhanced,colback=poporange,colframe=popink,boxrule=2.6pt,arc=11pt,
      drop shadow=popink,left=15pt,right=15pt,top=12pt,bottom=11pt,before skip=3pt,after skip=18pt]
    \poppill{#2}{popink}{popcream}\par\vspace{9pt}
    {\raggedright\hyphenpenalty=10000\exhyphenpenalty=10000\popdisplay\fontsize{22}{24}\selectfont\color{popcream}\MakeUppercase{#1}\par}
  \end{tcolorbox}
}
% Section numérotée : pastille orange avec le numéro + titre Archivo
\newcounter{popsec}
\newcommand{\coursec}[1]{%
  \stepcounter{popsec}\setcounter{popsub}{0}%
  \par\vspace{16pt}\noindent
  \begin{minipage}{\linewidth}
  \tikz[baseline=-0.7ex]\node[draw=popink,line width=1.6pt,fill=poporange,rounded corners=4pt,minimum size=22pt,inner sep=0]{\popdisplay\fontsize{12}{12}\selectfont\color{popcream}\thepopsec};%
  \hspace{8pt}{\popdisplay\fontsize{15}{17}\selectfont\color{popink}\MakeUppercase{#1}}\par
  \vspace{4pt}\noindent{\color{poporange}\rule{36pt}{3.6pt}}
  \end{minipage}\par\nopagebreak\vspace{8pt}
}
\newcounter{popsub}
\newcommand{\courssub}[1]{%
  \stepcounter{popsub}%
  \par\vspace{9pt}\noindent{\popxbold\fontsize{11.5}{13}\selectfont\color{popink}{\color{poporange}\thepopsec.\thepopsub}\hspace{6pt}#1}\par\nopagebreak\vspace{4pt}
}

% ============================================================
% Boîtes pédagogiques — même recette : fond blanc, bordure épaisse colorée,
% ombre dure encre, badge plein. Titre optionnel [#1].
% ============================================================
\tcbset{popbase/.style={breakable,enhanced,colback=white,boxrule=2.3pt,arc=9pt,
  shadow={3pt}{-3pt}{0pt}{popink},left=14pt,right=14pt,top=11pt,bottom=11pt,before skip=11pt,after skip=15pt,
  fontupper=\popsemi\fontsize{10.6}{14.5}\selectfont\color{popink},
  fontlower=\popsemi\fontsize{10.6}{14.5}\selectfont\color{popink}}}
% Coche dessinée (le glyphe ✓ n'existe pas dans les polices du document)
\newcommand{\popcheck}[1]{\tikz[baseline=-0.5ex]\draw[#1,line width=1.6pt,line cap=round,line join=round](-0.13,0.0)--(-0.04,-0.09)--(0.13,0.1);}
\newcommand{\popboxhead}[3]{\poppill{#1}{#2}{white}\ifx\relax#3\relax\else\hspace{6pt}{\popxbold\fontsize{10.6}{12}\selectfont\color{popink}#3}\fi\par\vspace{7pt}}

\newtcolorbox{aretenir}[1][]{popbase,colframe=popgreen,before upper={\popboxhead{À retenir}{popgreen}{#1}}}
\newtcolorbox{exemplebox}[1][]{popbase,colframe=poporange,before upper={\popboxhead{Exemple}{poporange}{#1}}}
\newtcolorbox{definition}[1][]{popbase,colframe=popblue,before upper={\popboxhead{Définition}{popblue}{#1}}}
\newtcolorbox{propriete}[1][]{popbase,colframe=poppurple,before upper={\popboxhead{Propriété}{poppurple}{#1}}}
\newtcolorbox{methode}[1][]{popbase,colframe=popgold,before upper={\popboxhead{Méthode}{popgold}{#1}}}
\newtcolorbox{attention}[1][]{popbase,colframe=poppink,before upper={\popboxhead{Attention}{poppink}{#1}}}
\newtcolorbox{experience}[1][]{popbase,colframe=popblue,colback=popblueL,before upper={\popboxhead{Expérience}{popblue}{#1}}}
% « Le savais-tu ? » : carte violette pleine (contexte, culture, Cameroun)
\newtcolorbox{savaistu}[1][]{popbase,colback=poppurple,colframe=popink,
  fontupper=\popsemi\fontsize{10.4}{14.2}\selectfont\color{white},
  before upper={\poppill{Le savais-tu\,?}{white}{poppurple}\ifx\relax#1\relax\else\hspace{6pt}{\popxbold\color{white}#1}\fi\par\vspace{7pt}}}
% Exercice résolu : énoncé en haut, \tcblower puis la solution sur fond crème
\newcounter{popexo}\newcounter{popexr}
\newtcolorbox{exoresolu}[1][]{popbase,colframe=poporange,colbacklower=popcream,
  segmentation style={popink,line width=1.2pt,dashed},
  before upper={\stepcounter{popexr}\popboxhead{Exercice résolu \thepopexr}{poporange}{#1}},
  before lower={\poppill{Solution}{popink}{popcream}\par\vspace{6pt}}}
% Exercice (non résolu en place) — la correction est dans la partie Corrigés
\newtcolorbox{exercice}[1][]{popbase,colframe=popink,
  before upper={\stepcounter{popexo}\popboxhead{Exercice \thepopexo}{popink}{#1}}}
% Corrigé
\newtcolorbox{corrige}[1][]{popbase,colframe=popgreen,colback=popgreenL,
  before upper={\popboxhead{Corrigé}{popgreen}{#1}}}
% Objectifs / prérequis / bilan
\newtcolorbox{objectifs}{popbase,colframe=popink,colback=poporangeL,
  before upper={\popboxhead{Objectifs de la leçon}{popink}{}}}
\newtcolorbox{prerequis}{popbase,colframe=popink,colback=popblueL,
  before upper={\popboxhead{Prérequis}{popblue}{}}}
\newtcolorbox{bilan}{popbase,colframe=popink,colback=white,boxrule=2.8pt,
  before upper={\popboxhead{Fiche bilan}{poporange}{L'essentiel à connaître pour l'examen}}}
% Figure encadrée avec légende
\newtcolorbox{popfigure}[1][]{enhanced,breakable=false,colback=white,colframe=popline,boxrule=1.4pt,arc=8pt,
  left=10pt,right=10pt,top=10pt,bottom=8pt,before skip=10pt,after skip=12pt,halign=center,
  lower separated=false,halign lower=center,
  fontlower=\popsemi\fontsize{9}{11.5}\selectfont\color{popmuted},
  #1}
\newcommand{\legende}[1]{\par\vspace{6pt}{\popsemi\fontsize{9}{11.5}\selectfont\color{popmuted}#1}}

% ============================================================
% Signature OnBuch+
% ============================================================
\newcommand{\poplogoinline}[1]{{\poplogo\fontsize{#1}{#1}\selectfont\color{popink}OnBuch\color{poporange}+}}
\newcommand{\signatureonbuch}{%
  \begin{tcolorbox}[enhanced,colback=popink,colframe=popink,boxrule=2.2pt,arc=10pt,
      drop shadow=poporange,left=14pt,right=14pt,top=10pt,bottom=10pt,before skip=0pt,after skip=0pt]
    \tikz[baseline=-0.7ex]\node[circle,fill=popgreen,draw=popcream,line width=1.3pt,minimum size=22pt,inner sep=0]{\popcheck{white}};%
    \hspace{9pt}\begin{minipage}[c]{0.62\linewidth}
      {\popxbold\fontsize{12}{13}\selectfont\color{popcream}Signé\ \textbullet\ {\poplogo OnBuch\color{poporange}+}}\par\vspace{2pt}
      {\raggedright\popsemi\fontsize{8}{10.5}\selectfont\color{popline}\MakeUppercase{Équipe pédagogique OnBuch+}\\\MakeUppercase{Conforme au programme officiel MINESEC}\par}
    \end{minipage}\hfill
    {\poplogo\fontsize{22}{22}\selectfont\color{popcream}OnBuch\color{poporange}+}
  \end{tcolorbox}%
}

% ============================================================
% Page de garde
% #1 titre · #2 matière · #3 niveau · #4 module · #5 n° leçon · #6 durée ·
% #7 description / objectifs (texte ou itemize)
% ============================================================
\newcommand{\pagegarde}[7]{%
  \thispagestyle{empty}
  \newgeometry{a4paper,margin=1.7cm,top=1.6cm,bottom=1.4cm}%
  \begin{tikzpicture}[remember picture,overlay]
    \fill[poporange!18] ([xshift=-3.2cm,yshift=-3.6cm]current page.north east) circle (4.6cm);
    \fill[poppurple!14] ([xshift=2.4cm,yshift=7.6cm]current page.south west) circle (3.8cm);
  \end{tikzpicture}%
  {\poplogo\fontsize{30}{30}\selectfont\color{popink}OnBuch\color{poporange}+}\hfill
  \raisebox{0.6ex}{\poppill{Cours signé \textbullet\ Premium}{popink}{popcream}}\par
  \vspace{1pt}\popsquiggle{poppurple}\par
  \vspace{8pt}
  \begin{tcolorbox}[enhanced,colback=poporange,colframe=popink,boxrule=2.8pt,arc=13pt,
      drop shadow=popink,left=18pt,right=18pt,top=12pt,bottom=13pt,after skip=10pt]
    \poppill{#2 \textbullet\ #3}{popink}{popcream}\hspace{5pt}\poppill{#4}{white}{popink}\par\vspace{8pt}
    {\popdisplay\fontsize{14}{14}\selectfont\color{popink}LEÇON #5}\par\vspace{4pt}
    {\raggedright\hyphenpenalty=10000\exhyphenpenalty=10000\popdisplay\fontsize{25}{27}\selectfont\color{popcream}\MakeUppercase{#1}\par}
    \vspace{8pt}
    {\popxbold\fontsize{9.5}{9.5}\selectfont\color{popink}\MakeUppercase{Durée conseillée : #6}}
  \end{tcolorbox}
  \begin{tcolorbox}[enhanced,colback=white,colframe=popink,boxrule=2.2pt,arc=10pt,
      drop shadow=popink,left=16pt,right=16pt,top=10pt,bottom=10pt,after skip=8pt]
    \poppill{Dans cette leçon}{poppurple}{white}\par\vspace{5pt}
    \popsemi\fontsize{9.3}{12.6}\selectfont\color{popink}#7
  \end{tcolorbox}
  \noindent\begin{minipage}[t]{0.487\linewidth}
    \begin{tcolorbox}[enhanced,colback=popgreen,colframe=popink,boxrule=2.2pt,arc=10pt,
        drop shadow=popink,left=11pt,right=11pt,top=10pt,bottom=10pt,height=3.1cm,valign=center]
      \tcbox[colback=white,colframe=popink,boxrule=1.3pt,arc=5pt,boxsep=0pt,nobeforeafter,
        left=3pt,right=3pt,top=3pt,bottom=3pt,valign=center]{\qrcode[height=1.9cm]{https://play.google.com/store/apps/details?id=com.luvvix.onbuch&hl=fr}}%
      \hspace{8pt}\begin{minipage}[c]{0.5\linewidth}
        {\popdisplay\fontsize{11}{12.5}\selectfont\color{white}\MakeUppercase{Télécharge OnBuch}}\par\vspace{4pt}
        {\raggedright\popsemi\fontsize{8.4}{11}\selectfont\color{white}Gratuit \textbullet\ Google Play\\Tuteur IA, annales, concours, résultats\par}
      \end{minipage}
    \end{tcolorbox}
  \end{minipage}\hfill
  \begin{minipage}[t]{0.487\linewidth}
    \begin{tcolorbox}[enhanced,colback=poppurple,colframe=popink,boxrule=2.2pt,arc=10pt,
        drop shadow=popink,left=11pt,right=11pt,top=10pt,bottom=10pt,height=3.1cm,valign=center]
      \tikz[baseline=-0.7ex]\node[circle,fill=white,draw=popink,line width=1.3pt,minimum size=1.6cm,inner sep=0]{\popdisplay\fontsize{24}{24}\selectfont\color{poppurple}+};%
      \hspace{8pt}\begin{minipage}[c]{0.6\linewidth}
        {\popdisplay\fontsize{11}{12.5}\selectfont\color{white}\MakeUppercase{Passe à OnBuch+}}\par\vspace{4pt}
        {\raggedright\popsemi\fontsize{8.4}{11}\selectfont\color{white}Tous les cours signés, Tuteur IA illimité, fiches PDF — onglet OnBuch+ de l'app\par}
      \end{minipage}
    \end{tcolorbox}
  \end{minipage}
  \vspace{8pt}
  \popcontenu
  \vfill
  \signatureonbuch
  \restoregeometry
  \newpage
}

% Rangée « Ce cours contient » (page de garde)
\newcommand{\poptile}[3]{% #1 couleur #2 gros texte #3 légende
  \begin{tcolorbox}[enhanced,colback=#1,colframe=popink,boxrule=2pt,arc=9pt,shadow={2.5pt}{-2.5pt}{0pt}{popink},
      left=6pt,right=6pt,top=9pt,bottom=9pt,halign=center,valign=center,height=1.75cm,width=\linewidth,nobeforeafter]
    {\popdisplay\fontsize{12.5}{13}\selectfont\color{white}\MakeUppercase{#2}}\par\vspace{3pt}
    {\centering\popsemi\fontsize{7.8}{9.5}\selectfont\color{white}#3\par}
  \end{tcolorbox}}
\newcommand{\popcontenu}{%
  \par\noindent{\popxboldls\fontsize{7.5}{7.5}\selectfont\color{popmuted}CE COURS SIGNÉ CONTIENT}\par\vspace{7pt}
  \noindent\begin{minipage}[t]{0.235\linewidth}\poptile{popblue}{Cours}{complet, illustré, pas à pas}\end{minipage}\hfill
  \begin{minipage}[t]{0.235\linewidth}\poptile{poporange}{Méthodes}{et exercices résolus}\end{minipage}\hfill
  \begin{minipage}[t]{0.235\linewidth}\poptile{poppink}{Exercices}{type examen + corrigés}\end{minipage}\hfill
  \begin{minipage}[t]{0.235\linewidth}\poptile{popgreen}{Fiche bilan}{l'essentiel à retenir}\end{minipage}\par}

% Sommaire
\newtcolorbox{sommaire}{enhanced,colback=white,colframe=popink,boxrule=2.2pt,arc=10pt,
  drop shadow=popink,left=18pt,right=18pt,top=13pt,bottom=13pt,before skip=6pt,after skip=20pt,
  before upper={\poppill{Sommaire}{poppurple}{white}\par\vspace{9pt}\popsemi\fontsize{10.6}{14.5}\selectfont\color{popink}}}

% Signature de fin de cours — poussée tout en bas de la dernière page
\newcommand{\findecours}{%
  \par\vfill\nopagebreak
  \begin{center}\popsquiggle{poporange}\end{center}\vspace{4pt}
  \signatureonbuch
}
% Glyphes absents de Fira Math : redéfinis à partir de glyphes disponibles
\AtBeginDocument{%
  \def\setminus{\mathbin{\backslash}}\let\smallsetminus\setminus
  \def\vdots{\mathord{\vbox{\baselineskip3.2pt\lineskiplimit0pt\kern4pt\hbox{.}\hbox{.}\hbox{.}}}}%
  \def\ddots{\mathinner{\mkern1mu\raise6.5pt\hbox{.}\mkern2mu\raise3.6pt\hbox{.}\mkern2mu\raise0.7pt\hbox{.}\mkern1mu}}%
  \def\triangle{\mathord{\scalebox{0.9}{$\symup{\Delta}$}}}%
  \def\nabla{\mathord{\raisebox{\depth}{\scalebox{1}[-1]{$\symup{\Delta}$}}}}%
  \def\checkmark{\popcheck{popdark}}%
  \def\star{\mathbin{\ast}}\def\bigstar{\mathord{\ast}}%
  \def\diamond{\mathbin{\scalebox{0.6}{$\Diamond$}}}%
  \def\bigcup{\mathop{\mathchoice{\raisebox{-0.3em}{\scalebox{1.7}{$\cup$}}}{\scalebox{1.25}{$\cup$}}{\cup}{\cup}}\displaylimits}%
  \def\bigcap{\mathop{\mathchoice{\raisebox{-0.3em}{\scalebox{1.7}{$\cap$}}}{\scalebox{1.25}{$\cap$}}{\cap}{\cap}}\displaylimits}%
  \def\lesseqqgtr{\mathrel{\lessgtr}}\def\bigoplus{\mathop{\oplus}\displaylimits}%
}
\providecommand{\ssi}{si et seulement si} % abréviation fréquente

%% FIGFIX-BEGIN v3 — correctifs globaux des figures (docs/onbuch-generation/tools/figfix)
\usepackage{adjustbox}\usepackage{etoolbox}
% toute figure TikZ/pgfplots plus large que la ligne est réduite à la largeur disponible
\BeforeBeginEnvironment{tikzpicture}{\begin{adjustbox}{max width=\linewidth}}
\AfterEndEnvironment{tikzpicture}{\end{adjustbox}}
% légendes pgfplots lisibles par-dessus les courbes
\pgfplotsset{every axis/.append style={legend style={fill=white,fill opacity=0.92,text opacity=1,draw=black!15,inner sep=3pt}}}
% étiquettes d'arêtes (syntaxe edge["..."]) avec fond blanc
\tikzset{every edge quotes/.append style={fill=white,inner sep=1.2pt}}
% bulles de cartes mentales : texte à la ligne dans la bulle, bulle un peu plus grande
\tikzset{every concept/.append style={text width=2.0cm,align=center,minimum size=2.3cm,inner sep=2pt}}
%% FIGFIX-END
\begin{document}

\pagegarde{Force des acides et des bases, couples acide/base}{Chimie}{Tle TI}{Module 2 · Acides et bases}{07}{6 h (cours + TP)}{Cette leçon établit la distinction entre acides et bases forts (totalement ionisés) et faibles (partiellement ionisés) à partir de mesures de pH, puis introduit la constante d'acidité Ka et le pKa pour classer les couples acide/base et prévoir les domaines de prédominance des espèces. Elle culmine par l'étude expérimentale d'un acide faible et la détermination du degré d'acidité d'un vinaigre commercial.
\vspace{4pt}
\begin{multicols}{2}\small
\begin{itemize}
\item À la fin de cette leçon, je sais montrer qu'un acide fort ou une base forte est totalement ionisé(e) en solution à partir du pH et de la concentration
\item À la fin de cette leçon, je sais montrer qu'un acide faible ou une base faible n'est pas totalement ionisé(e) en solution
\item À la fin de cette leçon, je sais utiliser le tableau d'avancement pour étudier un équilibre chimique et calculer un taux d'avancement final
\item À la fin de cette leçon, je sais identifier les couples acide/base mis en jeu dans une équation d'équilibre chimique
\item À la fin de cette leçon, je sais définir et calculer la constante d'acidité Ka, la constante de basicité Kb et le pKa d'un couple
\item À la fin de cette leçon, je sais classer des couples acide/base selon leur Ka ou leur pKa et prévoir le sens d'une réaction acido-basique
\end{itemize}\end{multicols}}

\begin{sommaire}
\begin{multicols}{2}
\begin{enumerate}
\item Acides forts et bases fortes
\item Acides faibles et bases faibles
\item Étude quantitative d'un équilibre : avancement et taux d'avancement
\item Constante d'acidité Ka, constante de basicité Kb et pKa
\item Classification des couples et prévision des réactions
\item Relation pH–pKa et domaines de prédominance
\item Applications : détermination expérimentale de Ka et degré d'acidité du vinaigre
\item Travaux pratiques
\item Méthodes et exercices résolus
\item Exercices
\item Corrigés des exercices
\item Fiche bilan
\end{enumerate}
\end{multicols}
\end{sommaire}

\begin{objectifs}
\begin{itemize}
\item À la fin de cette leçon, je sais montrer qu'un acide fort ou une base forte est totalement ionisé(e) en solution à partir du pH et de la concentration
\item À la fin de cette leçon, je sais montrer qu'un acide faible ou une base faible n'est pas totalement ionisé(e) en solution
\item À la fin de cette leçon, je sais utiliser le tableau d'avancement pour étudier un équilibre chimique et calculer un taux d'avancement final
\item À la fin de cette leçon, je sais identifier les couples acide/base mis en jeu dans une équation d'équilibre chimique
\item À la fin de cette leçon, je sais définir et calculer la constante d'acidité Ka, la constante de basicité Kb et le pKa d'un couple
\item À la fin de cette leçon, je sais classer des couples acide/base selon leur Ka ou leur pKa et prévoir le sens d'une réaction acido-basique
\item À la fin de cette leçon, je sais établir la relation pH = pKa + log([base]/[acide]) et déterminer les domaines de prédominance
\item À la fin de cette leçon, je sais déterminer expérimentalement la constante d'acidité d'un couple et le coefficient d'ionisation d'une espèce
\end{itemize}
\end{objectifs}

\begin{prerequis}
\begin{itemize}
\item Définition du pH et relation pH = −log[H₃O⁺] (Première / début de Terminale)
\item Couples acide/base selon Brønsted et réactions acido-basiques (leçon précédente)
\item Produit ionique de l'eau Ke = 10⁻¹⁴ à 25 °C et relation pH + pOH = 14
\item Tableau d'avancement d'une réaction chimique et notion d'avancement maximal
\item Utilisation de la fonction logarithme décimal et de la calculatrice scientifique
\end{itemize}
\end{prerequis}

\newpage

\coursec{Acides forts et bases fortes}

Au marché Mokolo de Yaoundé, une vendeuse arrose ses légumes de vinaigre pour les conserver : l'acidité du vinaigre empêche le développement des microbes. Quelques kilomètres plus loin, au laboratoire du lycée, le technicien débouche un robinet entartré avec de l'acide chlorhydrique dilué, en prenant soin de porter des gants. Pourtant, si l'on préparait deux solutions de \textbf{même concentration}, l'une d'acide éthanoïque (l'acide du vinaigre), l'autre d'acide chlorhydrique, le pH-mètre afficherait des valeurs très différentes : environ $2,9$ pour le vinaigre dilué, et $2,0$ pour l'acide chlorhydrique. À concentration égale, l'un est nettement plus « agressif » que l'autre. Pourquoi ? Toutes les molécules d'acide ne se transforment-elles pas de la même façon dans l'eau ?

C'est toute la problématique de cette leçon : \emph{comment classer les acides et les bases selon leur force, et comment prévoir le sens d'une réaction acido-basique ?} Nous commençons par le cas le plus simple, celui des acides et des bases dits \textbf{forts}, qui réagissent totalement avec l'eau.

\courssub{Rappel : le pH, thermomètre de l'acidité}

\begin{definition}[pH d'une solution aqueuse]
Le \cle{pH} (potentiel hydrogène) d'une solution aqueuse diluée est défini par :
\[ \mathrm{pH} = -\log\,[\ce{H3O+}] \]
où $[\ce{H3O+}]$ est la concentration molaire en ions oxonium, exprimée en \si{mol.L^{-1}}. Réciproquement : $[\ce{H3O+}] = 10^{-\mathrm{pH}}$.
\end{definition}

Quelques repères à \SI{25}{\celsius} : une solution est \textbf{acide} si $\mathrm{pH} < 7$, \textbf{neutre} si $\mathrm{pH} = 7$, \textbf{basique} si $\mathrm{pH} > 7$. Le pH se mesure expérimentalement avec un \cle{pH-mètre} préalablement étalonné à l'aide de solutions étalons (par exemple pH $= 4,0$ et pH $= 7,0$). L'électrode de verre, plongée dans la solution, délivre une tension proportionnelle au pH.

\begin{attention}[Précision des mesures de pH]
Un pH-mètre de laboratoire donne le pH à $0,05$ unité près au mieux. Il est donc illusoire d'écrire un pH avec trois décimales. De plus, la relation $\mathrm{pH} = -\log[\ce{H3O+}]$ n'est rigoureuse que pour les solutions \textbf{diluées} ($[\ce{H3O+}] \leq \SI{0,05}{mol.L^{-1}}$ environ).
\end{attention}

\courssub{Étude expérimentale de solutions d'acide chlorhydrique}

\begin{experience}[pH de solutions d'acide chlorhydrique de concentrations connues]
\textbf{Protocole.} On prépare par dilutions successives, à partir d'une solution mère d'acide chlorhydrique, quatre solutions filles de concentrations $C$ connues. On mesure le pH de chacune avec un pH-mètre étalonné, en rinçant l'électrode à l'eau distillée entre deux mesures.

\textbf{Résultats} (à \SI{25}{\celsius}) :
\begin{center}
\begin{tabularx}{\linewidth}{|l|X|X|X|X|}
\hline
\rowcolor{popblueL}
$C$ (\si{mol.L^{-1}}) & $1,0\times10^{-1}$ & $1,0\times10^{-2}$ & $2,0\times10^{-3}$ & $1,0\times10^{-4}$ \\
\hline
pH mesuré & $1,0$ & $2,0$ & $2,7$ & $4,0$ \\
\hline
$-\log C$ & $1,0$ & $2,0$ & $2,7$ & $4,0$ \\
\hline
\end{tabularx}
\end{center}

\textbf{Observation.} Pour chaque solution, le pH mesuré est \emph{égal} à $-\log C$ (aux incertitudes de mesure près). On en déduit que $[\ce{H3O+}] = C$ : chaque molécule de \ce{HCl} apportée a produit exactement un ion \ce{H3O+}.
\end{experience}

\begin{popfigure}
\begin{center}
\begin{tikzpicture}
\begin{axis}[
    width=0.85\linewidth, height=6.5cm,
    xlabel={$-\log C$}, ylabel={pH mesuré},
    xmin=0, xmax=4.5, ymin=0, ymax=4.5,
    grid=both, grid style={popline!40},
    legend pos=north west,
    legend style={font=\small}
]
\addplot[popcurve, domain=0:4.5, samples=50]{x};
\addlegendentry{$\mathrm{pH} = -\log C$ (droite de pente 1)}
\addplot[only marks, mark=*, mark size=2.5pt, color=poporange] coordinates {(1,1) (2,2) (2.7,2.7) (4,4)};
\addlegendentry{Points expérimentaux (HCl)}
\end{axis}
\end{tikzpicture}
\end{center}
\legende{Le pH mesuré des solutions d'acide chlorhydrique coïncide avec $-\log C$ : les points expérimentaux sont alignés sur la droite de pente 1. C'est la signature d'une ionisation totale.}
\end{popfigure}

\courssub{Conclusion : l'acide chlorhydrique est un acide fort}

L'égalité $[\ce{H3O+}] = C$ signifie que la réaction entre le chlorure d'hydrogène et l'eau est \textbf{totale} : il ne reste pratiquement plus de molécules \ce{HCl} dans la solution. On écrit donc cette réaction avec une \textbf{flèche simple} :

\[ \ce{HCl + H2O -> H3O+ + Cl-} \]

\begin{popfigure}
\begin{center}
\begin{tikzpicture}[scale=0.9, every node/.style={font=\small}, node distance=1.2cm]
% Avant
\node[font=\bfseries, text=popblue] (titreA) at (0,3) {Avant mise dans l'eau};
\node[draw, circle, fill=popblueL, minimum size=1.2cm] (mol1) at (-2,1) {\ce{HCl}};
\node[draw, circle, fill=popblueL, minimum size=1.2cm] (mol2) at (0,0) {\ce{HCl}};
\node[draw, circle, fill=popblueL, minimum size=1.2cm] (mol3) at (2,1.5) {\ce{HCl}};
\node[draw, circle, fill=popblueL, minimum size=1.2cm] (mol4) at (0.5,2.5) {\ce{HCl}};
% Flèche
\draw[-{Stealth[length=4mm]}, line width=1.5pt, poporange] (3.5,1.5) -- (5.5,1.5) node[midway, above, popdark]{ionisation totale};
% Après
\node[font=\bfseries, text=popgreen] (titreB) at (9,3) {Après (solution aqueuse)};
\node[draw, circle, fill=popgreenL, minimum size=1.2cm] (io1) at (6.5,2.2) {\ce{H3O+}};
\node[draw, circle, fill=popgreenL, minimum size=1.2cm] (io2) at (8.5,0.8) {\ce{H3O+}};
\node[draw, circle, fill=popgreenL, minimum size=1.2cm] (io3) at (10.5,2) {\ce{H3O+}};
\node[draw, circle, fill=popgreenL, minimum size=1.2cm] (io4) at (7.5,2.8) {\ce{H3O+}};
\node[draw, circle, fill=poppinkL, minimum size=1cm] (io5) at (7.3,1) {\ce{Cl-}};
\node[draw, circle, fill=poppinkL, minimum size=1cm] (io6) at (9.3,2.6) {\ce{Cl-}};
\node[draw, circle, fill=poppinkL, minimum size=1cm] (io7) at (10.7,0.8) {\ce{Cl-}};
\node[draw, circle, fill=poppinkL, minimum size=1cm] (io8) at (6.3,0.4) {\ce{Cl-}};
\end{tikzpicture}
\end{center}
\legende{Ionisation totale de \ce{HCl} dans l'eau : toutes les molécules de chlorure d'hydrogène sont transformées en ions \ce{H3O+} et \ce{Cl-}. Il ne subsiste aucune molécule \ce{HCl} en solution.}
\end{popfigure}

\begin{definition}[Acide fort]
Un \cle{acide fort} est un acide dont la réaction avec l'eau est \textbf{totale} : il est entièrement ionisé en solution aqueuse. Pour une solution d'acide fort de concentration $C$ (solution diluée) :
\[ [\ce{H3O+}] = C \qquad \text{donc} \qquad \boxed{\mathrm{pH} = -\log C} \]
Exemples d'acides forts : l'acide chlorhydrique \ce{HCl}, l'acide nitrique \ce{HNO3}, l'acide bromhydrique \ce{HBr}, l'acide sulfurique \ce{H2SO4} (pour sa première acidité : \ce{H2SO4 + H2O -> H3O+ + HSO4-}).
\end{definition}

\begin{methode}[Vérifier expérimentalement qu'un acide est fort]
\begin{enumerate}
\item Mesurer le pH de la solution avec un pH-mètre étalonné.
\item Calculer $-\log C$, où $C$ est la concentration de l'acide apporté.
\item Comparer : si $\mathrm{pH} = -\log C$ (à $0,1$ unité près), alors $[\ce{H3O+}] = C$ et l'ionisation est totale : l'acide est \textbf{fort}. Si $\mathrm{pH} > -\log C$, il reste des molécules non ionisées : l'acide est \textbf{faible}.
\end{enumerate}
\end{methode}

\begin{exemplebox}[Calcul du pH d'une solution d'acide chlorhydrique]
Une solution d'acide chlorhydrique a pour concentration $C = \SI{2,0e-3}{mol.L^{-1}}$. Comme \ce{HCl} est un acide fort :
\[ \mathrm{pH} = -\log C = -\log(2,0\times10^{-3}) = 3 - \log 2 \approx 2,7. \]
Vérification : $[\ce{H3O+}] = 10^{-2,7} \approx \SI{2,0e-3}{mol.L^{-1}} = C$. L'ionisation est bien totale.
\end{exemplebox}

\courssub{Étude des solutions d'hydroxyde de sodium : la base forte}

Menons la même démarche avec la soude (hydroxyde de sodium \ce{NaOH}), solide ionique blanc très utilisé au laboratoire et dans les savonneries artisanales.

\begin{experience}[pH de solutions de soude de concentrations connues]
On prépare des solutions d'hydroxyde de sodium de concentrations $C$ connues et on mesure leur pH à \SI{25}{\celsius} :
\begin{center}
\begin{tabularx}{\linewidth}{|l|X|X|X|}
\hline
\rowcolor{popblueL}
$C$ (\si{mol.L^{-1}}) & $1,0\times10^{-1}$ & $1,0\times10^{-2}$ & $5,0\times10^{-4}$ \\
\hline
pH mesuré & $13,0$ & $12,0$ & $10,7$ \\
\hline
$14 + \log C$ & $13,0$ & $12,0$ & $10,7$ \\
\hline
\end{tabularx}
\end{center}
\textbf{Observation.} Le pH mesuré est égal à $14 + \log C$. Or le produit ionique de l'eau impose $[\ce{H3O+}]\cdot[\ce{HO-}] = K_e = 10^{-14}$ à \SI{25}{\celsius}. On en tire :
\[ [\ce{HO-}] = \frac{K_e}{[\ce{H3O+}]} = \frac{10^{-14}}{10^{-\mathrm{pH}}} = C. \]
Chaque unité \ce{NaOH} apportée a donc libéré exactement un ion hydroxyde.
\end{experience}

La soude est un cristal ionique $\{\ce{Na+} ; \ce{HO-}\}$ : sa dissolution dans l'eau est une \textbf{dissociation totale} :
\[ \ce{NaOH(s) ->[H2O] Na+(aq) + HO-(aq)} \]

\begin{definition}[Base forte]
Une \cle{base forte} est une base dont la réaction avec l'eau est \textbf{totale}. Pour une solution de base forte de concentration $C$ (solution diluée, à \SI{25}{\celsius}) :
\[ [\ce{HO-}] = C \qquad \text{donc} \qquad \boxed{\mathrm{pH} = 14 + \log C} \]
Exemples de bases fortes : l'hydroxyde de sodium \ce{NaOH} (soude) et l'hydroxyde de potassium \ce{KOH} (potasse), très utilisés dans la fabrication artisanale du savon à partir de l'huile de palme.
\end{definition}

\begin{exoresolu}[pH d'une solution de potasse]
On dissout \SI{1,12}{g} d'hydroxyde de potassium \ce{KOH} ($M = \SI{56,0}{g.mol^{-1}}$) dans de l'eau distillée pour obtenir \SI{500}{mL} de solution. Calculer le pH de cette solution à \SI{25}{\celsius}.
\tcblower
\textbf{Étape 1 : quantité de matière dissoute.}
\[ n(\ce{KOH}) = \frac{m}{M} = \frac{1,12}{56,0} = \SI{2,0e-2}{mol}. \]
\textbf{Étape 2 : concentration de la solution.}
\[ C = \frac{n}{V} = \frac{2,0\times10^{-2}}{0,500} = \SI{4,0e-2}{mol.L^{-1}}. \]
\textbf{Étape 3 : pH.} \ce{KOH} est une base forte, totalement dissociée : $[\ce{HO-}] = C$. Donc :
\[ \mathrm{pH} = 14 + \log C = 14 + \log(4,0\times10^{-2}) = 14 - 1,4 = 12,6. \]
La solution est fortement basique, ce qui est cohérent : $\mathrm{pH} > 7$.
\end{exoresolu}

\courssub{Conséquence remarquable : l'effet nivelant de l'eau}

Puisque tout acide fort est intégralement transformé en ions \ce{H3O+}, la seule espèce acide réellement présente dans l'eau est l'ion oxonium \ce{H3O+}, quel que soit l'acide fort de départ. Il en résulte une conséquence capitale :

\begin{propriete}[Effet nivelant de l'eau]
Deux solutions d'acides forts de \textbf{même concentration} ont le \textbf{même pH}, quelle que soit la nature de l'acide (\ce{HCl}, \ce{HNO3}, \ce{HBr}\ldots). On dit que l'eau exerce un \cle{effet nivelant} : l'ion \ce{H3O+} est l'acide le plus fort qui puisse exister en solution aqueuse, et \ce{HO-} la base la plus forte. Toutes les différences de force entre acides forts sont « nivelées » par le solvant.
\end{propriete}

Ainsi, une solution d'acide chlorhydrique et une solution d'acide nitrique, toutes deux à $C = \SI{1,0e-2}{mol.L^{-1}}$, affichent toutes les deux $\mathrm{pH} = 2,0$.

\begin{attention}[Ne pas confondre « fort » et « concentré »]
C'est l'erreur la plus fréquente au baccalauréat !
\begin{itemize}
\item La \textbf{force} d'un acide décrit sa capacité à s'ioniser dans l'eau : un acide fort est \emph{totalement} ionisé, quelle que soit sa concentration.
\item La \textbf{concentration} décrit la quantité d'acide dissoute par litre de solution.
\end{itemize}
On peut donc avoir un \emph{acide fort dilué} (peu de \ce{HCl}, mais tout est ionisé) et un \emph{acide faible concentré} (beaucoup de vinaigre, mais peu de molécules ionisées). « Fort » ne veut pas dire « dangereux » et « concentré » ne veut pas dire « fort » : ce sont deux notions indépendantes.
\end{attention}

\begin{savaistu}[L'acide chlorhydrique dans ton estomac]
L'estomac humain sécrète en permanence de l'acide chlorhydrique : le suc gastrique a un pH voisin de $1,5$, soit $[\ce{H3O+}] \approx \SI{3e-2}{mol.L^{-1}}$. Cet acide fort remplit deux missions : il tue la plupart des microbes avalés avec les aliments (une barrière sanitaire essentielle lorsqu'on consomme de la street food !) et il active la pepsine, l'enzyme qui découpe les protéines. La paroi de l'estomac se protège elle-même grâce à un mucus riche en hydrogénocarbonates \ce{HCO3-}, des ions basiques qui neutralisent localement l'acide. Quand cette protection faiblit, surviennent les brûlures d'estomac, soulagées par des antiacides.
\end{savaistu}

\begin{aretenir}[L'essentiel de la section]
\begin{itemize}
\item $\mathrm{pH} = -\log[\ce{H3O+}]$ ; le pH se mesure au pH-mètre étalonné.
\item Un \textbf{acide fort} est totalement ionisé dans l'eau : $\ce{HCl + H2O -> H3O+ + Cl-}$, d'où $\mathrm{pH} = -\log C$.
\item Une \textbf{base forte} est totalement dissociée : $\ce{NaOH ->[H2O] Na+ + HO-}$, d'où $\mathrm{pH} = 14 + \log C$ à \SI{25}{\celsius}.
\item Critère expérimental : un acide est fort si le pH mesuré est égal à $-\log C$.
\item Effet nivelant : dans l'eau, tous les acides forts se ramènent à \ce{H3O+} ; à concentration égale, ils ont le même pH.
\item Ne jamais confondre force (ionisation totale ou non) et concentration (quantité dissoute).
\end{itemize}
\end{aretenir}

\begin{exercice}[À toi de jouer]
On mesure à \SI{25}{\celsius} le pH de deux solutions de même concentration $C = \SI{5,0e-3}{mol.L^{-1}}$ : la solution A (acide chlorhydrique) a un pH de $2,3$ ; la solution B (acide éthanoïque) a un pH de $3,5$.
\begin{enumerate}
\item Calculer $-\log C$.
\item Montrer que l'acide chlorhydrique est un acide fort.
\item L'acide éthanoïque est-il totalement ionisé ? Justifier et conclure sur sa force.
\end{enumerate}
\end{exercice}

\begin{corrige}[Exercice : À toi de jouer]
\textbf{1.} $-\log C = -\log(5,0\times10^{-3}) = 3 - \log 5 \approx 2,3$.

\textbf{2.} Pour la solution A : $\mathrm{pH} = 2,3 = -\log C$, donc $[\ce{H3O+}] = C$ : l'ionisation de \ce{HCl} est totale, c'est un acide fort.

\textbf{3.} Pour la solution B : $\mathrm{pH} = 3,5 > -\log C = 2,3$, donc $[\ce{H3O+}] = 10^{-3,5} \approx \SI{3,2e-4}{mol.L^{-1}} < C = \SI{5,0e-3}{mol.L^{-1}}$. Seule une petite partie des molécules d'acide éthanoïque s'est ionisée : l'ionisation est \textbf{partielle}, l'acide éthanoïque est un \textbf{acide faible}. C'est exactement ce qui explique que le vinaigre du marché Mokolo soit moins agressif que l'acide chlorhydrique du laboratoire, à concentration égale.
\end{corrige}

\coursec{Acides faibles et bases faibles}

Dans la section précédente, nous avons rencontré des acides et des bases \emph{sans compromis} : l'acide chlorhydrique et l'hydroxyde de sodium se transforment \textbf{totalement} en ions dans l'eau, si bien que leur pH se calcule directement à partir de la concentration ($\mathrm{pH} = -\log C$ pour l'acide fort, $\mathrm{pH} = 14 + \log C$ pour la base forte). Mais la nature est plus nuancée. Le vinaigre que tu verses sur ta salade, le jus de citron, le bissap légèrement acidulé contiennent des acides qui, eux, ne se montrent pas aussi radicaux. Cette section est consacrée à ces espèces \emph{timides} : les \cle{acides faibles} et les \cle{bases faibles}.

\courssub{Étude expérimentale d'une solution d'acide éthanoïque}

\textbf{Une mesure qui dérange.}\par

Reprenons la démarche qui a si bien fonctionné pour l'acide chlorhydrique. On prépare une solution d'\cle{acide éthanoïque} \ce{CH3COOH} (l'acide du vinaigre) de concentration molaire apportée $C = \SI{1,0e-2}{mol.L^{-1}}$, on étalonne un pH-mètre avec des solutions tampons, puis on mesure le pH de la solution à \SI{25}{\celsius}.

\begin{experience}[Mesure du pH d'une solution d'acide éthanoïque]
\textbf{Dispositif et protocole :} on dissout une quantité connue d'acide éthanoïque dans l'eau distillée de façon à obtenir une solution de concentration $C = \SI{1,0e-2}{mol.L^{-1}}$. On plonge l'électrode combinée d'un pH-mètre étalonné dans la solution maintenue à \SI{25}{\celsius} et on lit la valeur stabilisée.

\textbf{Observation :} le pH-mètre affiche $\mathrm{pH} = 3{,}4$.

\textbf{Confrontation :} si l'acide éthanoïque se comportait comme un acide fort, on aurait $\mathrm{pH} = -\log C = -\log(1{,}0\times 10^{-2}) = 2{,}0$. Or la valeur mesurée ($3{,}4$) est \emph{nettement supérieure} à $2{,}0$.
\end{experience}

\textbf{Interprétation.}\par

Que signifie ce pH trop élevé ? Traduisons-le en concentration :
\[
[\ce{H3O+}] = 10^{-\mathrm{pH}} = 10^{-3{,}4} \approx \SI{4,0e-4}{mol.L^{-1}}.
\]
Or $C = \SI{1,0e-2}{mol.L^{-1}}$. On constate donc que
\[
[\ce{H3O+}] \approx \SI{4,0e-4}{mol.L^{-1}} \ll C = \SI{1,0e-2}{mol.L^{-1}}.
\]
La solution contient beaucoup moins d'ions oxonium que ce qu'aurait produit une ionisation totale. Conclusion inévitable : \textbf{toutes les molécules d'acide éthanoïque n'ont pas cédé leur proton à l'eau}. La grande majorité des molécules \ce{CH3COOH} est restée intacte ; seule une petite fraction s'est ionisée. L'acide éthanoïque est un \cle{acide faible}.

\begin{definition}[Acide faible]
Un \cle{acide faible} est un acide qui ne s'ionise que \textbf{partiellement} dans l'eau. Sa réaction avec l'eau est \textbf{limitée} et conduit à un \textbf{équilibre chimique} : en solution, la forme acide non ionisée coexiste avec la forme basique et les ions oxonium.
\end{definition}

\textbf{Équation de la réaction avec l'eau.}\par

Puisque la transformation n'est pas totale, on l'écrit avec une \textbf{double flèche}, symbole de l'équilibre :
\[
\ce{CH3COOH + H2O <=> CH3COO- + H3O+}
\]
Lisons cette équation : une molécule d'acide éthanoïque cède un proton \ce{H+} à une molécule d'eau ; il se forme l'ion \cle{éthanoate} \ce{CH3COO-} (la base conjuguée) et l'ion oxonium \ce{H3O+}. Mais la réaction inverse existe aussi : l'ion éthanoate peut capter un proton de \ce{H3O+} pour régénérer l'acide. À l'équilibre, les deux réactions se compensent et les concentrations ne varient plus. Cette réaction se déroule à température ambiante, sans catalyseur, et elle est \textbf{limitée} (non totale) : c'est sa caractéristique essentielle.

\begin{attention}[La double flèche est obligatoire]
Erreur fréquente au baccalauréat : écrire $\ce{CH3COOH + H2O -> CH3COO- + H3O+}$ avec une flèche simple. Une flèche simple signifie une réaction \emph{totale}, ce qui est faux pour un acide faible. Pour toute réaction d'un acide faible (ou d'une base faible) avec l'eau, il faut impérativement la double flèche \ce{<=>}. Le correcteur y sera attentif !
\end{attention}

\begin{popfigure}
\begin{center}
\begin{tikzpicture}[scale=0.9]
% Bécher
\draw[thick,popdark] (-4.2,-0.2) -- (-4.2,-3.2) arc[start angle=180,end angle=360,radius=2.1] -- (4.2,-0.2);
\draw[thick,popdark] (-4.2,-0.2) -- (-4.5,-0.2);
\draw[thick,popdark] (4.2,-0.2) -- (4.5,-0.2);
% Eau
\fill[popblue!15] (-4.15,-0.6) rectangle (4.15,-3.05);
\draw[popblue,thick] (-4.15,-0.6) -- (4.15,-0.6);
\node[popblue,anchor=west] at (-4.0,-0.95) {\small eau};
% Molécules intactes (majoritaires)
\foreach \x/\y in {-3.2/-1.6, -1.8/-2.4, -0.6/-1.5, 0.6/-2.5, 1.7/-1.7, 2.9/-2.5, -2.7/-2.8, 0.1/-2.0, 3.3/-1.4, -1.1/-2.9}{
  \node[circle,draw=popgreen,fill=popgreenL,inner sep=1.2pt,minimum size=9pt] at (\x,\y) {\tiny \ce{CH3COOH}};
}
% Quelques ions
\node[circle,draw=poporange,fill=poporangeL,inner sep=1pt,minimum size=8pt] at (-2.3,-1.2) {\tiny \ce{H3O+}};
\node[circle,draw=poppink,fill=poppinkL,inner sep=1pt,minimum size=8pt] at (1.2,-1.1) {\tiny \ce{CH3COO-}};
\node[circle,draw=poporange,fill=poporangeL,inner sep=1pt,minimum size=8pt] at (2.4,-2.9) {\tiny \ce{H3O+}};
\node[circle,draw=poppink,fill=poppinkL,inner sep=1pt,minimum size=8pt] at (-0.9,-2.2) {\tiny \ce{CH3COO-}};
\end{tikzpicture}
\end{center}
\legende{Ionisation partielle de l'acide éthanoïque dans l'eau : les molécules \ce{CH3COOH} intactes (vert) sont largement majoritaires ; seules quelques-unes se sont transformées en ions \ce{H3O+} (orange) et \ce{CH3COO-} (rose).}
\end{popfigure}

\courssub{Le coefficient d'ionisation $\alpha$}

Pour quantifier le caractère \og faible \fg{} d'un acide, on introduit une grandeur simple et parlante : la proportion de molécules qui se sont effectivement ionisées.

\begin{definition}[Coefficient d'ionisation]
Le \cle{coefficient d'ionisation} (ou taux d'ionisation) d'un acide est le quotient de la quantité d'acide ionisé par la quantité d'acide introduite :
\[
\alpha = \frac{n_{\text{ionisé}}}{n_{\text{introduit}}} = \frac{[\ce{H3O+}]_{\text{eq}}}{C}.
\]
On l'exprime souvent en pourcentage. Pour un acide \textbf{fort}, $\alpha = 1$ (ionisation totale, $100\,\%$). Pour un acide \textbf{faible}, $\alpha < 1$ : plus $\alpha$ est petit, plus l'acide est faible.
\end{definition}

\begin{exemplebox}[Calcul du coefficient d'ionisation de l'acide éthanoïque]
Reprenons notre solution à $C = \SI{1,0e-2}{mol.L^{-1}}$ et $\mathrm{pH} = 3{,}4$.
\begin{enumerate}
\item Concentration en ions oxonium : $[\ce{H3O+}] = 10^{-3{,}4} \approx \SI{4,0e-4}{mol.L^{-1}}$.
\item Coefficient d'ionisation :
\[
\alpha = \frac{[\ce{H3O+}]}{C} = \frac{4{,}0\times 10^{-4}}{1{,}0\times 10^{-2}} = 4{,}0\times 10^{-2} \approx 4\,\%.
\]
\end{enumerate}
Seulement 4 molécules d'acide éthanoïque sur 100 se sont ionisées ! L'acide éthanoïque est bien un acide faible, et même nettement faible.
\end{exemplebox}

\courssub{Comparaison graphique avec l'acide chlorhydrique}

La différence de comportement entre un acide fort et un acide faible apparaît de façon éclatante lorsqu'on trace le pH en fonction de $-\log C$ pour différentes concentrations.

\begin{popfigure}
\begin{center}
\begin{tikzpicture}
\begin{axis}[
  width=0.85\linewidth, height=7cm,
  xlabel={$-\log C$}, ylabel={pH},
  xmin=0, xmax=5, ymin=0, ymax=7,
  grid=both, grid style={popline!40},
  legend style={at={(0.03,0.03)},anchor=south west,font=\small,draw=popline},
  axis line style={popdark},
]
% HCl : pH = -log C
\addplot[popcurve,popblue,domain=0:5,samples=50]{x};
\addlegendentry{\ce{HCl} (acide fort) : $\mathrm{pH} = -\log C$}
% CH3COOH : pH = 0.5*(pKa - log C), pKa = 4.8 ; x = -log C => pH = 0.5*(4.8 + x)
\addplot[popcurve,poporange,domain=0:5,samples=50]{0.5*(4.8+x)};
\addlegendentry{\ce{CH3COOH} (acide faible, $\mathrm{p}K_{\mathrm{a}}=4{,}8$)}
% Point mesuré
\addplot[only marks,mark=*,mark size=2pt,popdark] coordinates {(2,3.4)};
\node[popdark,font=\small,anchor=south west] at (axis cs:2.05,3.4) {$C = 10^{-2}$, pH $= 3{,}4$};
\end{axis}
\end{tikzpicture}
\end{center}
\legende{pH en fonction de $-\log C$ à \SI{25}{\celsius}. Pour l'acide fort \ce{HCl}, les points se placent sur la droite $\mathrm{pH} = -\log C$. Pour l'acide faible \ce{CH3COOH}, la courbe se situe \emph{au-dessus} : à concentration égale, l'acide faible donne un pH plus élevé (moins d'ions \ce{H3O+}).}
\end{popfigure}

La lecture de ce graphique est un outil d'analyse puissant : si les points expérimentaux $(-\log C\,;\,\mathrm{pH})$ d'un acide se placent sur la première bissectrice, l'acide est fort ; s'ils se placent au-dessus, il est faible.

\begin{methode}[Montrer expérimentalement qu'un acide est faible]
\begin{enumerate}
\item Mesurer le pH de la solution de concentration $C$ connue avec un pH-mètre étalonné.
\item Calculer $-\log C$.
\item Comparer :
  \begin{itemize}
  \item si $\mathrm{pH} = -\log C$ (c'est-à-dire $[\ce{H3O+}] = C$), l'acide est \textbf{fort} ;
  \item si $\mathrm{pH} > -\log C$ (c'est-à-dire $[\ce{H3O+}] < C$), l'acide est \textbf{faible}.
  \end{itemize}
\item Conclure et, si l'acide est faible, écrire l'équation de sa réaction avec l'eau avec une double flèche, puis calculer éventuellement $\alpha = \dfrac{[\ce{H3O+}]}{C}$.
\end{enumerate}
\end{methode}

\courssub{Étude d'une solution de base faible : l'éthanoate de sodium}

\textbf{Une mesure symétrique.}\par

L'ion éthanoate \ce{CH3COO-}, formé lors de l'ionisation de l'acide éthanoïque, est capable de \emph{capter} un proton : c'est une base au sens de Brønsted. Étudions une solution d'\cle{éthanoate de sodium} \ce{CH3COONa} de concentration $C = \SI{1,0e-2}{mol.L^{-1}}$. Le sel se dissout totalement en ions \ce{Na+} (spectateur) et \ce{CH3COO-}. Le pH mesuré à \SI{25}{\celsius} vaut environ $8{,}4$.

Si l'ion éthanoate était une base \emph{forte}, il capterait totalement les protons de l'eau et l'on aurait $[\ce{HO-}] = C$, soit $\mathrm{pH} = 14 + \log C = 14 + \log(10^{-2}) = 12{,}0$. Or le pH mesuré ($8{,}4$) est \textbf{inférieur} à $12{,}0$. Calculons :
\[
[\ce{HO-}] = \frac{K_e}{[\ce{H3O+}]} = \frac{10^{-14}}{10^{-8{,}4}} = 10^{-5{,}6} \approx \SI{2,5e-6}{mol.L^{-1}} \ll C.
\]
La quantité d'ions hydroxyde formés est très inférieure à la quantité d'ions éthanoate introduits : la réaction de l'ion éthanoate avec l'eau est \textbf{partielle}. L'ion éthanoate est une \cle{base faible}. Il en va de même de l'ammoniac \ce{NH3}, base faible bien connue des produits ménagers.

\begin{definition}[Base faible]
Une \cle{base faible} est une base dont la réaction avec l'eau (captage d'un proton, ou \emph{protonation}) est \textbf{partielle} et conduit à un \textbf{équilibre chimique}. En solution, la forme basique coexiste avec sa forme acide conjuguée et les ions hydroxyde.
\end{definition}

\textbf{Équations des réactions avec l'eau.}\par

Pour l'ion éthanoate :
\[
\ce{CH3COO- + H2O <=> CH3COOH + HO-}
\]
L'ion éthanoate arrache un proton à une molécule d'eau ; il se forme de l'acide éthanoïque et des ions hydroxyde, responsables du caractère basique de la solution (pH $> 7$). Réaction limitée, à température ambiante, sans catalyseur.

Pour l'ammoniac, autre base faible de référence :
\[
\ce{NH3 + H2O <=> NH4+ + HO-}
\]
Remarque la symétrie parfaite avec les acides faibles : même double flèche, même caractère limité, même coexistence des formes à l'équilibre.

\begin{methode}[Montrer expérimentalement qu'une base est faible]
\begin{enumerate}
\item Mesurer le pH de la solution de concentration $C$ connue.
\item Calculer $14 + \log C$ (valeur attendue pour une base forte).
\item Comparer :
  \begin{itemize}
  \item si $\mathrm{pH} = 14 + \log C$ (c'est-à-dire $[\ce{HO-}] = C$), la base est \textbf{forte} ;
  \item si $\mathrm{pH} < 14 + \log C$ (c'est-à-dire $[\ce{HO-}] < C$), la base est \textbf{faible}.
  \end{itemize}
\item Si la base est faible, écrire l'équation avec une double flèche et calculer le taux de protonation $\alpha = \dfrac{[\ce{HO-}]}{C}$.
\end{enumerate}
\end{methode}

\courssub{Influence de la dilution sur l'ionisation}

Voici un résultat qui surprend souvent : quand on \textbf{dilue} une solution d'acide faible, son coefficient d'ionisation $\alpha$ \textbf{augmente}. Autrement dit, plus l'acide est dilué, plus la \emph{proportion} de molécules ionisées est grande — même si, en valeur absolue, la concentration en \ce{H3O+} diminue et le pH augmente.

\begin{exemplebox}[Dilution et ionisation de l'acide éthanoïque]
Pour l'acide éthanoïque à \SI{25}{\celsius}, les mesures donnent :
\begin{center}
\begin{tabularx}{\linewidth}{|l|X|X|X|}
\hline
\rowcolor{popblueL}
$C$ (mol.L$^{-1}$) & pH mesuré & $[\ce{H3O+}]$ (mol.L$^{-1}$) & $\alpha = [\ce{H3O+}]/C$ \\
\hline
$1{,}0\times 10^{-1}$ & $2{,}9$ & $1{,}3\times 10^{-3}$ & $1{,}3\,\%$ \\
\hline
$1{,}0\times 10^{-2}$ & $3{,}4$ & $4{,}0\times 10^{-4}$ & $4{,}0\,\%$ \\
\hline
$1{,}0\times 10^{-3}$ & $3{,}9$ & $1{,}3\times 10^{-4}$ & $13\,\%$ \\
\hline
\end{tabularx}
\end{center}
Quand la concentration est divisée par 100, le coefficient d'ionisation est multiplié par 10 : l'ionisation, bien que toujours partielle, devient proportionnellement plus importante à mesure que l'on dilue.
\end{exemplebox}

\begin{propriete}[Loi de dilution d'Ostwald (qualitative)]
Pour un acide faible (ou une base faible), le coefficient d'ionisation $\alpha$ \textbf{augmente lorsque la concentration diminue}. À dilution infinie, $\alpha$ tend vers 1 : un acide faible très dilué se comporte presque comme un acide fort. Attention cependant : le pH, lui, augmente toujours lors d'une dilution, car la concentration totale en \ce{H3O+} diminue.
\end{propriete}

\begin{attention}[{Ne pas confondre $\alpha$ et $[\ce{H3O+}]$}]
Deux pièges classiques :
\begin{itemize}
\item \og Diluer augmente $\alpha$, donc le pH diminue \fg{} : \textbf{faux} ! Même si la \emph{proportion} ionisée augmente, la quantité totale d'acide par litre diminue davantage, donc $[\ce{H3O+}]$ diminue et le pH \emph{augmente}.
\item \og Un acide faible, c'est un acide dilué \fg{} : \textbf{faux} ! \emph{Faible} qualifie la force (caractère intrinsèque de l'espèce, lié à l'équilibre de sa réaction avec l'eau) ; \emph{dilué} qualifie la concentration. On peut avoir une solution concentrée d'acide faible (vinaigre pur) ou diluée d'acide fort.
\end{itemize}
\end{attention}

\begin{exoresolu}[Identification d'une base faible]
\textbf{Énoncé.} On dissout du méthanoate de sodium \ce{HCOONa} dans l'eau distillée ; on obtient une solution de concentration $C = \SI{5,0e-2}{mol.L^{-1}}$. Le pH mesuré à \SI{25}{\celsius} vaut $8{,}4$.
\begin{enumerate}
\item Montrer que l'ion méthanoate \ce{HCOO-} est une base faible.
\item Écrire l'équation de sa réaction avec l'eau.
\item Calculer son coefficient de protonation $\alpha$.
\end{enumerate}
\tcblower
\textbf{Solution détaillée.}
\begin{enumerate}
\item Si \ce{HCOO-} était une base forte, on aurait $[\ce{HO-}] = C$, soit
\[
\mathrm{pH} = 14 + \log C = 14 + \log(5{,}0\times 10^{-2}) = 14 - 1{,}3 = 12{,}7.
\]
Or le pH mesuré ($8{,}4$) est \textbf{inférieur} à $12{,}7$, donc $[\ce{HO-}] < C$ : la protonation est partielle, l'ion méthanoate est une \textbf{base faible}.
\item Équation de la réaction avec l'eau (double flèche, équilibre limité) :
\[
\ce{HCOO- + H2O <=> HCOOH + HO-}
\]
Il se forme de l'acide méthanoïque \ce{HCOOH} et des ions hydroxyde.
\item Calcul de $\alpha$ :
\[
[\ce{HO-}] = \frac{K_e}{[\ce{H3O+}]} = \frac{10^{-14}}{10^{-8{,}4}} = 10^{-5{,}6} \approx \SI{2,5e-6}{mol.L^{-1}},
\]
\[
\alpha = \frac{[\ce{HO-}]}{C} = \frac{2{,}5\times 10^{-6}}{5{,}0\times 10^{-2}} = 5{,}0\times 10^{-5} \approx 0{,}005\,\%.
\]
La protonation est extrêmement limitée : l'ion méthanoate est une base faible, et même très faible.
\end{enumerate}
\end{exoresolu}

\begin{savaistu}[Le vinaigre : un acide faible que l'on peut boire]
Le vinaigre est une solution aqueuse d'acide éthanoïque à environ $6$ à $8\,\%$ en masse, obtenu par oxydation de l'éthanol des boissons fermentées (vin, bière, vin de palme laissé à l'air libre !) sous l'action de bactéries du genre \emph{Acetobacter}. C'est précisément parce que l'acide éthanoïque est un acide \emph{faible} — seulement quelques pourcents de ses molécules libèrent des ions \ce{H3O+} — que le vinaigre est consommable et agréable en assaisonnement. Imagine une solution d'acide chlorhydrique (acide fort) à la même concentration : elle attaquerait gravement la bouche et l'œsophage ! La \og faiblesse \fg{} d'un acide n'est donc pas un défaut : c'est une propriété qui conditionne ses usages, de la cuisine à la pharmacopée traditionnelle.
\end{savaistu}

\begin{aretenir}[L'essentiel de la section]
\begin{itemize}
\item Un \textbf{acide faible} ne s'ionise que partiellement dans l'eau : $\ce{AH + H2O <=> A- + H3O+}$ (double flèche, équilibre).
\item Critère expérimental : $\mathrm{pH} > -\log C$, c'est-à-dire $[\ce{H3O+}] < C$.
\item Une \textbf{base faible} ne se protone que partiellement : $\ce{A- + H2O <=> AH + HO-}$. Critère : $\mathrm{pH} < 14 + \log C$, c'est-à-dire $[\ce{HO-}] < C$.
\item Le \textbf{coefficient d'ionisation} $\alpha = \dfrac{\text{quantité ionisée}}{\text{quantité introduite}}$ est inférieur à 1 pour un acide ou une base faible ; il vaut 1 pour un acide ou une base forte.
\item La \textbf{dilution augmente} le coefficient d'ionisation d'un acide faible (loi de dilution d'Ostwald), mais le pH augmente quand même.
\item Exemples de référence : acide éthanoïque \ce{CH3COOH} (acide faible, $\alpha \approx 4\,\%$ à $10^{-2}$ mol.L$^{-1}$), ion éthanoate \ce{CH3COO-} et ammoniac \ce{NH3} (bases faibles).
\end{itemize}
\end{aretenir}

Maintenant que nous savons distinguer qualitativement les acides et bases forts des faibles, il nous faut un outil \emph{quantitatif} pour mesurer et comparer ces forces : ce sera l'objet des sections suivantes, avec l'avancement d'une réaction, puis la constante d'acidité $K_a$.

\coursec{Étude quantitative d'un équilibre : avancement et taux d'avancement}

Dans la section précédente, nous avons découvert que l'acide éthanoïque, contrairement à l'acide chlorhydrique, ne s'ionise pas totalement dans l'eau : sa solution de concentration \SI{1,0e-2}{mol.L^{-1}} affiche un pH d'environ $3,4$ au lieu de $2$. Nous avons dit qualitativement : « l'acide est faible, la réaction est limitée ». Fort bien. Mais en sciences, une affirmation qualitative ne suffit pas : il faut la \cle{quantifier}. De combien la réaction a-t-elle avancé ? Quelle fraction des molécules d'acide s'est réellement transformée en ions ? C'est toute l'ambition de cette section : construire un outil de mesure universel, le \cle{taux d'avancement final} $\tau$, qui nous permettra de classer rigoureusement les acides et les bases.

\courssub{Rappel : le tableau d'avancement}

Reprenons les fondamentaux vus en Première. Pour étudier une transformation chimique, on décrit l'évolution du système à l'aide d'une grandeur : l'\cle{avancement}, noté $x$, exprimé en moles.

\begin{definition}[Avancement d'une réaction]
L'\cle{avancement} $x$ d'une réaction chimique est une grandeur extensive (unité : \si{mol}) qui permet de quantifier l'état d'avancement de la réaction. À chaque instant, la quantité de matière de chaque espèce s'obtient en ajoutant (produits) ou en retranchant (réactifs) $x$ multiplié par son nombre stœchiométrique (signé) à la quantité initiale. L'avancement varie de $x = 0$ (état initial) à une valeur finale $x_f$ (état final).
\end{definition}

Le \cle{tableau d'avancement} est l'outil comptable de la réaction. Il comporte trois lignes essentielles :

\begin{itemize}
\item \textbf{État initial} ($x = 0$) : on y inscrit les quantités de matière introduites ;
\item \textbf{État intermédiaire} ($x$ quelconque) : chaque quantité s'exprime en fonction de $x$ ;
\item \textbf{État final} ($x = x_f$) : on remplace $x$ par sa valeur finale.
\end{itemize}

Il faut bien distinguer deux avancements finaux possibles :

\begin{definition}[Avancement maximal et avancement final]
L'\cle{avancement maximal} $x_{max}$ est l'avancement qu'atteindrait la réaction si elle était \textbf{totale} : il est fixé par le réactif limitant. L'\cle{avancement final} $x_f$ est l'avancement \textbf{réellement} atteint lorsque le système n'évolue plus macroscopiquement (équilibre chimique ou réaction totale). Pour une réaction totale, $x_f = x_{max}$ ; pour une réaction limitée (équilibre), $x_f < x_{max}$.
\end{definition}

\courssub{Application à l'ionisation de l'acide éthanoïque dans l'eau}

Considérons un volume $V$ de solution d'acide éthanoïque de concentration apportée $C$. La quantité d'acide introduite est donc $n_0 = C \cdot V$. L'équation de la réaction avec l'eau s'écrit :

\[ \ce{CH3COOH + H2O <=> CH3COO- + H3O+} \]

Commentons cette équation : elle se déroule à température ambiante, sans catalyseur, et la double flèche $\ce{<=>}$ traduit son caractère \textbf{limité} : c'est un équilibre chimique. L'eau, en très large excès (c'est le solvant), joue le rôle de base : elle capte le proton cédé par l'acide. Le couple mis en jeu côté acide est \ce{CH3COOH}/\ce{CH3COO-}, et côté eau \ce{H3O+}/\ce{H2O}.

Dressons le tableau d'avancement complet :

\begin{popfigure}
\begin{center}
\renewcommand{\arraystretch}{1,6}
\begin{tabularx}{\linewidth}{|l|>{\centering\arraybackslash}X|>{\centering\arraybackslash}X|>{\centering\arraybackslash}X|>{\centering\arraybackslash}X|>{\centering\arraybackslash}X|}
\hline
\rowcolor{popblueL} \textbf{État} & \textbf{Avancement} & \ce{CH3COOH} & \ce{H2O} & \ce{CH3COO-} & \ce{H3O+} \\
\hline
Initial & $x = 0$ & $C \cdot V$ & excès & $0$ & $0$ \\
\hline
Intermédiaire & $x$ & $C \cdot V - x$ & excès & $x$ & $x$ \\
\hline
Final (équilibre) & $x = x_f$ & $C \cdot V - x_f$ & excès & $x_f$ & $x_f$ \\
\hline
Maximal (si totale) & $x = x_{max}$ & $C \cdot V - x_{max} = 0$ & excès & $x_{max}$ & $x_{max}$ \\
\hline
\end{tabularx}
\end{center}
\legende{Tableau d'avancement de la réaction de l'acide éthanoïque avec l'eau, pour un volume $V$ de solution de concentration $C$. On néglige les ions \ce{H3O+} provenant de l'autoprotolyse de l'eau.}
\end{popfigure}

Deux lectures de ce tableau sont capitales :

\begin{itemize}
\item \textbf{Si la réaction était totale} (hypothèse de l'acide fort), tout l'acide disparaîtrait : $C \cdot V - x_{max} = 0$, d'où :
\[ x_{max} = C \cdot V \]
\item \textbf{Dans la réalité}, la réaction s'arrête à l'équilibre avec $x_f < x_{max}$, et la concentration finale en ions oxonium vaut :
\[ [\ce{H3O+}]_f = \frac{x_f}{V} \]
\end{itemize}

\courssub{Mesurer l'avancement final grâce au pH}

Voici l'idée géniale de cette section : les ions \ce{H3O+} formés sont exactement ceux que mesure le pH-mètre ! Or par définition du pH :

\[ [\ce{H3O+}]_f = 10^{-\text{pH}} \]

En combinant avec la relation précédente, on obtient l'avancement final :

\begin{propriete}[Avancement final déduit du pH]
Pour la réaction d'un acide avec l'eau, dans un volume $V$ de solution, l'avancement final est :
\[ x_f = [\ce{H3O+}]_f \times V = 10^{-\text{pH}} \times V \]
avec $x_f$ en \si{mol}, $V$ en \si{L} et le pH sans unité.
\end{propriete}

Autrement dit : une simple mesure de pH nous permet de \emph{compter} les moles d'ions oxonium réellement formées à l'équilibre, donc de savoir « jusqu'où » la réaction est allée.

\begin{methode}[Déterminer $x_f$ puis $\tau$ à partir du pH mesuré]
\begin{enumerate}
\item Relever le pH de la solution au pH-mètre et noter le volume $V$ et la concentration $C$.
\item Calculer la concentration en ions oxonium : $[\ce{H3O+}]_f = 10^{-\text{pH}}$ (en \si{mol.L^{-1}}).
\item En déduire l'avancement final : $x_f = 10^{-\text{pH}} \times V$ (en \si{mol}).
\item Calculer l'avancement maximal : $x_{max} = C \times V$ (en \si{mol}).
\item Conclure par le taux d'avancement final : $\tau = \dfrac{x_f}{x_{max}} = \dfrac{10^{-\text{pH}}}{C}$.
\item Interpréter : si $\tau = 1$, l'acide est fort ; si $\tau < 1$, l'acide est faible.
\end{enumerate}
\end{methode}

\begin{attention}[Le piège du volume]
L'erreur la plus fréquente au Baccalauréat : écrire $x_f = 10^{-\text{pH}}$ en oubliant de multiplier par le volume $V$ ! Attention : $10^{-\text{pH}}$ est une \textbf{concentration} (en \si{mol.L^{-1}}), pas une quantité de matière. Un avancement s'exprime en \textbf{moles}. Il faut donc impérativement écrire $x_f = 10^{-\text{pH}} \times V$. Bonne nouvelle : dans le rapport $\tau = \dfrac{x_f}{x_{max}} = \dfrac{10^{-\text{pH}} \cdot V}{C \cdot V}$, le volume $V$ se simplifie, d'où la formule simplifiée $\tau = \dfrac{10^{-\text{pH}}}{C}$. Mais si l'énoncé demande $x_f$ lui-même, le volume est indispensable !
\end{attention}

\courssub{Le taux d'avancement final : définition et critère de force}

\begin{definition}[Taux d'avancement final]
Le \cle{taux d'avancement final} d'une réaction, noté $\tau$ (lettre grecque « tau »), est le rapport de l'avancement final à l'avancement maximal :
\[ \tau = \frac{x_f}{x_{max}} \]
C'est un nombre sans unité, compris entre $0$ et $1$, souvent exprimé en pourcentage.
\end{definition}

Pour l'ionisation d'un acide de concentration $C$ dans l'eau, la simplification du volume donne immédiatement :

\[ \tau = \frac{[\ce{H3O+}]_f}{C} = \frac{10^{-\text{pH}}}{C} \]

Ce taux est le \textbf{critère quantitatif} qui remplace définitivement le vocabulaire vague de « réaction limitée » :

\begin{aretenir}[Critère de force d'un acide ou d'une base]
\begin{itemize}
\item $\tau = 1$ (soit $100\,\%$) : la réaction avec l'eau est \textbf{totale}. L'acide est un \cle{acide fort} (ou la base une \cle{base forte}) : toutes les molécules se sont ionisées.
\item $\tau < 1$ (soit moins de $100\,\%$) : la réaction est \textbf{limitée}, elle aboutit à un \textbf{équilibre chimique}. L'acide est un \cle{acide faible} (ou la base une \cle{base faible}) : seule une fraction des molécules s'est ionisée.
\end{itemize}
Plus $\tau$ est proche de $1$, plus l'acide est fort ; plus $\tau$ est proche de $0$, plus il est faible.
\end{aretenir}

\begin{popfigure}
\begin{center}
\begin{tikzpicture}[x=6.2cm,y=1.2cm]
% Axe des abscisses (tau)
\draw[popline,thick,->] (0,0) -- (1.1,0) node[below right,popdark]{\small $\tau$};
% Graduations et étiquettes
\foreach \x/\lab in {0/{0}, 0.25/{25\,\%}, 0.5/{50\,\%}, 0.75/{75\,\%}, 1/{100\,\%}} {
  \draw[popline] (\x,0.03) -- (\x,-0.03) node[below,popmuted]{\scriptsize \lab};
}
% Barre acide fort (centrée sur y=0.5)
\fill[popgreen!85] (0,0.35) rectangle (1,0.65);
\node[left,popdark,align=right] at (-0.1,0.5) {\small \ce{HCl} (acide fort)};
\node[right,popdark] at (1.02,0.5) {\scriptsize $\tau = 1$};
% Barre acide faible (centrée sur y=1.3)
\fill[poporange!85] (0,1.15) rectangle (0.04,1.45);
\draw[poporange,thick,->] (0.04,1.3) -- (0.16,1.3) node[right,popdark]{\scriptsize $\tau \approx 0{,}04$};
\node[left,popdark,align=right] at (-0.1,1.3) {\small \ce{CH3COOH} (acide faible)};
\end{tikzpicture}
\end{center}
\legende{Comparaison des taux d'avancement final à concentration égale $C = \SI{1,0e-2}{mol.L^{-1}}$ : l'acide chlorhydrique s'ionise totalement ($\tau = 100\,\%$), l'acide éthanoïque à peine ($\tau \approx 4\,\%$).}
\end{popfigure}

\courssub{Lien entre le taux d'avancement et le coefficient d'ionisation}

En Terminale TI, le programme utilise aussi le \cle{coefficient d'ionisation} $\alpha$ (lettre grecque « alpha »), défini comme la fraction des molécules d'acide (ou de base) qui se sont effectivement ionisées dans l'eau :

\[ \alpha = \frac{\text{quantité d'acide ionisé}}{\text{quantité d'acide introduite}} \]

Or, d'après le tableau d'avancement, la quantité d'acide ionisé à l'équilibre vaut exactement $x_f$, et la quantité introduite vaut $C \cdot V = x_{max}$. Par conséquent :

\begin{propriete}[Égalité entre $\tau$ et $\alpha$ pour un acide faible]
Pour la réaction d'un acide faible (ou d'une base faible) avec l'eau, le taux d'avancement final est égal au coefficient d'ionisation :
\[ \tau = \alpha = \frac{[\ce{H3O+}]_f}{C} \]
Les deux grandeurs expriment la même réalité : le pourcentage de molécules transformées en ions à l'équilibre.
\end{propriete}

Ainsi, dire « le coefficient d'ionisation de l'acide éthanoïque vaut $4\,\%$ à \SI{1,0e-2}{mol.L^{-1}} » ou « le taux d'avancement final de sa réaction avec l'eau vaut $0{,}04$ », c'est dire exactement la même chose. Retiens aussi que $\alpha$ (donc $\tau$) \textbf{dépend de la concentration} : plus on dilue un acide faible, plus il s'ionise (loi de dilution d'Ostwald) — mais il ne devient jamais fort pour autant !

\courssub{Exercice guidé : calcul du taux d'avancement de l'acide éthanoïque}

\begin{exoresolu}[Taux d'avancement d'une solution d'acide éthanoïque]
On prépare un volume $V = \SI{200}{mL}$ de solution d'acide éthanoïque de concentration $C = \SI{1,0e-2}{mol.L^{-1}}$. Le pH-mètre indique $\text{pH} = 3{,}4$.
\begin{enumerate}
\item Calculer l'avancement maximal $x_{max}$ de la réaction de l'acide avec l'eau.
\item Calculer l'avancement final $x_f$.
\item En déduire le taux d'avancement final $\tau$ et le coefficient d'ionisation $\alpha$. Conclure sur la force de l'acide.
\end{enumerate}
\tcblower
\textbf{1. Avancement maximal.} Si la réaction était totale, tout l'acide serait consommé :
\[ x_{max} = C \times V = \SI{1,0e-2}{mol.L^{-1}} \times \SI{0,200}{L} = \SI{2,0e-3}{mol} \]

\textbf{2. Avancement final.} La mesure du pH donne la concentration en ions oxonium :
\[ [\ce{H3O+}]_f = 10^{-\text{pH}} = 10^{-3{,}4} \approx \SI{4,0e-4}{mol.L^{-1}} \]
On multiplie par le volume pour obtenir une quantité de matière :
\[ x_f = [\ce{H3O+}]_f \times V = \SI{4,0e-4}{mol.L^{-1}} \times \SI{0,200}{L} = \SI{8,0e-5}{mol} \]

\textbf{3. Taux d'avancement final.}
\[ \tau = \frac{x_f}{x_{max}} = \frac{\SI{8,0e-5}{mol}}{\SI{2,0e-3}{mol}} = 0{,}04 \]
Vérification par la formule simplifiée (le volume se simplifie) :
\[ \tau = \frac{10^{-\text{pH}}}{C} = \frac{10^{-3{,}4}}{10^{-2}} = 10^{-1{,}4} \approx 0{,}04 \]
Pour un acide faible, $\alpha = \tau = 0{,}04$, soit $\mathbf{4\,\%}$.

\textbf{Conclusion.} $\tau = 0{,}04 < 1$ : la réaction de l'acide éthanoïque avec l'eau est \textbf{limitée}. Seules 4 molécules d'acide sur 100 se sont ionisées ; l'acide éthanoïque est bien un \textbf{acide faible}. À l'équilibre, la solution contient majoritairement des molécules \ce{CH3COOH} intactes, et minoritairement des ions \ce{CH3COO-} et \ce{H3O+}.
\end{exoresolu}

\begin{exemplebox}[Et si c'était un acide fort ?]
Reprenons exactement les mêmes conditions ($C = \SI{1,0e-2}{mol.L^{-1}}$) mais avec de l'acide chlorhydrique. Le pH mesuré vaut $2{,}0$. Alors :
\[ \tau = \frac{10^{-\text{pH}}}{C} = \frac{10^{-2}}{10^{-2}} = 1 \]
$\tau = 100\,\%$ : la réaction $\ce{HCl + H2O -> H3O+ + Cl-}$ est \textbf{totale} (on la note d'ailleurs avec une flèche simple), et \ce{HCl} est un acide fort. Le même calcul, deux comportements radicalement différents : voilà toute la puissance du taux d'avancement.
\end{exemplebox}

\begin{attention}[Deux pièges supplémentaires]
\begin{itemize}
\item \textbf{Ne jamais conclure sur le pH seul.} Un pH de $3{,}4$ ne prouve rien à lui seul : c'est la \emph{comparaison} entre $10^{-\text{pH}}$ et $C$ qui conclut. Un acide fort très dilué peut avoir le même pH qu'un acide faible concentré !
\item \textbf{$\tau$ n'est pas une constante.} Contrairement à la constante d'acidité $K_a$ que nous étudierons dans la prochaine section, $\tau$ varie avec la concentration : il caractérise \emph{cette solution}, pas l'acide en général.
\end{itemize}
\end{attention}

\begin{savaistu}[Le pH-mètre, un « compteur » d'ions]
La mesure du pH est le moyen le plus simple de « compter » les ions \ce{H3O+} formés à l'équilibre : une électrode de verre, quelques secondes, et l'on connaît l'avancement final d'une réaction mettant en jeu des milliards de milliards de molécules ! Dans les savonneries artisanales du Cameroun, on contrôle la saponification de l'huile de palme en suivant le pH de la lessive : tant que le mélange reste très basique, la base forte \ce{NaOH} n'a pas été entièrement consommée par la réaction. Le même principe s'applique au contrôle du vinaigre, des jus de fruits comme le jus de bissap, ou des eaux de boisson : une mesure de pH, un calcul de $\tau$, et la chimie devient quantifiable.
\end{savaistu}

\begin{aretenir}[L'essentiel de la section]
\begin{itemize}
\item Le \textbf{tableau d'avancement} décrit les états initial, intermédiaire et final d'une transformation.
\item L'\textbf{avancement maximal} suppose la réaction totale : $x_{max} = C \cdot V$.
\item L'\textbf{avancement final} se déduit du pH mesuré : $x_f = 10^{-\text{pH}} \times V$.
\item Le \textbf{taux d'avancement final} $\tau = \dfrac{x_f}{x_{max}} = \dfrac{10^{-\text{pH}}}{C}$ quantifie le caractère total ou limité de la réaction.
\item $\tau = 1$ : acide (ou base) \textbf{fort(e)}, réaction totale ; $\tau < 1$ : acide (ou base) \textbf{faible}, équilibre chimique.
\item Pour un acide ou une base faible, $\tau$ s'identifie au \textbf{coefficient d'ionisation} $\alpha$.
\end{itemize}
\end{aretenir}

Nous disposons désormais d'un outil de mesure expérimental, $\tau$. Mais il dépend de la concentration, ce qui le rend peu pratique pour comparer les acides entre eux de manière absolue. Il nous faut une grandeur \emph{intrinsèque} à chaque couple acide/base : ce sera la \cle{constante d'acidité} $K_a$ et son compagnon le $pK_a$, objets de la section suivante.

\coursec{Constante d'acidité $K_a$, constante de basicité $K_b$ et $pK_a$}

Dans la section précédente, tu as appris à étudier quantitativement un équilibre chimique grâce à l'\cle{avancement} et au \cle{taux d'avancement} $\tau$. Tu as vu, par exemple, qu'une solution d'acide éthanoïque de concentration $c = \SI{1,0e-2}{mol.L^{-1}}$ présente un pH de $3,4$, ce qui correspond à un taux d'avancement d'environ $4\,\%$ seulement : l'acide réagit donc \emph{partiellement} avec l'eau. Mais $4\,\%$, est-ce beaucoup ou peu ? Par rapport à quoi ? Pour comparer entre eux tous les acides faibles (et toutes les bases faibles), il nous faut un \textbf{outil de mesure universel}, un « thermomètre de la force » des couples acide/base. Cet outil existe : c'est la \cle{constante d'acidité} $K_a$, et son compagnon logarithmique, le \cle{$pK_a$}. C'est l'objet de cette section, qui est le véritable cœur du chapitre.

\courssub{L'équilibre général d'un couple acide/base dans l'eau}

\textbf{Intuition.} Plonge un acide quelconque $\ce{AH}$ dans l'eau. Il peut céder son proton $\ce{H+}$ à une molécule d'eau, qui joue alors le rôle de base. Mais la réaction n'est jamais totale pour un acide faible : un équilibre s'établit entre les formes « acide » et « base » du couple.

\begin{definition}[Équilibre d'ionisation d'un acide dans l'eau]
Tout acide $\ce{AH}$ réagit avec l'eau selon l'équilibre :
\[ \ce{AH + H2O <=> A- + H3O+} \]
Cet équilibre met en jeu \textbf{deux couples} acide/base :
\begin{itemize}
\item le couple $\ce{AH}/\ce{A-}$ : $\ce{AH}$ cède un proton et devient sa base conjuguée $\ce{A-}$ ;
\item le couple $\ce{H3O+}/\ce{H2O}$ : l'eau capte le proton et devient l'ion oxonium $\ce{H3O+}$.
\end{itemize}
La réaction est \emph{limitée} (double flèche $\ce{<=>}$) : elle conduit à un état d'équilibre dynamique où les quatre espèces coexistent.
\end{definition}

\begin{exemplebox}[L'acide éthanoïque dans le vinaigre]
Le vinaigre consommé au Cameroun (vinaigre de vin ou vinaigre blanc utilisé pour assaisonner le poisson braisé) contient de l'acide éthanoïque \ce{CH3COOH}. Dans l'eau, il s'ionise partiellement :
\[ \ce{CH3COOH + H2O <=> CH3COO- + H3O+} \]
Les deux couples mis en jeu sont $\ce{CH3COOH}/\ce{CH3COO-}$ et $\ce{H3O+}/\ce{H2O}$. À \SI{25}{\celsius}, seule une petite fraction des molécules d'acide est ionisée : c'est un \cle{acide faible}.
\end{exemplebox}

\courssub{La constante d'acidité $K_a$}

\textbf{Intuition.} À l'équilibre, les concentrations ne bougent plus. Le quotient de réaction à l'équilibre, noté $K_a$, est une \textbf{constante} qui ne dépend que de la température. Plus l'acide « aime » céder son proton, plus l'équilibre est déplacé vers la droite, plus les produits $\ce{A-}$ et $\ce{H3O+}$ sont abondants, et plus $K_a$ est \emph{grand}.

\begin{definition}[Constante d'acidité]
La \cle{constante d'acidité} $K_a$ d'un couple $\ce{AH}/\ce{A-}$ est la constante d'équilibre de la réaction de l'acide avec l'eau :
\[ \ce{AH + H2O <=> A- + H3O+} \qquad K_a = \frac{[\ce{A-}]\,[\ce{H3O+}]}{[\ce{AH}]} \]
Les concentrations sont exprimées en \si{mol.L^{-1}} et prises \textbf{à l'équilibre}. $K_a$ est \emph{sans dimension} (grandeurs rapportées à la concentration standard $c^\circ = \SI{1}{mol.L^{-1}}$) et ne dépend que de la \textbf{température}.
\end{definition}

\begin{attention}[L'eau n'apparaît pas dans $K_a$ !]
Erreur très fréquente au baccalauréat : écrire $K_a = \dfrac{[\ce{A-}][\ce{H3O+}]}{[\ce{AH}][\ce{H2O}]}$. C'est \textbf{faux}. L'eau est le \cle{solvant} : elle est en très large excès et sa concentration reste pratiquement constante (environ $\SI{55,5}{mol.L^{-1}}$). Son \emph{activité} vaut $1$ par convention, donc $[\ce{H2O}]$ n'apparaît \textbf{jamais} dans l'expression de la constante d'équilibre. De même, les solides purs et les liquides purs n'y figurent pas.
\end{attention}

\begin{propriete}[Signification physique de $K_a$]
\begin{itemize}
\item Plus $K_a$ est \textbf{grand}, plus l'équilibre $\ce{AH + H2O <=> A- + H3O+}$ est déplacé vers la droite, plus l'acide est \textbf{fort} (fortement ionisé).
\item Plus $K_a$ est \textbf{petit}, plus l'acide est \textbf{faible} (peu ionisé).
\item Pour un acide fort (comme \ce{HCl}), la réaction avec l'eau est totale : $K_a$ est très grand et on ne le définit plus utilement en solution aqueuse (effet de nivellement).
\end{itemize}
\end{propriete}

\courssub{Le $pK_a$ : une échelle plus commode}

\textbf{Intuition.} Les valeurs de $K_a$ s'étendent sur des ordres de grandeur immenses (de $10^{7}$ à $10^{-14}$ !), exactement comme les concentrations en $\ce{H3O+}$. On fait donc le même choix que pour le pH : passer au \textbf{logarithme}.

\begin{definition}[Le $pK_a$]
On définit le \cle{$pK_a$} d'un couple acide/base par :
\[ pK_a = -\log K_a \qquad \text{soit de manière équivalente} \qquad K_a = 10^{-pK_a} \]
Le $pK_a$ est sans dimension. Plus le $pK_a$ est \textbf{petit}, plus l'acide est \textbf{fort} ; plus le $pK_a$ est \textbf{grand}, plus la base conjuguée est \textbf{forte}.
\end{definition}

\begin{exemplebox}[Le couple de l'acide éthanoïque]
Pour le couple $\ce{CH3COOH}/\ce{CH3COO-}$, à \SI{25}{\celsius} :
\[ K_a \approx 1{,}8 \times 10^{-5} \qquad pK_a = -\log(1{,}8 \times 10^{-5}) \approx 4{,}8 \]
Vérification inverse : $K_a = 10^{-4{,}8} \approx 1{,}6 \times 10^{-5}$, cohérent avec la valeur tabulée (les petites différences viennent des arrondis).
\end{exemplebox}

\begin{exemplebox}[Le couple de l'ion ammonium]
Pour le couple $\ce{NH4+}/\ce{NH3}$, $pK_a = 9{,}2$ à \SI{25}{\celsius}. On en déduit :
\[ K_a = 10^{-pK_a} = 10^{-9{,}2} \approx 6{,}3 \times 10^{-10} \]
Ce $K_a$ très petit confirme que $\ce{NH4+}$ est un acide \emph{beaucoup plus faible} que $\ce{CH3COOH}$ : il s'ionise très peu dans l'eau.
\end{exemplebox}

\begin{popfigure}
\begin{center}
\begin{tikzpicture}[x=0.72cm]
% Axe
\draw[-{Stealth},thick,popdark] (-2.5,0) -- (15.5,0) node[below left=1pt and 0pt,font=\small] {$pK_a$};
\foreach \x in {-2,-1,0,1,2,3,4,5,6,7,8,9,10,11,12,13,14}{
  \draw[popdark] (\x,0.08) -- (\x,-0.08) node[below,font=\scriptsize] {\x};
}
% Zones
\fill[popblue!15] (-2.5,0.25) rectangle (0,1.55);
\fill[popgreen!12] (0,0.25) rectangle (14,1.55);
\node[font=\scriptsize\itshape,text=popblue] at (-1.2,1.85) {acides forts};
\node[font=\scriptsize\itshape,text=popgreen!60!black] at (7,1.85) {acides faibles / bases faibles};
% Couples
\draw[popblue,very thick] (-1.5,0) -- (-1.5,1.2);
\node[above,font=\scriptsize,text=popblue,align=center] at (-1.5,1.2) {$\ce{HCl}/\ce{Cl-}$\\ $pK_a<0$};
\draw[popink,very thick] (0,0) -- (0,1.2);
\node[above,font=\scriptsize,text=popink,align=center] at (0.6,1.35) {$\ce{H3O+}/\ce{H2O}$\\ $pK_a=0$};
\draw[poporange,very thick] (4.8,0) -- (4.8,1.2);
\node[above,font=\scriptsize,text=poporange,align=center] at (4.8,1.2) {$\ce{CH3COOH}/\ce{CH3COO-}$\\ $pK_a=4{,}8$};
\draw[poppurple,very thick] (9.2,0) -- (9.2,1.2);
\node[above,font=\scriptsize,text=poppurple,align=center] at (9.2,1.2) {$\ce{NH4+}/\ce{NH3}$\\ $pK_a=9{,}2$};
\draw[popgold!80!black,very thick] (14,0) -- (14,1.2);
\node[above,font=\scriptsize,text=popgold!80!black,align=center] at (13.4,1.35) {$\ce{H2O}/\ce{HO-}$\\ $pK_a=14$};
% Flèches de lecture
\draw[-{Stealth},thick,popink] (-2.2,-0.9) -- (3,-0.9) node[midway,below,font=\scriptsize,text=popink] {force de l'acide croissante};
\draw[-{Stealth},thick,poppurple] (14.8,-1.7) -- (9,-1.7) node[midway,below,font=\scriptsize,text=poppurple] {force de la base conjuguée croissante};
\end{tikzpicture}
\end{center}
\legende{Échelle des $pK_a$ de quelques couples usuels à \SI{25}{\celsius}. Plus on se déplace vers la gauche, plus l'acide est fort ; plus on se déplace vers la droite, plus la base conjuguée est forte.}
\end{popfigure}

\courssub{La constante de basicité $K_b$}

\textbf{Intuition.} De la même façon qu'on mesure la force d'un acide par sa réaction \emph{avec l'eau}, on peut mesurer la force d'une base $\ce{A-}$ par sa réaction avec l'eau, où elle \emph{capture} un proton.

\begin{definition}[Constante de basicité]
La \cle{constante de basicité} $K_b$ d'un couple $\ce{AH}/\ce{A-}$ est la constante d'équilibre de la réaction de la base avec l'eau :
\[ \ce{A- + H2O <=> AH + HO-} \qquad K_b = \frac{[\ce{AH}]\,[\ce{HO-}]}{[\ce{A-}]} \]
On définit aussi $pK_b = -\log K_b$. Plus $K_b$ est grand (plus $pK_b$ est petit), plus la base est \textbf{forte}.
\end{definition}

\begin{exemplebox}[L'ion éthanoate, base faible]
L'ion éthanoate $\ce{CH3COO-}$, présent dans une solution d'éthanoate de sodium, réagit partiellement avec l'eau :
\[ \ce{CH3COO- + H2O <=> CH3COOH + HO-} \]
Cet équilibre, très peu déplacé vers la droite, rend la solution \emph{légèrement basique} : c'est pourquoi le pH d'une solution d'éthanoate de sodium est supérieur à $7$.
\end{exemplebox}

\courssub{Relation fondamentale entre $K_a$, $K_b$ et $K_e$}

Multiplions membre à membre les expressions de $K_a$ et $K_b$ d'un \emph{même} couple $\ce{AH}/\ce{A-}$ :
\[ K_a \times K_b = \frac{[\ce{A-}]\,[\ce{H3O+}]}{[\ce{AH}]} \times \frac{[\ce{AH}]\,[\ce{HO-}]}{[\ce{A-}]} = [\ce{H3O+}]\,[\ce{HO-}] = K_e \]

\begin{aretenir}[Relation fondamentale]
Pour tout couple acide/base $\ce{AH}/\ce{A-}$, à \SI{25}{\celsius} :
\[ K_a \times K_b = K_e = 10^{-14} \qquad \text{soit} \qquad \boxed{pK_a + pK_b = 14} \]
Conséquence immédiate : un acide fort (grand $K_a$, petit $pK_a$) possède une base conjuguée \textbf{très faible} (petit $K_b$, grand $pK_b$), et réciproquement. On dit que la base conjuguée d'un acide fort est \emph{indifférente} à l'eau (comme $\ce{Cl-}$, spectatrice).
\end{aretenir}

\begin{exemplebox}[Application au couple $\ce{NH4+}/\ce{NH3}$]
Avec $pK_a = 9{,}2$, on obtient pour l'ammoniac :
\[ pK_b(\ce{NH3}) = 14 - 9{,}2 = 4{,}8 \qquad K_b = 10^{-4{,}8} \approx 1{,}6 \times 10^{-5} \]
L'ammoniac est donc une base faible, mais nettement plus forte que l'ion éthanoate dont $pK_b = 14 - 4{,}8 = 9{,}2$.
\end{exemplebox}

\courssub{Les couples de l'eau : deux bornes de l'échelle}

L'eau est \cle{amphotère} : elle intervient dans \textbf{deux couples}, qui délimitent l'échelle des $pK_a$ en solution aqueuse.

\begin{propriete}[Les deux couples de l'eau]
\begin{itemize}
\item Couple $\ce{H3O+}/\ce{H2O}$ : $\ce{H3O+ + H2O <=> H2O + H3O+}$, avec $K_a = 1$ donc $pK_a = 0$. L'ion oxonium est \textbf{l'acide le plus fort} pouvant exister dans l'eau : tout acide plus fort (comme \ce{HCl}) est intégralement converti en $\ce{H3O+}$.
\item Couple $\ce{H2O}/\ce{HO-}$ : $\ce{H2O + H2O <=> HO- + H3O+}$, avec $K_a = K_e = 10^{-14}$ donc $pK_a = 14$. L'ion hydroxyde est \textbf{la base la plus forte} pouvant exister dans l'eau.
\end{itemize}
Tout couple acide faible / base faible a donc un $pK_a$ compris entre $0$ et $14$.
\end{propriete}

\courssub{Identifier les couples dans une équation d'équilibre}

\begin{methode}[Repérer les couples acide/base]
Pour identifier les deux couples mis en jeu dans une équation acido-basique :
\begin{enumerate}
\item Cherche les \textbf{deux espèces qui ne diffèrent que par un proton} $\ce{H+}$ : elles forment un couple (forme acide / forme basique).
\item Fais de même avec les deux autres espèces : c'est le second couple.
\item Le proton est \textbf{cédé} par l'acide du couple 1 et \textbf{capté} par la base du couple 2 : la réaction acido-basique est un \emph{transfert de proton}.
\end{enumerate}
\end{methode}

\begin{exemplebox}[Réaction de l'acide éthanoïque avec l'ammoniac]
\[ \ce{CH3COOH + NH3 <=> CH3COO- + NH4+} \]
\begin{itemize}
\item $\ce{CH3COOH}$ et $\ce{CH3COO-}$ ne diffèrent que d'un $\ce{H+}$ : couple $\ce{CH3COOH}/\ce{CH3COO-}$ ($pK_{a1} = 4{,}8$).
\item $\ce{NH4+}$ et $\ce{NH3}$ ne diffèrent que d'un $\ce{H+}$ : couple $\ce{NH4+}/\ce{NH3}$ ($pK_{a2} = 9{,}2$).
\item Le proton passe de $\ce{CH3COOH}$ (acide 1) vers $\ce{NH3}$ (base 2).
\end{itemize}
\end{exemplebox}

\begin{popfigure}
\begin{center}
\begin{tikzpicture}[font=\small]
% Couple 1
\node[draw=popink,thick,rounded corners=3pt,fill=popink!8,minimum width=2.6cm,minimum height=1cm,align=center] (a1) at (0,2) {\textbf{acide 1}\\ $\ce{CH3COOH}$};
\node[draw=popink,thick,rounded corners=3pt,fill=popink!8,minimum width=2.6cm,minimum height=1cm,align=center] (b1) at (7,2) {\textbf{base 1}\\ $\ce{CH3COO-}$};
% Couple 2
\node[draw=poppurple,thick,rounded corners=3pt,fill=poppurple!8,minimum width=2.6cm,minimum height=1cm,align=center] (b2) at (0,0) {\textbf{base 2}\\ $\ce{NH3}$};
\node[draw=poppurple,thick,rounded corners=3pt,fill=poppurple!8,minimum width=2.6cm,minimum height=1cm,align=center] (a2) at (7,0) {\textbf{acide 2}\\ $\ce{NH4+}$};
% Transfert de proton
\draw[-{Stealth},very thick,poporange,decorate,decoration={snake,amplitude=1.2pt,segment length=6pt}] (a1.south) -- (b2.north) node[midway,left=2pt,text=poporange,font=\small\bfseries] {$\ce{H+}$};
% Liaisons conjuguées
\draw[{Stealth}-{Stealth},thick,popink,dashed] (a1.east) -- (b1.west) node[midway,above,font=\scriptsize,text=popink] {couple 1 : $pK_a = 4{,}8$};
\draw[{Stealth}-{Stealth},thick,poppurple,dashed] (b2.east) -- (a2.west) node[midway,below,font=\scriptsize,text=poppurple] {couple 2 : $pK_a = 9{,}2$};
\node[font=\scriptsize\itshape,text=popmuted] at (3.5,-1) {$\ce{CH3COOH + NH3 <=> CH3COO- + NH4+}$};
\end{tikzpicture}
\end{center}
\legende{Transfert de proton $\ce{H+}$ entre les deux couples d'une réaction acido-basique : l'acide 1 cède son proton à la base 2.}
\end{popfigure}

\begin{exoresolu}[Calculer et exploiter $K_a$, $K_b$ et $pK_a$]
\textbf{Énoncé.} L'acide méthanoïque \ce{HCOOH} est présent dans le venin des fourmis. À \SI{25}{\celsius}, une solution de concentration $c = \SI{1,0e-2}{mol.L^{-1}}$ a un pH de $2{,}9$.
\begin{enumerate}
\item Écrire l'équilibre d'ionisation et identifier les couples.
\item Calculer $K_a$ et $pK_a$ du couple.
\item En déduire $pK_b$ de l'ion méthanoate.
\end{enumerate}
\tcblower
\textbf{Solution détaillée.}
\begin{enumerate}
\item Équilibre : $\ce{HCOOH + H2O <=> HCOO- + H3O+}$. Couples : $\ce{HCOOH}/\ce{HCOO-}$ et $\ce{H3O+}/\ce{H2O}$.
\item On détermine les concentrations à l'équilibre à l'aide d'un tableau d'avancement.
\begin{center}
\begin{tabularx}{\linewidth}{|c|c|c|c|}\hline
 & $\ce{HCOOH}$ & $\ce{H2O}$ & $\ce{HCOO-}$ + $\ce{H3O+}$ \\\hline
État initial & $c = \SI{1,0e-2}{M}$ & excès & $0$ \\\hline
État final & $c - x$ & excès & $x$ \\\hline
\end{tabularx}
\end{center}
Le pH donne $x = [\ce{H3O+}] = 10^{-\mathrm{pH}} = 10^{-2{,}9} \approx \SI{1,26e-3}{mol.L^{-1}}$.
Donc $[\ce{HCOO-}] = x = \SI{1,26e-3}{M}$ et $[\ce{HCOOH}] = c - x = \SI{1,00e-2}{M} - \SI{1,26e-3}{M} = \SI{8,74e-3}{M}$.
\[ K_a = \frac{[\ce{HCOO-}][\ce{H3O+}]}{[\ce{HCOOH}]} = \frac{(\SI{1,26e-3}{})^2}{\SI{8,74e-3}{}} \approx \SI{1,8e-4}{} \]
d'où $pK_a = -\log(\SI{1,8e-4}{}) \approx 3{,}7$. L'acide méthanoïque est un acide faible, légèrement plus fort que l'acide éthanoïque ($pK_a = 4{,}8$).
\item $pK_b(\ce{HCOO-}) = 14 - pK_a = 14 - 3{,}7 = 10{,}3$.
\end{enumerate}
\end{exoresolu}

\begin{attention}[Pièges à éviter]
\begin{itemize}
\item Ne confonds pas $K_a$ et $pK_a$ : un \emph{grand} $K_a$ correspond à un \emph{petit} $pK_a$ (acide fort). C'est contre-intuitif au début : entraîne-toi à convertir dans les deux sens avec $K_a = 10^{-pK_a}$.
\item La relation $pK_a + pK_b = 14$ n'est valable qu'à \SI{25}{\celsius} (car $pK_e = 14$ à cette température).
\item $K_a$ et $K_b$ d'un même couple ne décrivent pas la même réaction : $K_a$ concerne l'acide avec l'eau, $K_b$ la base avec l'eau.
\end{itemize}
\end{attention}

\begin{savaistu}[Le sang, une solution tamponnée grâce au $pK_a$]
Le pH du sang humain doit rester compris entre $7{,}35$ et $7{,}45$ : au-delà, la vie est en danger. Cette stabilité remarquable est assurée par le couple $\ce{H2CO3}/\ce{HCO3-}$ (acide carbonique / ion hydrogénocarbonate), dont le $pK_a$ apparent vaut $6{,}4$ à \SI{25}{\celsius} (proche du pH sanguin). Quand un effort musculaire produit de l'acide lactique, les ions $\ce{HCO3-}$ captent les protons en excès ; quand le sang devient trop basique, $\ce{H2CO3}$ libère des protons. Ce mécanisme de \emph{solution tampon}, que tu étudieras en détail dans la suite du module, est une application directe de la relation entre pH et $pK_a$. Les plantes médicinales de la pharmacopée camerounaise contiennent aussi de nombreux acides et bases faibles dont l'activité dépend de leur $pK_a$ !
\end{savaistu}

\begin{aretenir}[L'essentiel de la section]
\begin{itemize}
\item Constante d'acidité : $K_a = \dfrac{[\ce{A-}][\ce{H3O+}]}{[\ce{AH}]}$ pour l'équilibre $\ce{AH + H2O <=> A- + H3O+}$ ; l'eau, solvant, n'y figure pas.
\item $pK_a = -\log K_a$ et $K_a = 10^{-pK_a}$. Plus $pK_a$ est petit, plus l'acide est fort ; plus $pK_a$ est grand, plus la base conjuguée est forte.
\item Constante de basicité : $K_b = \dfrac{[\ce{AH}][\ce{HO-}]}{[\ce{A-}]}$, avec $K_a \times K_b = K_e = 10^{-14}$, soit $pK_a + pK_b = 14$ à \SI{25}{\celsius}.
\item Couples de l'eau : $\ce{H3O+}/\ce{H2O}$ ($pK_a = 0$) et $\ce{H2O}/\ce{HO-}$ ($pK_a = 14$) : ils bornent l'échelle des $pK_a$ en solution aqueuse.
\item Pour identifier les couples d'un équilibre, repère les espèces qui ne diffèrent que d'un proton $\ce{H+}$.
\end{itemize}
\end{aretenir}

\coursec{Classification des couples et prévision des réactions}

Dans la section précédente, nous avons appris à associer à chaque couple acide/base deux nombres remarquables : la constante d'acidité $K_a$ et le $pK_a$. Ces nombres ne sont pas de simples étiquettes : ils mesurent la \emph{force} de l'acide et de sa base conjuguée. Nous allons maintenant exploiter cette information pour répondre à deux questions essentielles au Baccalauréat TI : \emph{comment classer les couples entre eux ?} et \emph{quand on mélange un acide et une base, dans quel sens la réaction se produit-elle, et jusqu'où va-t-elle ?} C'est toute la puissance prédictive de la chimie acido-basique.

\courssub{Classement des couples sur une échelle de $pK_a$}

\textbf{Intuition.} Imagine une compétition : chaque couple acide/base est un « joueur » dont la force est mesurée par son $pK_a$. Pour comparer deux acides, il suffit de comparer les scores de leurs couples. Mais attention au sens du score : ici, \emph{le plus petit gagne}.

\begin{definition}[Classement des couples par les $pK_a$]
Deux couples acide/base peuvent être classés par la valeur de leur $pK_a$ (ou de leur $K_a$). L'\cle{acide le plus fort} appartient au couple de \cle{plus petit $pK_a$} (c'est-à-dire de plus grand $K_a$). Réciproquement, la \cle{base la plus forte} appartient au couple de \cle{plus grand $pK_a$}.
\end{definition}

\begin{propriete}[Force et $pK_a$]
\begin{itemize}
\item Plus le $pK_a$ d'un couple est \textbf{petit}, plus l'\textbf{acide} du couple est \textbf{fort} (il cède facilement son proton) et plus sa \textbf{base conjuguée} est \textbf{faible} (elle a peu tendance à capter un proton : on dit qu'elle est \emph{indifférente} à l'eau).
\item Plus le $pK_a$ d'un couple est \textbf{grand}, plus la \textbf{base} du couple est \textbf{forte} et plus l'acide conjugué est faible.
\end{itemize}
\end{propriete}

\textbf{Cas limites : acides forts, bases fortes et effet de nivellement.}
\begin{itemize}
\item Les \cle{acides forts} (\ce{HCl}, \ce{HNO3}, \ce{H2SO4} (premier proton), ion \ce{H3O+}) sont totalement dissociés dans l'eau : leur couple a un $pK_a < 0$ (pour \ce{HCl}/\ce{Cl-}, $pK_a \approx -7$ ; pour \ce{H3O+}/\ce{H2O}, $pK_a = 0$). Leur base conjuguée (\ce{Cl-}, \ce{NO3-}, \ce{H2O}) est \emph{spectatrice} : elle ne capte aucun proton dans l'eau.
\item Les \cle{bases fortes} (ion \ce{HO-}, et des bases plus fortes comme l'ion éthanolate \ce{C2H5O-} ou l'amide \ce{NH2-}) réagissent totalement avec l'eau. Leur couple a un $pK_a \ge 14$ (pour \ce{H2O}/\ce{HO-}, $pK_a = 14$ ; pour \ce{C2H5OH}/\ce{C2H5O-}, $pK_a \approx 16$). L'acide conjugué est alors extrêmement faible.
\item \textbf{Effet de nivellement :} Dans l'eau, aucun acide ne peut être plus fort que \ce{H3O+} (il donne son proton à l'eau) et aucune base ne peut être plus forte que \ce{HO-} (elle prend un proton à l'eau). L'eau « nivelle » les forces : \ce{HCl} et \ce{HNO3} apparaissent de force égale (totale) ; \ce{C2H5O-} et \ce{NH2-} réagissent instantanément avec l'eau pour donner \ce{HO-}.
\item Entre $pK_a = 0$ et $pK_a = 14$ se situent tous les \textbf{acides faibles} et les \textbf{bases faibles} : c'est le domaine des équilibres, celui du vinaigre, de l'ammoniac, du jus de bissap.
\end{itemize}

\begin{attention}[Le piège du sens de l'échelle]
L'erreur la plus fréquente au Bac est d'inverser l'échelle : croire qu'un \emph{grand} $pK_a$ signifie un acide \emph{fort}. C'est l'inverse ! Un \textbf{petit} $pK_a$ correspond à un acide \textbf{fort}. Astuce mnémotechnique : le $pK_a$ mesure la force de la \emph{base} du couple ; l'acide, lui, joue « à l'envers ». Vérifie toujours avec un exemple connu : \ce{HCl} est un acide fort et son $pK_a$ vaut environ $-7$, donc petit $pK_a$ = acide fort.
\end{attention}

\begin{popfigure}
\begin{center}
\begin{tikzpicture}[scale=1.0]
% Axe vertical des pKa
\draw[->, thick, popdark] (0,-0.6) -- (0,15.2) node[above, font=\small\bfseries] {$pK_a$};
\foreach \y in {0,2,4,6,8,10,12,14} {
  \draw[popdark] (-0.12,\y) -- (0.12,\y) node[left=2pt, font=\small] {\y};
}
% Couples : acide à gauche, base à droite (valeurs standards arrondies)
\node[font=\small, text=popink] at (-1.9,0) {\ce{H3O+}};
\node[font=\small, text=popblue] at (1.9,0) {\ce{H2O}};
\node[font=\small, text=popink] at (-1.9,3.7) {\ce{HCOOH}};
\node[font=\small, text=popblue] at (1.9,3.7) {\ce{HCOO-}};
\node[font=\small, text=popink] at (-1.9,4.8) {\ce{CH3COOH}};
\node[font=\small, text=popblue] at (1.9,4.8) {\ce{CH3COO-}};
\node[font=\small, text=popink] at (-1.9,6.4) {\ce{H2CO3}};
\node[font=\small, text=popblue] at (1.9,6.4) {\ce{HCO3-}};
\node[font=\small, text=popink] at (-1.9,9.2) {\ce{NH4+}};
\node[font=\small, text=popblue] at (1.9,9.2) {\ce{NH3}};
\node[font=\small, text=popink] at (-1.9,14) {\ce{H2O}};
\node[font=\small, text=popblue] at (1.9,14) {\ce{HO-}};
% Lignes pointillées reliant les couples
\foreach \y in {0,3.7,4.8,6.4,9.2,14} {
  \draw[dotted, popmuted] (-1.3,\y) -- (1.3,\y);
}
% Flèches de force
\draw[->, very thick, popink] (-3.6,6) -- (-3.6,0.4) node[midway, left, font=\small\bfseries, align=center] {force de\\l'acide};
\draw[->, very thick, popblue] (3.6,8) -- (3.6,13.6) node[midway, right, font=\small\bfseries, align=center] {force de\\la base};
% Règle du gamma : exemple CH3COOH + NH3
\draw[->, very thick, poporange, rounded corners=4pt]
  (-1.9,4.8) -- (-1.9,7.0) -- (1.9,7.0) -- (1.9,9.2);
\node[font=\small\bfseries, text=poporange, align=center] at (0,8.1) {réaction\\prépondérante};
\end{tikzpicture}
\end{center}
\legende{Échelle verticale des $pK_a$ (valeurs usuelles) : les acides à gauche, les bases conjuguées à droite. La flèche orange illustre la règle du gamma pour le mélange \ce{CH3COOH} + \ce{NH3} : la réaction part de l'acide le plus fort (en bas) vers la base la plus forte (en haut).}
\end{popfigure}

\begin{popfigure}
\begin{center}
\begin{tikzpicture}
\node[font=\small\bfseries, text=white, fill=popink, minimum width=3.4cm, minimum height=0.8cm] (a1) at (0,0) {Acides forts};
\node[font=\small, fill=poppinkL, minimum width=3.4cm, minimum height=1.5cm, anchor=north] at (a1.south) {\shortstack{$pK_a < 0$\\ \ce{HCl}, \ce{HNO3}, \ce{H3O+}\\ réaction totale avec l'eau}};
\node[font=\small\bfseries, text=white, fill=poporange, minimum width=3.4cm, minimum height=0.8cm] (a2) at (4.1,0) {Acides faibles};
\node[font=\small, fill=poporangeL, minimum width=3.4cm, minimum height=1.5cm, anchor=north] at (a2.south) {\shortstack{$0 < pK_a < 14$\\ \ce{CH3COOH}, \ce{NH4+}, \ce{H2CO3}\\ équilibre avec l'eau}};
\node[font=\small\bfseries, text=white, fill=popblue, minimum width=3.4cm, minimum height=0.8cm] (b1) at (8.2,0) {Bases faibles};
\node[font=\small, fill=popblueL, minimum width=3.4cm, minimum height=1.5cm, anchor=north] at (b1.south) {\shortstack{$0 < pK_a < 14$\\ \ce{CH3COO-}, \ce{NH3}, \ce{HCO3-}\\ équilibre avec l'eau}};
\node[font=\small\bfseries, text=white, fill=popgreen, minimum width=3.4cm, minimum height=0.8cm] (b2) at (12.3,0) {Bases fortes};
\node[font=\small, fill=popgreenL, minimum width=3.4cm, minimum height=1.5cm, anchor=north] at (b2.south) {\shortstack{$pK_a \ge 14$\\ \ce{HO-}, \ce{C2H5O-}*\\ *nivelées en \ce{HO-} dans l'eau}};
\end{tikzpicture}
\end{center}
\legende{Classification des couples acide/base selon la position du $pK_a$ par rapport aux bornes 0 et 14 du couple de l'eau.}
\end{popfigure}

\courssub{Prévision du sens d'une réaction acido-basique : la règle du gamma}

\textbf{Intuition.} Un proton, c'est un peu comme un objet convoité : il va toujours vers celui qui l'attire le plus fortement. Quand on mélange un acide $A_1$ et une base $B_2$, le proton quitte l'acide qui le retient le moins (l'acide le plus fort) pour rejoindre la base qui l'attire le plus (la base la plus forte).

\begin{propriete}[Réaction prépondérante]
Dans un mélange contenant plusieurs acides et plusieurs bases, la \cle{réaction prépondérante} est celle qui a lieu entre \textbf{l'acide le plus fort} (couple de plus petit $pK_a$) et \textbf{la base la plus forte} (couple de plus grand $pK_a$). Elle produit l'acide le plus faible et la base la plus faible.
\end{propriete}

\begin{methode}[Prévoir le sens d'une réaction : la règle du gamma]
\begin{enumerate}
\item \textbf{Identifier} les deux couples mis en jeu et écrire leurs $pK_a$.
\item \textbf{Placer} les deux couples sur une échelle verticale de $pK_a$ croissants vers le haut, acides à gauche, bases à droite.
\item \textbf{Entourer} les espèces effectivement introduites dans le mélange (l'acide d'un couple, la base de l'autre).
\item \textbf{Tracer} le « gamma » ($\gamma$) : partir de l'acide introduit, monter verticalement, tourner à angle droit vers la base introduite. Si le gamma se lit dans le sens « acide en bas, base en haut », la réaction écrite dans ce sens est la réaction prépondérante.
\item \textbf{Écrire} l'équation : l'acide cède son proton à la base ; chaque espèce donne sa forme conjuguée. On écrit d'abord une double flèche $\ce{<=>}$, puis on détermine $K$ pour savoir si on peut la remplacer par une flèche simple $\ce{->}$.
\end{enumerate}
\end{methode}

\begin{exemplebox}[Mélange acide éthanoïque + ammoniac]
On mélange une solution d'acide éthanoïque (couple \ce{CH3COOH}/\ce{CH3COO-}, $pK_{a1} = 4,8$) et une solution d'ammoniac (couple \ce{NH4+}/\ce{NH3}, $pK_{a2} = 9,2$). L'acide le plus fort est \ce{CH3COOH} (plus petit $pK_a$), la base la plus forte est \ce{NH3} (couple de plus grand $pK_a$). Le gamma part de \ce{CH3COOH} et monte vers \ce{NH3} : la réaction prépondérante s'écrit :
\[ \ce{CH3COOH + NH3 <=> CH3COO- + NH4+} \]
Transfert d'un proton de l'acide éthanoïque vers l'ammoniac ; réaction rapide, à température ambiante, sans catalyseur. Le calcul de $K$ (voir plus bas) montrera qu'elle est quasi totale.
\end{exemplebox}

\courssub{Constante d'équilibre d'une réaction acido-basique}

Savoir dans quel sens la réaction se produit ne suffit pas : il faut aussi savoir \emph{jusqu'où} elle va. Considérons la réaction générale entre l'acide $A_1$ du couple $A_1/B_1$ ($pK_{a1}$) et la base $B_2$ du couple $A_2/B_2$ ($pK_{a2}$) :
\[ \ce{A_1 + B_2 <=> B_1 + A_2} \qquad K = \frac{[B_1][A_2]}{[A_1][B_2]} \]
En multipliant numérateur et dénominateur par $[\ce{H3O+}]$, on fait apparaître le rapport des deux constantes d'acidité :
\[ K = \frac{K_{a1}}{K_{a2}} = 10^{pK_{a2} - pK_{a1}} \]

\begin{aretenir}[Constante d'équilibre et écart de $pK_a$]
Pour la réaction $\ce{A_1 + B_2 <=> B_1 + A_2}$ :
\[ \boxed{K = 10^{\,pK_a(\text{acide produit}) - pK_a(\text{acide réactif})}} \]
\begin{itemize}
\item Si $K > 10^4$ (écart de $pK_a$ supérieur à 4 unités) : la réaction est \textbf{quasi totale} (on utilise une flèche simple $\ce{->}$).
\item Si $K < 10^{-4}$ : la réaction est \textbf{très limitée}, on considère qu'elle n'a pas lieu.
\item Entre les deux : c'est un \textbf{équilibre} (double flèche $\ce{<=>}$), il faut utiliser un tableau d'avancement pour calculer la composition finale.
\end{itemize}
\end{aretenir}

\begin{exoresolu}[Réaction entre l'acide éthanoïque et l'ammoniac]
\textbf{Énoncé.} On mélange \SI{20,0}{mL} d'une solution d'acide éthanoïque à \SI{0,10}{mol.L^{-1}} et \SI{20,0}{mL} d'une solution d'ammoniac à \SI{0,10}{mol.L^{-1}}. Données : $pK_a(\ce{CH3COOH}/\ce{CH3COO-}) = 4,8$ ; $pK_a(\ce{NH4+}/\ce{NH3}) = 9,2$.
\begin{enumerate}
\item Écrire l'équation de la réaction prépondérante.
\item Calculer sa constante d'équilibre $K$ et conclure sur le caractère total ou non de la réaction.
\item Donner la composition finale du mélange (quantités de matière).
\end{enumerate}
\tcblower
\textbf{Solution détaillée.}
\begin{enumerate}
\item L'acide le plus fort est \ce{CH3COOH} ($pK_a = 4,8$), la base la plus forte est \ce{NH3} (couple de $pK_a = 9,2$). La règle du gamma donne :
\[ \ce{CH3COOH + NH3 <=> CH3COO- + NH4+} \]
Transfert d'un proton de l'acide éthanoïque vers l'ammoniac ; réaction rapide, à température ambiante, sans catalyseur.
\item L'acide produit est \ce{NH4+} ($pK_a = 9,2$), l'acide réactif est \ce{CH3COOH} ($pK_a = 4,8$) :
\[ K = 10^{9,2 - 4,8} = 10^{4,4} \approx 2,5 \times 10^{4} \]
Comme $K > 10^{4}$, la réaction est \textbf{quasi totale} : pratiquement tout l'acide éthanoïque et tout l'ammoniac introduits sont consommés. On peut écrire l'équation avec une flèche simple $\ce{->}$.
\item \textbf{Tableau d'avancement (réaction totale)} :
\begin{center}
\begin{tabularx}{\linewidth}{|c|X|X|X|X|}\hline
 & \ce{CH3COOH} & \ce{NH3} & \ce{CH3COO-} & \ce{NH4+} \\\hline
État initial & $n_0 = 2,0 \times 10^{-3}$ & $n_0 = 2,0 \times 10^{-3}$ & 0 & 0 \\\hline
État final & 0 & 0 & $2,0 \times 10^{-3}$ & $2,0 \times 10^{-3}$ \\\hline
\end{tabularx}
\end{center}
La solution obtenue contient essentiellement des ions éthanoate \ce{CH3COO-} et ammonium \ce{NH4+} (c'est une solution d'éthanoate d'ammonium) dans un volume total de \SI{40,0}{mL}. Leurs concentrations formelles sont \SI{0,050}{mol.L^{-1}}.
\end{enumerate}
\end{exoresolu}

\courssub{Lien avec la force des acides et des bases}

La règle du gamma éclaire d'un jour nouveau ce que nous avons appris sur les acides forts et faibles :

\begin{propriete}[Conséquences pratiques]
\begin{itemize}
\item Un \textbf{acide fort} ($pK_a < 0$) réagit de façon \textbf{totale avec toute base}, même très faible, car l'écart de $pK_a$ est toujours énorme. C'est pourquoi le dosage d'un acide fort par la soude donne un saut de pH franc.
\item Un \textbf{acide faible} ne réagit \textbf{totalement} qu'avec une base \textbf{suffisamment forte}, c'est-à-dire appartenant à un couple dont le $pK_a$ dépasse le sien d'au moins 4 unités. Avec une base trop faible, on obtient un simple équilibre.
\item Symétriquement, une \textbf{base forte} ($pK_a \ge 14$) capte le proton de \textbf{tout acide}, même très faible, y compris l'eau elle-même : c'est pourquoi on ne peut pas préparer une solution aqueuse d'ion éthanolate \ce{C2H5O-} sans qu'il réagisse totalement avec l'eau (effet de nivellement).
\end{itemize}
\end{propriete}

\begin{exemplebox}[Le vinaigre attaque-t-il le calcaire ?]
Le calcaire des bouilloires est du carbonate de calcium \ce{CaCO3}. L'ion carbonate est la base du couple \ce{HCO3-}/\ce{CO3^2-} ($pK_{a2} = 10,3$) et l'ion hydrogénocarbonate est la base du couple \ce{H2CO3}/\ce{HCO3-} ($pK_{a1} = 6,4$). L'acide éthanoïque du vinaigre ($pK_a = 4,8$) est plus fort que \ce{H2CO3} :
\begin{align*}
\text{1ère étape :}\quad & \ce{CH3COOH + CO3^2- -> CH3COO- + HCO3-} \qquad K_1 = 10^{10,3-4,8} = 10^{5,5} \gg 10^{4} \quad (\text{total}) \\
\text{2ème étape :}\quad & \ce{CH3COOH + HCO3- <=> CH3COO- + H2CO3} \qquad K_2 = 10^{6,4-4,8} = 10^{1,6} \approx 40 \quad (\text{équilibre})
\end{align*}
L'acide carbonique \ce{H2CO3} formé se décompose en \ce{CO2} gazeux (effervescence) et eau : $\ce{H2CO3 -> CO2 ^ + H2O}$. Cette évasion gazeuse déplace le second équilibre vers l'avant (principe de Le Chatelier), ce qui rend le détartrage globalement complet.
\end{exemplebox}

\begin{savaistu}[Pourquoi le vinaigre détartre-t-il ?]
C'est précisément parce que l'acide éthanoïque ($pK_a = 4,8$) est \emph{plus fort} que l'acide carbonique ($pK_{a1} = 6,4$) que le vinaigre fait effervescer le calcaire : la règle du gamma impose le transfert du proton vers le carbonate, et le \ce{CO2} qui se dégage chasse l'équilibre vers la dissolution complète du dépôt. À la maison comme au laboratoire, un détartrant n'est jamais qu'un acide dont le $pK_a$ est bien choisi : assez petit pour attaquer le calcaire, assez grand pour rester maniable sans danger. Les détartrants industriels utilisent souvent l'acide sulfamique ($pK_a \approx 1,0$) ou l'acide citrique ($pK_{a1} \approx 3,1$), pour le même type de raisonnement.
\end{savaistu}

\begin{attention}[Deux pièges classiques au Bac]
\begin{itemize}
\item \textbf{Ne pas confondre} $K$ de la réaction avec $K_a$ d'un couple : $K = 10^{\Delta pK_a}$ concerne une \emph{réaction entre deux couples}, alors que $K_a$ caractérise \emph{un seul couple} face à l'eau.
\item \textbf{Ne pas conclure trop vite} « réaction totale » : vérifie toujours le critère quantitatif $K > 10^{4}$. Un écart de $pK_a$ de 2 unités donne $K = 100$ : la réaction est favorisée mais \emph{pas} totale ; il faut alors un tableau d'avancement avec l'avancement final $x_f$ solution d'une équation du second degré.
\item \textbf{Attention aux couples amphotères} : \ce{HCO3-} peut agir comme acide (couple \ce{HCO3-}/\ce{CO3^2-}, $pK_a = 10,3$) ou comme base (couple \ce{H2CO3}/\ce{HCO3-}, $pK_a = 6,4$). Identifie bien le rôle joué dans le mélange.
\end{itemize}
\end{attention}

\begin{exercice}[À toi de jouer]
On considère les couples \ce{HCOOH}/\ce{HCOO-} ($pK_a = 3,7$) et \ce{CH3COOH}/\ce{CH3COO-} ($pK_a = 4,8$).
\begin{enumerate}
\item Quel est l'acide le plus fort ? La base la plus forte ?
\item Écrire l'équation de la réaction entre l'acide méthanoïque \ce{HCOOH} et l'ion éthanoate \ce{CH3COO-}.
\item Calculer la constante $K$ de cette réaction. Est-elle totale ?
\end{enumerate}
\end{exercice}

\begin{corrige}[Exercice : acide méthanoïque et ion éthanoate]
\begin{enumerate}
\item L'acide le plus fort est celui du couple de plus petit $pK_a$ : \ce{HCOOH} ($3,7 < 4,8$). La base la plus forte est la base du couple de plus grand $pK_a$ : \ce{CH3COO-}.
\item Règle du gamma : l'acide \ce{HCOOH} cède son proton à la base \ce{CH3COO-} :
\[ \ce{HCOOH + CH3COO- <=> HCOO- + CH3COOH} \]
\item $K = 10^{pK_a(\text{produit}) - pK_a(\text{réactif})} = 10^{4,8 - 3,7} = 10^{1,1} \approx 12,6$. Comme $10^{-4} < K < 10^{4}$, la réaction est \textbf{limitée} : c'est un équilibre, les quatre espèces coexistent en proportions comparables. Une flèche simple serait ici une erreur.
\end{enumerate}
\end{corrige}

\begin{aretenir}[L'essentiel de la section]
\begin{itemize}
\item Les couples se classent sur une échelle de $pK_a$ : petit $pK_a$ = acide fort (base conjuguée faible) ; grand $pK_a$ = base forte (acide conjugué faible).
\item Acides forts : $pK_a < 0$ ; bases fortes : $pK_a \ge 14$ (nivelées par l'eau) ; entre les deux, le royaume des acides et bases faibles.
\item La réaction prépondérante oppose l'acide le plus fort à la base la plus forte (règle du gamma) et produit l'acide et la base les plus faibles.
\item Sa constante vaut $K = 10^{\Delta pK_a}$ ; si $K > 10^{4}$, la réaction est quasi totale (flèche simple), sinon c'est un équilibre (double flèche + tableau d'avancement).
\end{itemize}
\end{aretenir}

\coursec{Relation pH--pKa et domaines de prédominance}

Dans la section précédente, nous avons classé les couples acide/base grâce à leur constante d'acidité $K_a$ (ou à leur $pK_a$) et appris à prévoir le sens des réactions acido-basiques. Mais une question demeure, et elle est d'une importance pratique immense : lorsqu'on place un couple acide/base en solution, \emph{quelle forme domine} ? L'acide $\ce{AH}$ ou sa base conjuguée $\ce{A-}$ ? C'est cette question qui permet de comprendre pourquoi un indicateur coloré change de teinte, pourquoi le sang reste stable autour de pH 7,4, ou pourquoi le vinaigre contient à la fois de l'acide éthanoïque et des ions éthanoate. Toute la réponse tient en une seule formule, célèbre dans le monde entier : la \cle{relation de Henderson-Hasselbalch}.

\courssub{Établissement de la relation pH = pKa + log([base]/[acide])}

\textbf{Point de départ : la constante d'acidité.} Considérons un couple acide/base quelconque $\ce{AH}/\ce{A-}$. En solution aqueuse, l'acide réagit partiellement avec l'eau selon l'équilibre :

\[ \ce{AH + H2O <=> A- + H3O+} \]

Cet équilibre est caractérisé par la constante d'acidité :

\[ K_a = \frac{[\ce{A-}]\,[\ce{H3O+}]}{[\ce{AH}]} \]

\textbf{Un peu d'algèbre.} Isolons la concentration en ions oxonium :

\[ [\ce{H3O+}] = K_a \times \frac{[\ce{AH}]}{[\ce{A-}]} \]

Appliquons maintenant la fonction $-\log$ aux deux membres de l'égalité. Rappelons deux propriétés du logarithme : $-\log(a \times b) = -\log a - \log b$ et $-\log(1/x) = \log x$. Il vient :

\[ -\log[\ce{H3O+}] = -\log K_a - \log\frac{[\ce{AH}]}{[\ce{A-}]} = -\log K_a + \log\frac{[\ce{A-}]}{[\ce{AH}]} \]

Or, par définition, $-\log[\ce{H3O+}] = pH$ et $-\log K_a = pK_a$. On obtient donc la relation fondamentale :

\begin{definition}[Relation de Henderson-Hasselbalch]
Pour tout couple acide/base $\ce{AH}/\ce{A-}$ en solution aqueuse, le pH de la solution est relié à la composition du couple par :

\[ \boxed{\; pH = pK_a + \log\frac{[\ce{A-}]}{[\ce{AH}]} = pK_a + \log\frac{[\text{base}]}{[\text{acide}]} \;} \]

Cette relation relie le \cle{pH} d'une solution à la \cle{composition} du couple acide/base présent en solution. Elle est valable quelles que soient les concentrations, dès lors que l'équilibre $\ce{AH + H2O <=> A- + H3O+}$ est établi.
\end{definition}

\textbf{Lecture intuitive.} Cette formule dit quelque chose de très simple : plus il y a de base $\ce{A-}$ par rapport à l'acide $\ce{AH}$, plus le rapport $\dfrac{[\ce{A-}]}{[\ce{AH}]}$ est grand, plus son logarithme est grand, et donc plus le pH est élevé. C'est cohérent : une solution riche en base est basique ! Inversement, si l'acide domine, le rapport est inférieur à 1, son logarithme est négatif, et le pH est inférieur au $pK_a$.

\begin{attention}[Erreur fréquente : le sens du logarithme]
L'erreur la plus classique au Baccalauréat est d'\textbf{inverser acide et base} dans le logarithme. Retiens bien :

\[ pH = pK_a + \log\frac{[\text{\textbf{base}}]}{[\text{\textbf{acide}}]} \qquad \text{et NON} \qquad pK_a + \log\frac{[\text{acide}]}{[\text{base}]} \]

\textbf{Astuce de vérification :} si la base prédomine ($[\ce{A-}] > [\ce{AH}]$), la solution doit être plus basique que le $pK_a$, donc $pH > pK_a$. Or le rapport base/acide est alors supérieur à 1, son logarithme est positif : la formule donne bien $pH > pK_a$. Si ton calcul te donne l'inverse, c'est que tu as inversé le rapport. Autre moyen mnémotechnique : « la base est \emph{en haut} (au numérateur), comme le pH qui \emph{monte} quand on ajoute de la base ».
\end{attention}

\courssub{Cas particulier fondamental : pH = pKa à la demi-équivalence}

Que se passe-t-il si les deux formes du couple sont présentes en \emph{quantités égales} ? Si $[\ce{AH}] = [\ce{A-}]$, alors le rapport vaut 1, et comme $\log 1 = 0$, la relation de Henderson-Hasselbalch se simplifie remarquablement :

\begin{propriete}[Égalité des concentrations]
Lorsque les concentrations des formes acide et basique d'un couple sont égales :

\[ [\ce{AH}] = [\ce{A-}] \quad \Longrightarrow \quad \boxed{\; pH = pK_a \;} \]

Cette situation est réalisée à la \cle{demi-équivalence} d'un dosage acido-basique : quand on a versé exactement la moitié du volume de base nécessaire pour atteindre l'équivalence, la moitié de l'acide a été transformée en base conjuguée, donc $[\ce{AH}] = [\ce{A-}]$ et le pH-mètre affiche exactement le $pK_a$ du couple.
\end{propriete}

Ce résultat est une mine d'or expérimentale : il suffit de repérer le point de demi-équivalence sur une courbe de dosage pour \emph{lire directement} le $pK_a$ sur l'axe des pH. Nous y reviendrons à la fin de cette section.

\courssub{Domaines de prédominance des formes acide et basique}

\textbf{Question concrète.} Tu verses de l'acide éthanoïque dans une solution de pH connu, disons pH = 2. Sous quelle forme se trouve-t-il majoritairement : $\ce{CH3COOH}$ ou $\ce{CH3COO-}$ ? La relation de Henderson-Hasselbalch répond instantanément.

Réécrivons-la sous la forme :

\[ \frac{[\ce{A-}]}{[\ce{AH}]} = 10^{\,pH - pK_a} \]

Trois cas se présentent :

\begin{itemize}
\item Si $pH = pK_a$ : le rapport vaut $10^0 = 1$ : les deux formes sont en concentrations égales.
\item Si $pH > pK_a$ : le rapport est supérieur à 1 : la forme \cle{basique} $\ce{A-}$ prédomine.
\item Si $pH < pK_a$ : le rapport est inférieur à 1 : la forme \cle{acide} $\ce{AH}$ prédomine.
\end{itemize}

Pour trancher nettement, on adopte une convention : une forme est dite \emph{prédominante} lorsqu'elle est au moins \textbf{10 fois plus concentrée} que l'autre, c'est-à-dire lorsque le rapport des concentrations dépasse 10 (ou est inférieur à 1/10). Comme $\log 10 = 1$, cela correspond à un écart de \textbf{1 unité de pH} de part et d'autre du $pK_a$.

\begin{aretenir}[Domaines de prédominance]
Pour un couple $\ce{AH}/\ce{A-}$ de $pK_a$ donné :

\begin{center}
\begin{tabularx}{\linewidth}{|l|X|}
\hline
\rowcolor{popblueL} \textbf{Condition sur le pH} & \textbf{Espèce prédominante} \\
\hline
$pH < pK_a - 1$ & La forme \textbf{acide} $\ce{AH}$ prédomine ($[\ce{AH}] > 10\,[\ce{A-}]$) \\
\hline
$pK_a - 1 < pH < pK_a + 1$ & Zone de \textbf{coexistence} : les deux formes sont présentes en proportions comparables \\
\hline
$pH > pK_a + 1$ & La forme \textbf{basique} $\ce{A-}$ prédomine ($[\ce{A-}] > 10\,[\ce{AH}]$) \\
\hline
\end{tabularx}
\end{center}
\end{aretenir}

\begin{methode}[Déterminer l'espèce prédominante en solution]
Face à un exercice demandant quelle espèce domine, applique systématiquement cette démarche :
\begin{enumerate}
\item \textbf{Identifier} le couple acide/base concerné et écrire son équilibre : $\ce{AH + H2O <=> A- + H3O+}$.
\item \textbf{Relever} le $pK_a$ du couple (donné dans l'énoncé ou dans une table).
\item \textbf{Comparer} le pH de la solution au $pK_a$ :
   \begin{itemize}
   \item $pH < pK_a$ : l'acide $\ce{AH}$ prédomine ;
   \item $pH > pK_a$ : la base $\ce{A-}$ prédomine.
   \end{itemize}
\item \textbf{Quantifier} si demandé : calculer le rapport $\dfrac{[\ce{A-}]}{[\ce{AH}]} = 10^{pH - pK_a}$.
\item \textbf{Conclure} en rédigeant une phrase complète avec l'espèce chimique écrite correctement.
\end{enumerate}
\end{methode}

\courssub{Diagramme de prédominance du couple acide éthanoïque / ion éthanoate}

Appliquons ce qui précède au couple $\ce{CH3COOH}/\ce{CH3COO-}$, de $pK_a = 4{,}8$ à \SI{25}{\celsius}. On représente les domaines de prédominance sur un axe horizontal gradué en pH : c'est le \cle{diagramme de prédominance}.

\begin{popfigure}
\begin{center}
\begin{tikzpicture}[x=0.85cm]
% Axe de pH
\draw[->, thick, popdark] (0,0) -- (14.8,0) node[below left=2pt and -30pt] {};
\node[below] at (14.3,-0.05) {pH};
% Graduations
\foreach \x in {0,1,...,14} {
  \draw[popdark] (\x,0.08) -- (\x,-0.08);
  \node[below, font=\small, popdark] at (\x,-0.1) {\x};
}
% Marqueurs pKa
\draw[very thick, poporange] (4.8,-0.35) -- (4.8,0.9);
\node[above, poporange, font=\small\bfseries] at (4.8,0.9) {$pK_a = 4{,}8$};
\draw[dashed, popmuted] (3.8,-0.35) -- (3.8,0.55);
\draw[dashed, popmuted] (5.8,-0.35) -- (5.8,0.55);
\node[below, popmuted, font=\footnotesize] at (3.8,-0.42) {$pK_a-1$};
\node[below, popmuted, font=\footnotesize] at (5.8,-0.42) {$pK_a+1$};
% Zone acide
\fill[popblue!25] (0,0.18) rectangle (3.8,0.62);
\node[font=\small\bfseries, popblue] at (1.9,0.4) {$\ce{CH3COOH}$ prédomine};
% Zone intermédiaire
\fill[poppurple!18] (3.8,0.18) rectangle (5.8,0.62);
\node[font=\footnotesize, poppurple] at (4.8,0.4) {coexistence};
% Zone basique
\fill[popgreen!25] (5.8,0.18) rectangle (14,0.62);
\node[font=\small\bfseries, popgreen!60!black] at (9.9,0.4) {$\ce{CH3COO-}$ prédomine};
\end{tikzpicture}
\end{center}
\legende{Diagramme de prédominance du couple $\ce{CH3COOH}/\ce{CH3COO-}$ ($pK_a = 4{,}8$). En dessous de pH 3,8, l'acide éthanoïque domine ; au-dessus de pH 5,8, c'est l'ion éthanoate qui domine ; entre les deux, les deux formes coexistent en proportions notables.}
\end{popfigure}

Ce diagramme se lit comme une carte : tu localises le pH de ta solution sur l'axe, et tu lis immédiatement quelle espèce règne. Par exemple, le jus de bissap (pH $\approx 3$) plonge l'acide éthanoïque sous sa forme $\ce{CH3COOH}$, tandis qu'une solution de pH 9 le transformerait presque totalement en ions $\ce{CH3COO-}$.

\courssub{Diagramme de distribution : les courbes en S}

Le diagramme de prédominance donne une réponse « tout ou rien ». Pour être plus fin, on trace le \cle{diagramme de distribution} : le pourcentage de chaque forme en fonction du pH. À partir de la relation de Henderson-Hasselbalch et de la conservation de la matière ($[\ce{AH}] + [\ce{A-}] = c$, constante), on démontre que :

\[ \%\,\ce{AH} = \frac{100}{1 + 10^{pH - pK_a}} \qquad \text{et} \qquad \%\,\ce{A-} = \frac{100}{1 + 10^{pK_a - pH}} \]

Ces deux fonctions ont l'allure de courbes sigmoïdes (en « S ») qui se croisent exactement à $pH = pK_a$, où chaque forme représente 50 \,\% du total.

\begin{popfigure}
\begin{center}
\begin{tikzpicture}
\begin{axis}[
  width=0.95\linewidth, height=7cm,
  xlabel={pH}, ylabel={Pourcentage (\%)},
  xmin=0, xmax=10, ymin=0, ymax=105,
  xtick={0,2,4,4.8,6,8,10},
  ytick={0,25,50,75,100},
  grid=both, grid style={popline!40},
  legend style={at={(0.5,-0.28)}, anchor=north, legend columns=2, draw=popline},
  axis line style={popdark},
  label style={popdark},
]
% %AH = 100/(1+10^(pH-4.8))
\addplot[popcurve, popblue, domain=0:10, samples=200] {100/(1+10^(x-4.8))};
\addlegendentry{$\%\,\ce{CH3COOH}$ (forme acide)}
% %A- = 100/(1+10^(4.8-pH))
\addplot[popcurve, popgreen!70!black, domain=0:10, samples=200] {100/(1+10^(4.8-x))};
\addlegendentry{$\%\,\ce{CH3COO-}$ (forme basique)}
% Croisement à pH = pKa
\draw[dashed, poporange, thick] (axis cs:4.8,0) -- (axis cs:4.8,50);
\draw[dashed, poporange, thick] (axis cs:0,50) -- (axis cs:4.8,50);
\fill[poporange] (axis cs:4.8,50) circle (2.2pt);
\node[poporange, font=\small\bfseries, anchor=south west] at (axis cs:4.9,52) {$pH = pK_a = 4{,}8$};
\node[poporange, font=\small, anchor=north] at (axis cs:4.8,-3) {};
\end{axis}
\end{tikzpicture}
\end{center}
\legende{Diagramme de distribution du couple $\ce{CH3COOH}/\ce{CH3COO-}$. Les deux courbes sigmoïdes se croisent à $pH = pK_a = 4{,}8$, où les deux formes représentent chacune 50 \,\% de l'espèce dissoute. À pH 2, plus de 99 \,\% de l'acide est sous forme $\ce{CH3COOH}$ ; à pH 8, plus de 99 \,\% est sous forme $\ce{CH3COO-}$.}
\end{popfigure}

Retiens les points de repère utiles pour les calculs rapides : à $pH = pK_a - 1$, la forme acide représente environ 91 \,\% ; à $pH = pK_a - 2$, environ 99 \,\%. Symétriquement, à $pH = pK_a + 1$, la forme basique atteint 91 \,\%, et 99 \,\% à $pK_a + 2$.

\courssub{Exemple chiffré : l'acide éthanoïque dans le sang}

\begin{exemplebox}[Quelle forme du couple dans le sang ?]
Le sang humain est remarquablement régulé à $pH = 7{,}4$. Si de l'acide éthanoïque (ou de l'éthanoate) y est introduit, sous quelle forme se trouve-t-il ?

\textbf{Donnée :} $pK_a(\ce{CH3COOH}/\ce{CH3COO-}) = 4{,}8$.

\textbf{Résolution.} Appliquons la relation de Henderson-Hasselbalch :

\[ \frac{[\ce{CH3COO-}]}{[\ce{CH3COOH}]} = 10^{pH - pK_a} = 10^{7{,}4 - 4{,}8} = 10^{2{,}6} \approx 400 \]

\textbf{Conclusion :} il y a \textbf{400 fois plus} d'ions éthanoate que de molécules d'acide éthanoïque dans le sang. La forme basique $\ce{CH3COO-}$ domine \emph{largement} : c'est cohérent avec le diagramme, puisque $pH = 7{,}4 > pK_a + 1 = 5{,}8$. Le pourcentage de forme basique vaut $\dfrac{400}{401} \times 100 \approx 99{,}75\,\%$.
\end{exemplebox}

\begin{exoresolu}[L'acide benzoïque, conservateur alimentaire]
L'acide benzoïque $\ce{C6H5COOH}$ (conservateur E210, utilisé dans certaines boissons gazeuses vendues à Douala et Yaoundé) appartient au couple $\ce{C6H5COOH}/\ce{C6H5COO-}$ de $pK_a = 4{,}2$. Une boisson a un pH de 3,0.
\begin{enumerate}
\item Quelle forme du couple prédomine dans cette boisson ?
\item Calculer le rapport $\dfrac{[\ce{C6H5COO-}]}{[\ce{C6H5COOH}]}$ et le pourcentage de forme acide.
\end{enumerate}
\tcblower
\textbf{1. Forme prédominante.} On compare le pH au $pK_a$ : $pH = 3{,}0 < pK_a = 4{,}2$. On est même dans le cas $pH < pK_a - 1 = 3{,}2$. Donc la forme \textbf{acide} $\ce{C6H5COOH}$ prédomine nettement.

\textbf{2. Rapport des concentrations.} D'après la relation de Henderson-Hasselbalch :

\[ \frac{[\ce{C6H5COO-}]}{[\ce{C6H5COOH}]} = 10^{pH - pK_a} = 10^{3{,}0 - 4{,}2} = 10^{-1{,}2} \approx 0{,}063 \]

Il y a donc environ $1/0{,}063 \approx 16$ fois plus d'acide que de base. Le pourcentage de forme acide est :

\[ \%\,\ce{C6H5COOH} = \frac{[\ce{C6H5COOH}]}{[\ce{C6H5COOH}] + [\ce{C6H5COO-}]} \times 100 = \frac{1}{1 + 0{,}063} \times 100 \approx 94\,\% \]

\textbf{Conclusion :} dans la boisson, environ 94 \,\% du conservateur est sous forme acide $\ce{C6H5COOH}$. C'est d'ailleurs cette forme acide, non ionisée, qui traverse les membranes des micro-organismes et exerce l'action conservatrice : le choix d'un pH acide pour les boissons n'est pas un hasard !
\end{exoresolu}

\courssub{Application aux indicateurs colorés}

\textbf{L'idée géniale.} Un \cle{indicateur coloré} acido-basique est lui-même un couple acide/base, noté $\ce{HInd}/\ce{Ind-}$, dont la particularité est que \textbf{la forme acide et la forme basique ont des couleurs différentes}. L'équilibre en jeu est :

\[ \ce{HInd + H2O <=> Ind- + H3O+} \]

La relation de Henderson-Hasselbalch s'applique : $pH = pK_{a,\text{ind}} + \log\dfrac{[\ce{Ind-}]}{[\ce{HInd}]}$. La couleur observée dépend donc du pH :

\begin{itemize}
\item Si $pH < pK_{a,\text{ind}} - 1$ : la forme acide $\ce{HInd}$ prédomine $\Rightarrow$ on voit la \textbf{teinte acide}.
\item Si $pH > pK_{a,\text{ind}} + 1$ : la forme basique $\ce{Ind-}$ prédomine $\Rightarrow$ on voit la \textbf{teinte basique}.
\item Entre les deux : mélange des deux formes $\Rightarrow$ \textbf{teinte intermédiaire} (zone de virage).
\end{itemize}

\begin{definition}[Zone de virage]
La \cle{zone de virage} d'un indicateur coloré est l'intervalle de pH, centré sur son $pK_a$, dans lequel l'indicateur passe progressivement de sa teinte acide à sa teinte basique. On la prend usuellement égale à $[pK_a - 1 \,;\, pK_a + 1]$.
\end{definition}

\begin{center}
\begin{tabularx}{\linewidth}{|l|X|X|X|l|}
\hline
\rowcolor{popblueL} \textbf{Indicateur} & \textbf{Teinte acide} & \textbf{Zone de virage} & \textbf{Teinte basique} & $pK_a$ \\
\hline
Hélianthine & rouge & 3,1 -- 4,4 & jaune orangé & $\approx 3{,}7$ \\
\hline
Bleu de bromothymol (BBT) & jaune & 6,0 -- 7,6 & bleu & $\approx 7{,}0$ \\
\hline
Phénolphtaléine & incolore & 8,2 -- 10,0 & rose fuchsia & $\approx 9{,}0$ \\
\hline
\end{tabularx}
\end{center}

\begin{popfigure}
\begin{center}
\begin{tikzpicture}[x=0.8cm]
% Axe pH
\draw[->, thick, popdark] (0,0) -- (14.8,0);
\node[below, popdark] at (14.4,-0.05) {pH};
\foreach \x in {0,1,...,14} {
  \draw[popdark] (\x,0.06) -- (\x,-0.06);
  \node[below, font=\footnotesize, popdark] at (\x,-0.08) {\x};
}
% Hélianthine : virage 3.1 - 4.4
\fill[red!70!orange] (0,0.55) rectangle (3.1,1.0);
\shade[left color=red!70!orange, right color=yellow!80!orange] (3.1,0.55) rectangle (4.4,1.0);
\fill[yellow!80!orange] (4.4,0.55) rectangle (14,1.0);
\node[font=\footnotesize\bfseries, popdark, anchor=west] at (14.1,0.78) {};
\node[font=\footnotesize, popdark] at (7,1.22) {\textbf{Hélianthine} ($pK_a \approx 3{,}7$)};
% BBT : virage 6.0 - 7.6
\fill[yellow!85!orange] (0,1.75) rectangle (6.0,2.2);
\shade[left color=yellow!85!orange, right color=popblue!80] (6.0,1.75) rectangle (7.6,2.2);
\fill[popblue!80] (7.6,1.75) rectangle (14,2.2);
\node[font=\footnotesize, popdark] at (7,2.42) {\textbf{Bleu de bromothymol} ($pK_a \approx 7{,}0$)};
% Phénolphtaléine : virage 8.2 - 10.0
\fill[popcream] (0,2.95) rectangle (8.2,3.4);
\shade[left color=popcream, right color=poppink!85] (8.2,2.95) rectangle (10.0,3.4);
\fill[poppink!85] (10.0,2.95) rectangle (14,3.4);
\draw[popline] (0,2.95) rectangle (14,3.4);
\node[font=\footnotesize, popdark] at (7,3.62) {\textbf{Phénolphtaléine} ($pK_a \approx 9{,}0$)};
\end{tikzpicture}
\end{center}
\legende{Zones de virage de trois indicateurs colorés usuels sur l'échelle de pH. Chaque indicateur change de couleur autour de son propre $pK_a$ : on choisit l'indicateur dont la zone de virage contient le pH à l'équivalence du dosage étudié.}
\end{popfigure}

\begin{attention}[Bien choisir son indicateur]
Un indicateur coloré n'est utile pour un dosage que si sa \textbf{zone de virage encadre le pH à l'équivalence}. Pour doser un acide faible par une base forte (équivalence vers pH 8--9), on choisit la phénolphtaléine, jamais l'hélianthine (qui virerait bien avant l'équivalence, vers pH 4). Autre piège : un indicateur est lui-même un acide ! On n'en verse que \textbf{quelques gouttes}, sinon il perturberait le dosage en consommant une partie du réactif titrant.
\end{attention}

\begin{savaistu}[Le BBT, détecteur du CO$_2$ que tu expires]
Le bleu de bromothymol vire du jaune au bleu autour de pH 7. Souffle avec une paille dans de l'eau pure additionnée de BBT : le dioxyde de carbone de ton souffle se dissout et forme de l'acide carbonique selon $\ce{CO2 + H2O <=> HCO3- + H3O+}$ ; le pH chute sous 6 et la solution devient \textbf{jaune} en quelques secondes. C'est le même principe que le fameux test à l'eau de chaux, qui se trouble en présence de $\ce{CO2}$ par formation de carbonate de calcium $\ce{CaCO3}$. Les enseignants camerounais utilisent souvent ces deux expériences, simples et spectaculaires, pour prouver que l'air expiré contient du dioxyde de carbone. Mieux : ce même équilibre $\ce{CO2}/\ce{HCO3-}$ ($pK_a \approx 6{,}4$) est le principal \emph{tampon} qui maintient le pH de ton sang à 7,4 !
\end{savaistu}

\courssub{Détermination expérimentale d'un pKa : la méthode de la demi-équivalence}

Nous avons vu que $pH = pK_a$ lorsque $[\ce{AH}] = [\ce{A-}]$. Cette égalité est réalisée à la \textbf{demi-équivalence} d'un dosage : c'est la base d'une méthode expérimentale simple et puissante de mesure d'un $pK_a$.

\begin{methode}[Mesurer un pKa par dosage pH-métrique]
\begin{enumerate}
\item \textbf{Doser} un volume $V_A$ d'acide faible (par exemple l'acide éthanoïque d'un vinaigre dilué) par une solution d'hydroxyde de sodium de concentration connue, en suivant le pH au pH-mètre : $\ce{CH3COOH + HO- -> CH3COO- + H2O}$ (réaction quasi totale).
\item \textbf{Tracer} la courbe $pH = f(V_B)$ et déterminer le volume équivalent $V_E$ (saut de pH, méthode des tangentes ou de la dérivée).
\item \textbf{Repérer} le point de demi-équivalence : $V = \dfrac{V_E}{2}$. À ce volume, la moitié de l'acide a été consommée et transformée en éthanoate : $[\ce{CH3COOH}] = [\ce{CH3COO-}]$.
\item \textbf{Lire} le pH en ce point : c'est directement le $pK_a$ du couple.
\item \textbf{Calculer} $K_a = 10^{-pK_a}$ et comparer à la valeur tabulée ($pK_a = 4{,}8$ pour l'acide éthanoïque à \SI{25}{\celsius}).
\end{enumerate}
\end{methode}

C'est exactement cette démarche que tu mettras en œuvre en TP lors du dosage du vinaigre : en mesurant le pH à la demi-équivalence, tu détermineras expérimentalement la constante d'acidité du couple de l'acide éthanoïque, avant d'exploiter l'équivalence pour remonter au degré d'acidité du vinaigre. La relation de Henderson-Hasselbalch, établie ici sur le papier, deviendra alors un véritable \emph{instrument de mesure} entre tes mains.

\begin{aretenir}[L'essentiel de la section]
\begin{itemize}
\item Relation de Henderson-Hasselbalch : $pH = pK_a + \log\dfrac{[\text{base}]}{[\text{acide}]}$.
\item Si $[\ce{AH}] = [\ce{A-}]$, alors $pH = pK_a$ : c'est la situation de demi-équivalence d'un dosage.
\item $pH < pK_a - 1$ : l'acide $\ce{AH}$ prédomine ; $pH > pK_a + 1$ : la base $\ce{A-}$ prédomine ; entre les deux, coexistence.
\item Les courbes de distribution se croisent à 50 \,\% pour $pH = pK_a$.
\item Un indicateur coloré est un couple $\ce{HInd}/\ce{Ind-}$ dont les deux formes ont des couleurs différentes ; sa zone de virage est centrée sur son $pK_a$.
\item Le $pK_a$ d'un couple se mesure expérimentalement en lisant le pH à la demi-équivalence d'un dosage pH-métrique.
\end{itemize}
\end{aretenir}

\coursec{Applications : détermination expérimentale de $K_\mathrm{a}$ et degré d'acidité du vinaigre}

Dans la section précédente, nous avons établi la relation fondamentale $\mathrm{pH} = \mathrm{p}K_\mathrm{a} + \log\dfrac{[\text{base}]}{[\text{acide}]}$ et identifié les domaines de prédominance des formes acide et basique d'un couple. Cette relation n'est pas seulement un outil de prévision : elle devient une \textbf{méthode expérimentale} pour déterminer la constante d'acidité d'un couple inconnu. De plus, la réaction totale entre un acide faible et une base forte permet de doser l'acide éthanoïque présent dans le vinaigre et d'en exprimer la richesse sous forme de \emph{degré d'acidité}, une grandeur très utilisée dans l'industrie agroalimentaire camerounaise. Nous allons voir comment passer de la théorie au laboratoire, puis du laboratoire à l'étiquette du produit.

\courssub{Détermination de $K_\mathrm{a}$ et du coefficient d'ionisation par mesure de pH}

\paragraph{Intuition.}
Si l'on prépare une solution d'un acide faible AH de concentration connue $C$, l'équilibre $\ce{AH + H2O <=> A- + H3O+}$ s'établit. Si l'on mesure le pH de cette solution, on accède directement à la concentration en ions oxydronium $[\ce{H3O+}]$ à l'équilibre. Connaissant $C$ et $[\ce{H3O+}]_\text{eq}$, on peut remonter à $K_\mathrm{a}$ grâce au tableau d'avancement.

\begin{methode}[Détermination de $K_\mathrm{a}$ par pH-métrie directe]
\begin{enumerate}
\item Préparer une solution de l'acide faible AH de concentration molaire $C$ connue (par exemple \SI{0,10}{mol.L^{-1}}).
\item Mesurer le pH de cette solution à température constante (généralement \SI{25}{\celsius}).
\item Calculer $[\ce{H3O+}]_\text{eq} = 10^{-\mathrm{pH}}$.
\item Établir le tableau d'avancement de l'ionisation de AH.
\item En déduire $K_\mathrm{a} = \dfrac{[\ce{H3O+}]_\text{eq}^2}{C - [\ce{H3O+}]_\text{eq}}$.
\item Calculer le coefficient d'ionisation $\alpha = \dfrac{[\ce{H3O+}]_\text{eq}}{C}$.
\end{enumerate}
\end{methode}

\begin{exemplebox}[Acide éthanoïque \SI{0,10}{mol.L^{-1}}]
On mesure le pH d'une solution d'acide éthanoïque \SI{0,10}{mol.L^{-1}} à \SI{25}{\celsius} : $\mathrm{pH} = 2,88$.
\begin{itemize}
\item $[\ce{H3O+}]_\text{eq} = 10^{-2,88} = \SI{1,32e-3}{mol.L^{-1}}$.
\item Tableau d'avancement :
\begin{center}
\begin{tabular}{|c|c|c|c|}
\hline
 & $\ce{CH3COOH}$ & $\ce{H2O}$ & $\ce{CH3COO-}$ + $\ce{H3O+}$ \\
\hline
État initial & $C = 0,10$ & excès & 0 \\
\hline
État final & $C - x$ & excès & $x$ \\
\hline
\end{tabular}
\end{center}
avec $x = [\ce{H3O+}]_\text{eq} = \SI{1,32e-3}{mol.L^{-1}}$.
\item $K_\mathrm{a} = \dfrac{x^2}{C - x} = \dfrac{(\SI{1,32e-3}{})^2}{0,10 - \SI{1,32e-3}{}} = \SI{1,76e-5}{}$.
\item $\mathrm{p}K_\mathrm{a} = -\log K_\mathrm{a} = 4,75$.
\item $\alpha = \dfrac{x}{C} = \dfrac{\SI{1,32e-3}{}}{0,10} = \SI{1,32}{\%}$.
\end{itemize}
L'acide éthanoïque est bien un acide faible : seulement \SI{1,3}{\%} des molécules sont ionisées.
\end{exemplebox}

\begin{attention}[Piège : approximation $C - x \approx C$]
On a souvent tendance à écrire $K_\mathrm{a} \approx \dfrac{[\ce{H3O+}]^2}{C}$ en négligeant $x$ devant $C$. Cette approximation n'est valable que si $\alpha < 5\%$ (soit $x < 0,05 C$). Ici $\alpha = 1,32\%$, l'approximation est correcte ($K_\mathrm{a} \approx \SI{1,74e-5}{}$). Mais pour un acide moins faible ou plus concentré, il faut garder le terme exact $C - x$. \textbf{Vérifiez toujours le critère $\alpha < 5\%$ avant d'approximer.}
\end{attention}

\courssub{Détermination du $\mathrm{p}K_\mathrm{a}$ par dosage pH-métrique : la méthode de la demi-équivalence}

\paragraph{Intuition.}
Lors du dosage d'un acide faible par une base forte, on suit l'évolution du pH en fonction du volume de base ajouté. Au point de \textbf{demi-équivalence}, la moitié de l'acide initial a été transformée en base conjuguée : $n_\text{AH} = n_\text{A-}$. D'après la relation $\mathrm{pH} = \mathrm{p}K_\mathrm{a} + \log\dfrac{[\ce{A-}]}{[\ce{AH}]}$, le rapport des concentrations vaut 1, son logarithme vaut 0, donc \textbf{pH = pKa}. C'est une méthode très précise car elle ne nécessite pas de connaître la concentration exacte de l'acide, seulement le volume à l'équivalence.

\begin{definition}[Point de demi-équivalence]
Lors du dosage d'un acide faible par une base forte (ou d'une base faible par un acide fort), le point de demi-équivalence correspond au volume $V_{\text{E}}/2$ pour lequel la quantité de matière d'acide restant est égale à la quantité de matière de base conjuguée formée. À ce point, $\mathrm{pH} = \mathrm{p}K_\mathrm{a}$ du couple $\ce{AH}/\ce{A-}$.
\end{definition}

\begin{experience}[Dosage pH-métrique d'une solution d'acide éthanoïque \SI{0,10}{mol.L^{-1}} par une solution de soude \SI{0,10}{mol.L^{-1}}]
\textbf{Dispositif :} Voir figure ci-dessous (burette graduée, bécher, pH-mètre étalonné, agitateur magnétique).
\textbf{Protocole :}
\begin{enumerate}
\item Verser \SI{20,0}{mL} de solution d'acide éthanoïque \SI{0,10}{mol.L^{-1}} (préparée par dilution du vinaigre commercial si nécessaire) dans un bécher, ajouter \SI{50}{mL} d'eau distillée et une barre aimantée.
\item Plonger l'électrode pH-métrique (préalablement étalonnée avec deux tampons pH 4,01 et 7,00 ou 9,21).
\item Titrer avec une solution de soude de concentration connue $C_\text{B} = \SI{0,100}{mol.L^{-1}}$ en ajoutant par fractions (ex. \SI{1,0}{mL} puis \SI{0,5}{mL} près de l'équivalence) et en notant le pH après chaque addition.
\item Tracer la courbe $\mathrm{pH} = f(V_\text{B})$.
\end{enumerate}
\textbf{Observations :} La courbe présente un saut de pH à l'équivalence $V_\text{E} = \SI{20,0}{mL}$ (car $C_\text{AH}V_\text{AH} = C_\text{B}V_\text{E}$). Le pH à l'équivalence est basique ($\approx 8,7$) car l'éthanoate formé est une base faible. Le point de demi-équivalence se repère à l'abscisse $V_\text{E}/2 = \SI{10,0}{mL}$.
\textbf{Interprétation :} La réaction de dosage est $\ce{CH3COOH + HO- -> CH3COO- + H2O}$. Elle est totale ($K$ très grand) car la base forte $\ce{HO-}$ réagit complètement avec l'acide faible. À la demi-équivalence, $n_{\ce{CH3COOH}} = n_{\ce{CH3COO-}}$, donc $\mathrm{pH} = \mathrm{p}K_\mathrm{a} \approx 4,75$.
\end{experience}

\begin{popfigure}
\begin{tikzpicture}
\begin{axis}[
    width=\linewidth,
    height=0.55\linewidth,
    xlabel={Volume de soude $V_\text{B}$ (mL)},
    ylabel={pH},
    xmin=0, xmax=22,
    ymin=2, ymax=12,
    grid=major,
    grid style={popline},
    axis lines=middle,
    tick label style={font=\small},
    label style={font=\small},
    clip=false,
    domain=0:20,
    samples=400
]
% Courbe sigmoïde réaliste pour 0.10 M CH3COOH / 0.10 M NaOH
% pH initial = 2.88, pH équiv = 8.70, V_E = 20 mL, V_1/2 = 10 mL
\addplot[popcurve, thick] 
({x}, 
{2.88 + 5.82/(1+exp(-2.8*(x-10)))} 
);

% Equivalence point
\draw[poporange, dashed, thick] (axis cs:20,2) -- (axis cs:20,8.7) node[pos=0.5, right, font=\small, poporange] {$V_\text{E} = \SI{20,0}{mL}$};
\draw[poporange, thick, fill=poporange] (axis cs:20,8.7) circle (2.5pt);

% Half-equivalence point
\draw[popgreen, dashed, thick] (axis cs:10,2) -- (axis cs:10,5.75) node[pos=0.5, left, font=\small, popgreen] {$V_\text{E}/2 = \SI{10,0}{mL}$};
\draw[popgreen, thick, fill=popgreen] (axis cs:10,5.75) circle (2.5pt);
\node[popgreen, anchor=south west, font=\small] at (axis cs:10,5.75) {$\mathrm{pH} = \mathrm{p}K_\mathrm{a} = 4,75$};

% Initial pH
\draw[popblue, dashed] (axis cs:0,2.88) -- (axis cs:2,2.88);
\node[popblue, anchor=east, font=\small] at (axis cs:0,2.88) {pH initial $\approx 2,9$};

% Equivalence pH
\draw[poppurple, dashed] (axis cs:20,8.7) -- (axis cs:18,8.7);
\node[poppurple, anchor=east, font=\small] at (axis cs:20,8.7) {pH équiv. $\approx 8,7$};

\end{axis}
\end{tikzpicture}
\legende{Courbe de dosage pH-métrique d'un acide faible (acide éthanoïque \SI{0,10}{mol.L^{-1}}, \SI{20,0}{mL}) par une base forte (soude \SI{0,10}{mol.L^{-1}}). Repérage de l'équivalence $V_\text{E}$ et de la demi-équivalence où $\mathrm{pH} = \mathrm{p}K_\mathrm{a}$.}
\end{popfigure}

\begin{popfigure}
\begin{tikzpicture}[scale=0.9, every node/.style={font=\small}]
% Support
\draw[thick, popdark] (-3,-2) -- (3,-2);
\draw[thick, popdark] (0,-2) -- (0,4);
% Burette
\draw[thick, popblue] (2.5, -1) -- (2.5, 3.5);
\draw[thick, popblue] (2.2, -1) -- (2.8, -1);
\draw[thick, popblue] (2.2, 3.5) -- (2.8, 3.5);
\draw[popblue, fill=popblueL] (2.5, 3.0) rectangle (2.8, 3.2); % robinet
\node[popblue, right] at (2.8, 3.1) {Robinet};
\draw[popblue, fill=popblue!20] (2.2, 1.5) rectangle (2.8, 2.8); % solution
\node[popblue, right] at (2.8, 2.15) {Soude \SI{0,10}{mol.L^{-1}}};
% Graduations burette
\foreach \y in {0, 0.5, ..., 3} {
    \draw[popblue, thin] (2.2, -1+\y) -- (2.0, -1+\y);
}
% Bécher
\draw[thick, popdark] (-1.5, -1.5) -- (-1.5, 1.5) -- (1.5, 1.5) -- (1.5, -1.5);
\draw[thick, popdark] (-1.7, -1.5) -- (1.7, -1.5); % base
% Solution dans bécher
\draw[poporange, fill=poporangeL] (-1.4, -1.4) rectangle (1.4, 0.5);
\node[poporange] at (0, -0.5) {Acide éthanoïque \SI{0,10}{M} + eau};
% Barre aimantée
\draw[popgold, thick, fill=popgoldL] (-0.3, -1.2) -- (0.3, -1.2) -- (0.3, -1.0) -- (-0.3, -1.0) -- cycle;
\node[popgold, below] at (0, -1.3) {Barre aimantée};
% Plaque agitateur
\draw[popmuted, thick, fill=popmuted!20] (-2, -1.8) rectangle (2, -1.6);
\node[popmuted] at (0, -1.7) {Agitateur magnétique};
% pH-mètre
\draw[thick, poppurple] (-0.2, 0.5) -- (-0.2, 2.5) -- (0.2, 2.5) -- (0.2, 0.5);
\draw[poppurple, fill=poppurpleL] (-0.5, 2.5) rectangle (0.5, 2.8);
\node[poppurple] at (0, 2.65) {Électrode pH};
\node[poppurple, left] at (-0.5, 1.5) {pH-mètre};
% Flèches agitation
\draw[->, popgreen, thick] (-1.0, 0.8) arc (180:0:0.5);
\draw[->, popgreen, thick] (1.0, 0.8) arc (0:-180:0.5);
\end{tikzpicture}
\legende{Montage de dosage pH-métrique : burette (soude), bécher (acide éthanoïque \SI{0,10}{M} + eau), barre aimantée, agitateur magnétique, électrode pH-métrique connectée au pH-mètre.}
\end{popfigure}

\courssub{Degré d'acidité du vinaigre : définition et détermination}

\paragraph{Contexte camerounais.}
Au Cameroun, le vinaigre (souvent issu de la fermentation acétique du vin de palme ou du bili-bili) est un conservateur traditionnel essentiel pour les piments (condiments « piment vinaigre »), les légumes (carottes, choux) et les poissons fumés. Son efficacité antibactérienne repose sur son acidité (pH < 4). La réglementation (norme NC 123) impose un \textbf{degré d'acidité} compris entre \textbf{6° et 8°} pour le vinaigre de table.

\begin{definition}[Degré d'acidité d'un vinaigre]
Le \textbf{degré d'acidité} (noté en degrés, symbole °) d'un vinaigre est la \textbf{masse d'acide éthanoïque (en grammes) contenue dans \SI{100}{g} de vinaigre}.
\[
\text{Degré} = \frac{m_{\ce{CH3COOH}}\text{ (en g)}}{m_{\text{vinaigre}}\text{ (en g)}} \times 100
\]
C'est un pourcentage massique. Un vinaigre à 6° contient \SI{6}{g} d'acide éthanoïque pur pour \SI{100}{g} de vinaigre.
\end{definition}

\paragraph{Stratégie de calcul.}
On dose l'acide éthanoïque par la soude. La réaction est totale :
\[
\ce{CH3COOH_{(aq)} + HO^{-}_{(aq)} -> CH3COO^{-}_{(aq)} + H2O_{(l)}}
\]
À l'équivalence : $n_{\ce{CH3COOH}} = n_{\ce{HO-}}$ soit $C_\text{AH} \times V_\text{AH} = C_\text{B} \times V_\text{E}$.
On en déduit la concentration molaire $C_\text{AH}$ de la solution \textbf{diluée}, puis on remonte à la concentration de la solution \textbf{mère} (vinaigre pur) en tenant compte du facteur de dilution, puis on convertit en degré massique.

\begin{methode}[Détermination du degré d'acidité d'un vinaigre commercial]
\begin{enumerate}
\item \textbf{Dilution connue :} Prélever $V_\text{prélevé}$ de vinaigre mère (pipette jaugée) et compléter à $V_\text{fiole}$ (fiole jaugée). Facteur de dilution $f = \dfrac{V_\text{fiole}}{V_\text{prélevé}}$.
\item \textbf{Dosage :} Prélèvement $V_\text{AH}$ de la solution diluée, dosage par soude $C_\text{B}$, volume à l'équivalence $V_\text{E}$.
\item \textbf{Concentration solution diluée :} $C_\text{diluée} = \dfrac{C_\text{B} \times V_\text{E}}{V_\text{AH}}$.
\item \textbf{Concentration vinaigre mère :} $C_\text{mère} = f \times C_\text{diluée}$.
\item \textbf{Masse volumique :} On utilise la masse volumique du vinaigre $\rho \approx \SI{1,01}{g.mL^{-1}}$ (proche de l'eau).
\item \textbf{Masse d'acide dans 1 L :} $m_{\ce{CH3COOH}} = C_\text{mère} \times M_{\ce{CH3COOH}}$ avec $M_{\ce{CH3COOH}} = \SI{60,05}{g.mol^{-1}}$.
\item \textbf{Masse de 1 L de vinaigre :} $m_{\text{vinaigre}} = \rho \times 1000$ (en g).
\item \textbf{Degré :} $\text{Degré} = \dfrac{m_{\ce{CH3COOH}}}{m_{\text{vinaigre}}} \times 100$.
\end{enumerate}
\end{methode}

\begin{exoresolu}[Vérification de l'étiquette d'un vinaigre commercial]
\textbf{Énoncé.} On souhaite vérifier le degré d'acidité annoncé sur une bouteille de vinaigre « Vinaigre de palme - 7° ». 
\begin{itemize}
\item Dilution : \SI{10,0}{mL} de vinaigre mère $\rightarrow$ fiole jaugée \SI{100,0}{mL} (facteur $f=10$).
\item Dosage : \SI{20,0}{mL} de solution diluée titrés par soude \SI{0,100}{mol.L^{-1}}. Volume à l'équivalence $V_\text{E} = \SI{19,8}{mL}$.
\item Masse volumique du vinaigre : $\rho = \SI{1,01}{g.mL^{-1}}$.
\item $M_{\ce{CH3COOH}} = \SI{60,05}{g.mol^{-1}}$.
\end{itemize}
Déterminer le degré d'acidité et conclure sur la conformité.

\tcblower
\textbf{Solution détaillée.}

\underline{1. Quantité de matière de soude versée à l'équivalence}
$n_{\ce{HO-}} = C_\text{B} \times V_\text{E} = \SI{0,100}{mol.L^{-1}} \times \SI{19,8e-3}{L} = \SI{1,98e-3}{mol}$.

\underline{2. Quantité de matière d'acide éthanoïque dans l'aliquot dosé}
Réaction totale $\ce{CH3COOH + HO- -> CH3COO- + H2O}$ : $n_{\ce{CH3COOH}}(\text{dilué}) = n_{\ce{HO-}} = \SI{1,98e-3}{mol}$.

\underline{3. Concentration de la solution diluée}
$C_\text{diluée} = \dfrac{n_{\ce{CH3COOH}}(\text{dilué})}{V_\text{AH}} = \dfrac{\SI{1,98e-3}{mol}}{\SI{20,0e-3}{L}} = \SI{0,0990}{mol.L^{-1}}$.

\underline{4. Concentration du vinaigre mère (facteur de dilution $f=10$)}
$C_\text{mère} = f \times C_\text{diluée} = 10 \times \SI{0,0990}{mol.L^{-1}} = \SI{0,990}{mol.L^{-1}}$.

\underline{5. Masse d'acide éthanoïque dans 1 L de vinaigre mère}
$m_{\ce{CH3COOH}} = C_\text{mère} \times M_{\ce{CH3COOH}} = \SI{0,990}{mol.L^{-1}} \times \SI{60,05}{g.mol^{-1}} = \SI{59,45}{g.L^{-1}}$.

\underline{6. Masse de 1 L de vinaigre mère}
$m_{\text{vinaigre}} = \rho \times 1000 = \SI{1,01}{g.mL^{-1}} \times 1000 = \SI{1010}{g}$.

\underline{7. Degré d'acidité}
$\text{Degré} = \dfrac{\SI{59,45}{g}}{\SI{1010}{g}} \times 100 = \textbf{5,89°} \approx \textbf{5,9°}$.

\underline{Conclusion.} Le degré trouvé (\SI{5,9}{°}) est \textbf{légèrement inférieur} à la valeur annoncée (7°) et se situe \textbf{en dessous de la norme minimale (6°)}. Ce vinaigre n'est pas conforme. Il est possible qu'il ait été trop dilué lors de la fabrication ou qu'il ait subi une évaporation d'acide (l'acide éthanoïque est volatil) lors du stockage.
\end{exoresolu}

\begin{attention}[Erreur fatale : oublier le facteur de dilution !]
C'est l'erreur \textbf{numéro 1} au Bac. Si l'on oublie de multiplier par $f=10$ à l'étape 4, on trouve $C_\text{mère} = \SI{0,099}{mol.L^{-1}}$, puis un degré de \SI{0,59}{°}... ce qui n'a aucun sens pour du vinaigre. \textbf{Relisez toujours l'énoncé pour identifier la dilution préalable (souvent 1/10 ou 1/20) et appliquez le facteur à la concentration molaire, pas au volume final.}
\end{attention}

\begin{savaistu}[Le vinaigre au Cameroun : tradition et science]
Au Cameroun, le vinaigre artisanal s'obtient par fermentation acétique du vin de palme (vin de raphia ou de palmier à huile) ou du bili-bili (bière de mil/maïs). Les bactéries acétiques (\emph{Acetobacter}) transforment l'éthanol en acide éthanoïque en présence d'oxygène : $\ce{CH3CH2OH + O2 -> CH3COOH + H2O}$. 
Le degré d'acidité traditionnel se mesure parfois « à l'ancienne » : on ajoute de la soude caustique (lessive) goutte à goutte jusqu'à ce que le goût acide disparaisse, ou on observe la réaction avec du bicarbonate (effervescence). 
La norme \textbf{NC 123} (vinaigre de table) impose : degré entre 6° et 8°, pH < 4,0, absence de métaux lourds. Un vinaigre trop faible (< 5°) ne conserve pas bien les aliments (risque de botulisme, moisissures) ; un vinaigre trop fort (> 10°) est corrosif pour l'œsophage et l'émail dentaire. 
Dans l'industrie (SOCAPALM, brasseries), le contrôle qualité utilise exactement la méthode pH-métrique ou par titrage colorimétrique (indicateur phénolphtaléine) que nous venons d'étudier.
\end{savaistu}

\courssub{Bilan de la leçon : carte mentale unificatrice}

Pour clore ce module sur les acides et bases, voici une synthèse visuelle qui relie tous les concepts abordés : force, pH, constante d'acidité, avancement, et domaines de prédominance. Cette carte mentale est votre « fiche de révision » pour le Bac.

\begin{popfigure}
\begin{tikzpicture}[
    scale=0.85,
    every node/.style={font=\small},
    concept/.style={circle, draw=popblue, thick, fill=popblueL, minimum size=2.2cm, text width=2cm, align=center},
    subconcept/.style={rectangle, rounded corners, draw=poporange, thick, fill=poporangeL, minimum width=2.5cm, minimum height=1cm, text width=2.3cm, align=center},
    link/.style={-Stealth, thick, popdark},
    bidilink/.style={Stealth-Stealth, thick, poppurple}
]
% Nœud central
\node[concept] (centre) at (0,0) {\textbf{COUPLE ACIDE/BASE}\\\ce{AH}/\ce{A-}};

% Branche 1 : Force
\node[subconcept] (force) at (-5, 3) {\textbf{FORCE}\\\textit{Acide fort / Base forte}};
\node[subconcept] (faible) at (-5, 1) {\textbf{FAIBLESSE}\\\textit{Acide faible / Base faible}};
\node[subconcept] (ionisation) at (-5, -1) {\textbf{IONISATION}\\\textit{Partielle ($\alpha < 1$)}};
\node[subconcept] (totale) at (-5, -3) {\textbf{IONISATION}\\\textit{Totale ($\alpha = 1$)}};

\draw[link] (centre) -- (force);
\draw[link] (force) -- (totale);
\draw[link] (centre) -- (faible);
\draw[link] (faible) -- (ionisation);

% Branche 2 : Constantes
\node[subconcept] (Ka) at (5, 3) {$K_\mathrm{a}$ \textit{(acidité)}};
\node[subconcept] (Kb) at (5, 1) {$K_\mathrm{b}$ \textit{(basicité)}};
\node[subconcept] (pKa) at (5, -1) {$\mathrm{p}K_\mathrm{a} = -\log K_\mathrm{a}$};
\node[subconcept] (relation) at (5, -3) {$K_\mathrm{a} \times K_\mathrm{b} = K_\text{e}$};

\draw[link] (centre) -- (Ka);
\draw[link] (Ka) -- (pKa);
\draw[link] (centre) -- (Kb);
\draw[link] (Kb) -- (relation);
\draw[link] (pKa) -- (relation);

% Branche 3 : pH et Avancement
\node[subconcept] (pH) at (0, 4) {\textbf{pH = $-\log[\ce{H3O+}]$}};
\node[subconcept] (avancement) at (0, -4) {\textbf{AVANCEMENT $\xi$}\\\textit{Tableau d'avancement}};
\node[subconcept] (tau) at (3, -4) {$\tau = \xi_\text{eq}/\xi_\text{max}$};
\node[subconcept] (alpha) at (-3, -4) {$\alpha = \xi_\text{eq}/n_\text{initial}$};

\draw[link] (centre) -- (pH);
\draw[link] (centre) -- (avancement);
\draw[link] (avancement) -- (tau);
\draw[link] (avancement) -- (alpha);

% Branche 4 : Prédominance
\node[subconcept] (predom) at (0, -6) {\textbf{DOMAINES DE PRÉDOMINANCE}};
\node[subconcept] (acide) at (-3, -7) {$\mathrm{pH} < \mathrm{p}K_\mathrm{a}$ : \ce{AH} major.};
\node[subconcept] (base) at (3, -7) {$\mathrm{pH} > \mathrm{p}K_\mathrm{a}$ : \ce{A-} major.};
\node[subconcept] (egal) at (0, -7) {$\mathrm{pH} = \mathrm{p}K_\mathrm{a}$ : $[\ce{AH}] = [\ce{A-}]$};

\draw[link] (centre) -- (predom);
\draw[link] (predom) -- (acide);
\draw[link] (predom) -- (base);
\draw[link] (predom) -- (egal);

% Branche 5 : Applications
\node[subconcept] (dosage) at (-7, -3) {\textbf{DOSAGE pH-métrique}};
\node[subconcept] (demi) at (-7, -4.5) {Demi-équivalence : $\mathrm{pH}=\mathrm{p}K_\mathrm{a}$};
\node[subconcept] (vin) at (-7, -6) {Degré vinaigre : g/100g};

\draw[link] (pKa) -- (demi);
\draw[link] (demi) -- (dosage);
\draw[link] (dosage) -- (vin);

% Liens transverses
\draw[bidilink] (faible) -- (Ka);
\draw[bidilink] (ionisation) -- (alpha);
\draw[bidilink] (totale) -- (tau);
\end{tikzpicture}
\legende{Carte mentale récapitulative de la leçon 7 : force, constantes, pH, avancement, prédominance et applications (dosage, vinaigre).}
\end{popfigure}

\begin{aretenir}[Synthèse des formules-clés pour le Bac]
\begin{itemize}
\item \textbf{Acide fort / Base forte :} Ionisation totale. $[\ce{H3O+}] = C_\text{acide}$ (monoprotique) ; $\mathrm{pH} = -\log C_\text{acide}$.
\item \textbf{Acide faible / Base faible :} Équilibre. Tableau d'avancement $\rightarrow$ $K_\mathrm{a} = \dfrac{[\ce{H3O+}]_\text{eq}^2}{C - [\ce{H3O+}]_\text{eq}}$.
\item \textbf{Coefficient d'ionisation :} $\alpha = \dfrac{[\ce{H3O+}]_\text{eq}}{C} = \dfrac{\xi_\text{eq}}{n_\text{initial}}$.
\item \textbf{Taux d'avancement :} $\tau = \dfrac{\xi_\text{eq}}{\xi_\text{max}}$.
\item \textbf{Relation pH–pKa :} $\mathrm{pH} = \mathrm{p}K_\mathrm{a} + \log\dfrac{[\ce{A-}]}{[\ce{AH}]}$.
\item \textbf{Demi-équivalence :} $\mathrm{pH} = \mathrm{p}K_\mathrm{a}$ (dosage acide faible/base forte).
\item \textbf{Degré d'acidité (vinaigre) :} $\text{Degré} = \dfrac{C_\text{mère} \times M_{\ce{CH3COOH}}}{10 \times \rho}$ (avec $C_\text{mère}$ en mol/L, $M$ en g/mol, $\rho$ en g/mL). \textbf{N'oubliez pas le facteur de dilution !}
\end{itemize}
\end{aretenir}

\begin{exercice}[Entraînement Bac : Vinaigre de bissap]
On souhaite contrôler la teneur en acide éthanoïque d'un vinaigre de bissap (hibiscus) artisanal.
\begin{enumerate}
\item On dilue \SI{5,0}{mL} de vinaigre mère dans une fiole jaugée de \SI{100,0}{mL}.
\item On prélève \SI{20,0}{mL} de cette solution diluée que l'on dose par une solution de soude \SI{0,050}{mol.L^{-1}}.
\item L'équivalence est atteinte pour $V_\text{E} = \SI{18,4}{mL}$ de soude.
\item La masse volumique du vinaigre est $\rho = \SI{1,00}{g.mL^{-1}}$.
\end{enumerate}
Calculer le degré d'acidité de ce vinaigre. Est-il conforme à la norme (6°–8°) ?
\begin{corrige}[Exercice : Vinaigre de bissap]
\textbf{1. Facteur de dilution :} $f = \dfrac{100,0}{5,0} = 20$.
\textbf{2. Quantité de soude :} $n_{\ce{HO-}} = \SI{0,050}{mol.L^{-1}} \times \SI{18,4e-3}{L} = \SI{9,2e-4}{mol}$.
\textbf{3. Quantité d'acide dans l'aliquot :} $n_{\ce{CH3COOH}}(\text{dilué}) = \SI{9,2e-4}{mol}$.
\textbf{4. Concentration diluée :} $C_\text{diluée} = \dfrac{\SI{9,2e-4}{mol}}{\SI{20,0e-3}{L}} = \SI{0,0460}{mol.L^{-1}}$.
\textbf{5. Concentration mère :} $C_\text{mère} = 20 \times \SI{0,0460}{mol.L^{-1}} = \SI{0,920}{mol.L^{-1}}$.
\textbf{6. Masse d'acide par litre :} $m = \SI{0,920}{mol.L^{-1}} \times \SI{60,05}{g.mol^{-1}} = \SI{55,25}{g.L^{-1}}$.
\textbf{7. Masse de 1 L vinaigre :} $m_\text{sol} = \SI{1,00}{g.mL^{-1}} \times 1000 = \SI{1000}{g}$.
\textbf{8. Degré :} $\dfrac{\SI{55,25}{g}}{\SI{1000}{g}} \times 100 = \textbf{5,5°}$.
\textbf{Conclusion :} Le vinaigre est à \textbf{5,5°}, \textbf{non conforme} (inférieur à 6°). Il est trop faiblement acide pour une conservation sûre.
\end{corrige}
\end{exercice}

Cette leçon s'achève. Vous maîtrisez maintenant la distinction fondamentale entre force et concentration, l'outil « tableau d'avancement » pour quantifier les équilibres, la puissance prédictive du couple $\mathrm{pH}/\mathrm{p}K_\mathrm{a}$, et l'application concrète au dosage du vinaigre — un savoir-faire directement valorisable dans l'industrie agroalimentaire camerounaise. Bonnes révisions pour le Baccalauréat !

\newpage

\coursec{Travaux pratiques}

\courssub{Objectifs du TP}

À l'issue de cette séance de travaux pratiques, tu seras capable de :

\begin{itemize}
\item vérifier expérimentalement le caractère \cle{faible} de l'acide éthanoïque contenu dans le vinaigre ;
\item réaliser un \cle{dosage pH-métrique} acide faible -- base forte et tracer la courbe $\mathrm{pH} = f(V_b)$ ;
\item déterminer le volume équivalent $V_E$ par la \cle{méthode des tangentes} et par le virage de la phénolphtaléine ;
\item lire le $\mathrm{p}K_a$ du couple \ce{CH3COOH}/\ce{CH3COO-} à la \cle{demi-équivalence} ;
\item calculer la concentration en acide éthanoïque du vinaigre commercial et son \cle{degré d'acidité}, puis comparer la valeur obtenue à celle portée sur l'étiquette.
\end{itemize}

\begin{aretenir}[Rappels utiles avant de commencer]
\begin{itemize}
\item Réaction de dosage : \[ \ce{CH3COOH + HO- -> CH3COO- + H2O} \] Elle est totale, rapide et unique : c'est bien une réaction de dosage.
\item À l'équivalence : $C_a \times V_a = C_b \times V_E$.
\item À la demi-équivalence ($V_b = V_E/2$) : $[\ce{CH3COOH}] = [\ce{CH3COO-}]$, donc $\mathrm{pH} = \mathrm{p}K_a$.
\item Le degré d'acidité d'un vinaigre est la masse d'acide éthanoïque (en grammes) contenue dans \SI{100}{g} de vinaigre (assimilé à \SI{100}{mL}).
\end{itemize}
\end{aretenir}

\courssub{Matériel et produits}

\begin{center}
\rowcolors{2}{popblueL}{white}
\begin{tabularx}{\linewidth}{|l|X|}
\hline
\rowcolor{popblue!40} \textbf{Matériel} & \textbf{Produits} \\
\hline
Burette graduée de \SI{25}{mL} & Solution d'hydroxyde de sodium (soude) à $C_b = \SI{0,10}{mol.L^{-1}}$ \\
Béchers de \SI{100}{mL} et \SI{250}{mL} & Vinaigre commercial du marché (degré indiqué sur l'étiquette, souvent $5^{\circ}$ ou $6^{\circ}$) \\
Pipette jaugée de \SI{20,0}{mL} + poire aspirante & Phénolphtaléine (solution alcoolique) \\
Fiole jaugée de \SI{100}{mL} & Eau distillée (ou eau de puits bouillie puis refroidie et filtrée) \\
pH-mètre étalonné (ou papier pH de précision) & \\
Agitateur magnétique ou agitateur en verre & \\
Lunettes de protection, gants, blouse & \\
\hline
\end{tabularx}
\end{center}

\begin{savaistu}[Et si le pH-mètre manque ?]
Dans beaucoup de lycées camerounais, le pH-mètre est rare ou en panne. Deux solutions :
\begin{itemize}
\item utiliser du \textbf{papier pH de précision} (bandes graduées de 0,5 unité de pH) : on prélève une goutte de solution avec un agitateur propre après chaque ajout de soude ;
\item réaliser un \textbf{dosage colorimétrique simple} avec la phénolphtaléine seule : on repère l'équivalence au virage (rose pâle persistant), puis on admet la valeur tabulée $\mathrm{p}K_a = 4{,}8$ pour la partie théorique.
\end{itemize}
Le vinaigre blanc du marché de Mokolo ou de Nkouloulou convient parfaitement ; évite les vinaigres colorés qui gênent l'observation du virage.
\end{savaistu}

\courssub{Sécurité}

\begin{attention}[Consignes de sécurité]
\begin{itemize}
\item \textbf{Soude \SI{0,10}{mol.L^{-1}}} : produit \emph{corrosif} (pictogramme : deux gouttes attaquant une main et une surface métallique). Porte obligatoirement \textbf{lunettes, gants et blouse}. En cas de contact avec la peau, rince abondamment à l'eau froide pendant plusieurs minutes et préviens l'enseignant.
\item \textbf{Phénolphtaléine en solution alcoolique} : liquide \emph{inflammable} (pictogramme : flamme). Tiens-la éloignée de toute flamme (bec Bunsen éteint pendant le dosage).
\item \textbf{Pipetage} : utilise toujours la \textbf{poire aspirante} (propipette) ; il est interdit de pipeter avec la bouche.
\item \textbf{Verrerie} : manipule la burette et la fiole jaugée à deux mains ; signale immédiatement toute casse.
\item Ne goûte \textbf{jamais} les solutions, même le vinaigre dilué, dans un laboratoire.
\item En fin de séance, neutralise les rejets (ils sont proches de la neutralité après dosage) avant de les verser à l'évier avec un grand volume d'eau.
\end{itemize}
\end{attention}

\courssub{Schéma du dispositif}

\begin{popfigure}
\begin{center}
\begin{tikzpicture}[scale=1.0, every node/.style={font=\small}]
% Support et burette
\draw[thick, popdark] (-3.2,0) -- (-3.2,5.6);
\draw[thick, popdark] (-3.9,0) -- (-2.5,0);
\draw[thick, popdark] (-3.2,4.6) -- (-2.2,4.6);
\fill[popdark] (-2.3,4.5) rectangle (-2.1,4.7);
% Burette
\draw[thick, popblue] (-1.6,5.4) -- (-1.6,2.6);
\draw[thick, popblue] (-1.2,5.4) -- (-1.2,2.6);
\draw[thick, popblue] (-1.6,2.6) -- (-1.5,2.2);
\draw[thick, popblue] (-1.2,2.6) -- (-1.3,2.2);
\draw[thick, popblue] (-1.5,2.2) -- (-1.5,1.9);
\draw[thick, popblue] (-1.3,2.2) -- (-1.3,1.9);
\draw[thick, popblue] (-1.5,1.9) -- (-1.3,1.9);
% Robinet
\draw[thick, popdark] (-1.75,2.35) -- (-1.05,2.35);
% Liquide dans burette
\fill[popblue!25] (-1.58,5.0) rectangle (-1.22,2.62);
\draw[popblue] (-1.58,5.0) -- (-1.22,5.0);
% Graduations
\foreach \y in {2.9,3.3,3.7,4.1,4.5,4.9} \draw[popdark] (-1.6,\y) -- (-1.5,\y);
\node[popblue, anchor=west] at (-0.9,4.6) {Soude $C_b = \SI{0,10}{mol.L^{-1}}$};
% Goutte
\fill[popblue!50] (-1.4,1.55) circle (0.06);
% Bécher
\draw[thick, popdark] (-2.6,0) -- (-2.6,1.5) -- (-0.2,1.5) -- (-0.2,0) -- cycle;
\fill[popgold!25] (-2.55,0.05) rectangle (-0.25,0.95);
\draw[popgold] (-2.55,0.95) -- (-0.25,0.95);
\node[popdark] at (-1.4,-0.35) {Vinaigre dilué $V_a = \SI{20,0}{mL}$ + phénolphtaléine};
% Sonde pH-mètre
\draw[thick, poppurple] (0.6,2.6) -- (0.6,0.4);
\draw[thick, poppurple] (0.85,2.6) -- (0.85,0.4);
\draw[thick, poppurple] (0.6,0.4) -- (0.85,0.4);
\fill[poppurple!30] (0.62,0.42) rectangle (0.83,0.7);
\draw[thick, poppurple] (0.72,2.6) -- (0.72,3.0);
% pH-mètre boîtier
\draw[thick, poppurple, rounded corners] (0.1,3.0) rectangle (1.9,3.9);
\node[poppurple] at (1.0,3.6) {pH-mètre};
\node[poppurple, font=\footnotesize] at (1.0,3.25) {pH = 2{,}9};
% Barreau magnétique
\fill[popdark] (-1.7,0.12) rectangle (-1.1,0.22);
% Agitateur
\draw[thick, popmuted] (-2.9,-0.75) rectangle (0.1,-0.5);
\node[popmuted, font=\footnotesize] at (-1.4,-0.62) {agitateur magnétique};
\end{tikzpicture}
\end{center}
\legende{Dispositif de dosage pH-métrique du vinaigre dilué par la soude. La sonde du pH-mètre plonge dans le bécher sans toucher le barreau aimanté.}
\end{popfigure}

\courssub{Protocole}

\textbf{Étape 1 -- Dilution du vinaigre (dilution au dixième).}\par
Le vinaigre commercial est trop concentré et trop coloré pour être dosé directement.
\begin{enumerate}
\item À l'aide d'une pipette jaugée de \SI{10,0}{mL} munie d'une poire, prélève \SI{10,0}{mL} de vinaigre commercial.
\item Verse dans une fiole jaugée de \SI{100,0}{mL}, complète avec de l'eau distillée jusqu'au trait de jauge, bouche et homogénéise en retournant dix fois. Tu obtiens la \textbf{solution S}, dix fois moins concentrée que le vinaigre.
\end{enumerate}

\textbf{Étape 2 -- Mise en place du dosage.}
\begin{enumerate}
\setcounter{enumi}{2}
\item Rince la burette avec un peu de solution de soude, puis remplis-la avec la soude à $C_b = \SI{0,10}{mol.L^{-1}}$. Ajuste le zéro en vérifiant l'absence de bulle dans le bec.
\item Pipette $V_a = \SI{20,0}{mL}$ de solution S et verse-les dans un bécher de \SI{100}{mL}. Ajoute 2 à 3 gouttes de phénolphtaléine et environ \SI{20}{mL} d'eau distillée pour que la sonde plonge correctement (l'ajout d'eau ne modifie pas l'équivalence).
\item Place le bécher sur l'agitateur, introduis la sonde du pH-mètre (préalablement étalonné avec les solutions tampons pH 4 et pH 7) et le barreau aimanté. Mets l'agitation en route, modérée.
\end{enumerate}

\textbf{Étape 3 -- Mesures.}
\begin{enumerate}
\setcounter{enumi}{5}
\item Relève le pH initial ($V_b = 0$).
\item Ajoute la soude par fractions de \SI{1}{mL} en relevant le pH après chaque ajout, en attendant sa stabilisation.
\item Dès que le pH commence à varier fortement (vers pH $\approx 5{,}5$--6), réduis les ajouts à \SI{0,2}{mL} (ou goutte à goutte) pour bien décrire le saut de pH. Note le volume au virage persistant de la phénolphtaléine (rose pâle).
\item Poursuis jusqu'à $V_b = \SI{25}{mL}$.
\end{enumerate}

\textbf{Étape 4 -- Exploitation.} Trace la courbe $\mathrm{pH} = f(V_b)$ sur papier millimétré, détermine $V_E$ par la méthode des tangentes, lis le pH à $V_E/2$, puis calcule.

\courssub{Observations et résultats attendus}

\begin{itemize}
\item La solution initiale est incolore ; son pH initial vaut environ $2{,}9$ : bien supérieur à celui d'un acide fort de même concentration ($\mathrm{pH} = 1$ pour \SI{0,1}{mol.L^{-1}}), ce qui traduit déjà le caractère faible de l'acide éthanoïque.
\item Le pH croît lentement, puis brutalement autour de l'équivalence (saut de pH de pH $\approx 7$ à pH $\approx 11$), puis se stabilise.
\item La phénolphtaléine vire au rose pâle persistant légèrement \emph{après} le début du saut, vers pH $\approx 8{,}5$ : sa zone de virage (8,2 -- 10) est bien adaptée car le pH à l'équivalence est supérieur à 7 (solution d'éthanoate de sodium, basique).
\end{itemize}

\begin{center}
\rowcolors{2}{popgoldL}{white}
\begin{tabularx}{\linewidth}{|c|X|c|X|}
\hline
\rowcolor{popgold!50} $V_b$ (mL) & pH & $V_b$ (mL) & pH \\
\hline
0,0 & 2,9 & 11,0 & 5,2 \\
2,0 & 3,5 & 11,5 & 5,4 \\
4,0 & 3,9 & 12,0 & 5,6 \\
6,0 & 4,3 & 12,5 & 5,8 \\
8,0 & 4,7 & 13,0 & 6,1 \\
8,3 & 4,8 & 14,0 & 6,6 \\
9,0 & 4,9 & 15,0 & 7,3 \\
10,0 & 5,0 & 16,0 & 8,7 \\
\hline
\end{tabularx}
\end{center}

Tableau de mesures type (vinaigre dilué 10 fois, $C_a \approx \SI{0,083}{mol.L^{-1}}$, $V_E \approx \SI{16,6}{mL}$). Le virage de la phénolphtaléine est observé vers $V_b \approx \SI{16,8}{mL}$.

\courssub{Exploitation}

\begin{methode}[Tableau d'avancement et détermination de $K_a$ à l'état initial]
Pour mettre en œuvre le savoir-faire « utiliser la notion d'avancement pour étudier un équilibre chimique », on modélise l'ionisation de l'acide éthanoïque dans la solution S avant tout ajout de soude.
\begin{enumerate}
\item \textbf{Équation de l'équilibre :} $\ce{CH3COOH + H2O <=> CH3COO- + H3O+}$.
\item \textbf{État initial :} $n_{\ce{CH3COOH}} = C_a V_a$ ; $n_{\ce{CH3COO-}} = 0$ ; $n_{\ce{H3O+}} \approx 0$ (négligeable devant $x$).
\item \textbf{Avancement $x$ :} à l'équilibre, $n_{\ce{CH3COOH}} = C_a V_a - x$ ; $n_{\ce{CH3COO-}} = x$ ; $n_{\ce{H3O+}} = x$.
\item \textbf{Concentrations à l'équilibre :} $[\ce{CH3COOH}] = \frac{C_a V_a - x}{V_a} \approx C_a - \frac{x}{V_a}$ ; $[\ce{CH3COO-}] = [\ce{H3O+}] = \frac{x}{V_a}$.
\item \textbf{Expression de $K_a$ :} $K_a = \frac{[\ce{CH3COO-}][\ce{H3O+}]}{[\ce{CH3COOH}]} = \frac{(x/V_a)^2}{C_a - x/V_a}$.
\item \textbf{Application numérique :} pH initial $= 2{,}9 \Rightarrow [\ce{H3O+}] = 10^{-2,9} \approx \SI{1,26e-3}{mol.L^{-1}} = x/V_a$.
$C_a \approx \SI{0,083}{mol.L^{-1}}$.
$K_a = \frac{(1,26 \times 10^{-3})^2}{0,083 - 1,26 \times 10^{-3}} \approx \frac{1,59 \times 10^{-6}}{0,0817} \approx \SI{1,95e-5}{}$.
$\mathrm{p}K_a = -\log K_a \approx 4{,}71$, valeur proche de celle trouvée par titration ($4{,}8$).
\end{enumerate}
\end{methode}

\begin{methode}[Tableau d'avancement de la réaction de dosage (totale)]
La réaction de dosage $\ce{CH3COOH + HO- -> CH3COO- + H2O}$ est totale. On utilise un tableau d'avancement pour justifier la relation à l'équivalence.
\begin{center}
\begin{tabular}{|c|c|c|c|c|}
\hline
 & \ce{CH3COOH} & \ce{HO-} & \ce{CH3COO-} & \ce{H2O} \\
\hline
État initial & $n_a = C_a V_a$ & $n_b = C_b V_b$ & 0 & --- \\
\hline
Changement & $-x$ & $-x$ & $+x$ & --- \\
\hline
État final & $n_a - x$ & $n_b - x$ & $x$ & --- \\
\hline
\end{tabular}
\end{center}
À l'équivalence, les réactifs sont introduits selon les proportions stœchiométriques : $x_{\max} = n_a = n_b$, d'où $C_a V_a = C_b V_E$.
\end{methode}

\begin{exoresolu}[Exploitation guidée du dosage]
\textbf{Énoncé.} On dose $V_a = \SI{20,0}{mL}$ de vinaigre dilué 10 fois par la soude à $C_b = \SI{0,10}{mol.L^{-1}}$. La méthode des tangentes donne $V_E = \SI{16,6}{mL}$ et le pH à la demi-équivalence vaut $4{,}8$. L'étiquette du vinaigre annonce $5^{\circ}$.
\begin{enumerate}
\item Pourquoi le pH initial ($2{,}9$) prouve-t-il que l'acide éthanoïque est un acide faible ?
\item Déduis du graphe le $\mathrm{p}K_a$ du couple \ce{CH3COOH}/\ce{CH3COO-}.
\item Calcule la concentration $C_a$ de la solution diluée S, puis la concentration $C_0$ du vinaigre commercial.
\item Calcule le degré d'acidité du vinaigre et compare-le à l'étiquette. Donnée : $M(\ce{CH3COOH}) = \SI{60}{g.mol^{-1}}$.
\end{enumerate}
\tcblower
\textbf{Solution détaillée.}
\begin{enumerate}
\item Si l'acide était fort et totalement ionisé, on aurait $\mathrm{pH} = -\log C_a$. Or la solution S, de concentration voisine de \SI{0,08}{mol.L^{-1}}, donnerait $\mathrm{pH} \approx 1{,}1$. Le pH mesuré ($2{,}9$) est nettement supérieur : $[\ce{H3O+}] = 10^{-2,9} \approx \SI{1,26e-3}{mol.L^{-1}} \ll C_a$. L'acide éthanoïque n'est donc que \textbf{partiellement ionisé} : c'est un \cle{acide faible}. Son coefficient d'ionisation est $\alpha = \dfrac{[\ce{H3O+}]}{C_a} \approx \dfrac{1,26 \times 10^{-3}}{0,083} \approx 1,5\,\%$.
\item À la demi-équivalence, $V_b = V_E/2 = \SI{8,3}{mL}$, les concentrations des deux formes du couple sont égales : $[\ce{CH3COOH}] = [\ce{CH3COO-}]$, donc $\mathrm{pH} = \mathrm{p}K_a$. La lecture graphique (tableau : $V_b = 8,3$ mL $\rightarrow$ pH $= 4,8$) donne $\boxed{\mathrm{p}K_a(\ce{CH3COOH}/\ce{CH3COO-}) = 4{,}8}$, en accord avec la valeur tabulée.
\item À l'équivalence : $C_a V_a = C_b V_E$, donc
\[ C_a = \frac{C_b \times V_E}{V_a} = \frac{0,10 \times 16,6}{20,0} = \SI{0,083}{mol.L^{-1}}. \]
Le vinaigre commercial est 10 fois plus concentré : $C_0 = 10 \times C_a = \SI{0,83}{mol.L^{-1}}$.
\item Concentration massique : $C_m = C_0 \times M = 0,83 \times 60 = \SI{49,8}{g.L^{-1}} \approx \SI{50}{g.L^{-1}}$, soit \SI{5,0}{g} d'acide éthanoïque pour \SI{100}{mL} de vinaigre. Le \textbf{degré d'acidité} est donc d'environ $\boxed{5^{\circ}}$, en parfait accord avec l'étiquette (écart relatif inférieur à $2\,\%$, imputable aux erreurs de lecture du volume équivalent).
\end{enumerate}
\end{exoresolu}

\begin{attention}[Erreurs fréquentes à éviter]
\begin{itemize}
\item \textbf{Oublier le facteur de dilution} : la concentration calculée à l'équivalence est celle de la solution \emph{diluée} ; il faut multiplier par 10 pour le vinaigre commercial.
\item \textbf{Confondre équivalence et neutralité} : à l'équivalence, le pH n'est pas 7 mais environ 8,7, car la solution contient la base faible \ce{CH3COO-} qui se hydrolyse : $\ce{CH3COO- + H2O <=> CH3COOH + HO-}$.
\item \textbf{Lire le $\mathrm{p}K_a$ au mauvais endroit} : c'est à $V_E/2$ et non à $V_E$ que $\mathrm{pH} = \mathrm{p}K_a$.
\item \textbf{Mal étalonner le pH-mètre} : sans étalonnage préalable (tampons pH 4 et 7), toutes les mesures sont faussées.
\item \textbf{Ajouter la soude trop vite près du saut} : on risque de « rater » l'équivalence ; réduis les volumes ajoutés dès que le pH varie sensiblement.
\end{itemize}
\end{attention}

\courssub{Conclusion}

Ce TP t'a permis de mettre en œuvre l'essentiel des compétences de la leçon. Le pH initial de la solution de vinaigre dilué, très supérieur à celui qu'aurait un acide fort de même concentration, confirme que l'acide éthanoïque est un \cle{acide faible}, partiellement ionisé dans l'eau ($\alpha \approx 1,5\,\%$). Le tracé de la courbe $\mathrm{pH} = f(V_b)$ montre les quatre régions caractéristiques du dosage d'un acide faible par une base forte : départ à pH modéré, zone tampon autour de la demi-équivalence, saut de pH moins étendu que pour un acide fort, et pH à l'équivalence supérieur à 7. La lecture à la demi-équivalence fournit $\mathrm{p}K_a = 4{,}8$ pour le couple \ce{CH3COOH}/\ce{CH3COO-}, valeur conforme aux tables. Enfin, le calcul du degré d'acidité ($\approx 5^{\circ}$) valide l'indication de l'étiquette du vinaigre du marché : la chimie acido-basique est ainsi un véritable outil de \textbf{contrôle de qualité} des produits de consommation courante, comme le font les laboratoires d'analyse des denrées alimentaires au Cameroun.

\newpage

\coursec{Méthodes et exercices résolus}

\coursec{Force des acides et des bases, couples acide/base}

\courssub{Montrer qu'un acide (ou une base) est fort(e) à partir du pH et de la concentration}

\begin{methode}[Démontrer le caractère fort d'un acide ou d'une base]
\begin{enumerate}
\item \textbf{Calculer la concentration théorique en ions.} Pour un monoacide de concentration $C$ : si l'acide était fort, $[\ce{H3O+}] = C$, donc $\mathrm{pH}_{\text{théorique}} = -\log C$. Pour une monobase forte : $[\ce{HO-}] = C$, d'où $\mathrm{pH} = 14 + \log C$ (à \SI{25}{\celsius}, $K_e = 10^{-14}$).
\item \textbf{Comparer avec le pH mesuré.} Si $\mathrm{pH}_{\text{mesuré}} = \mathrm{pH}_{\text{théorique}}$ (aux incertitudes près), l'acide (ou la base) est \cle{totalement ionisé(e)} : il est \cle{fort(e)}.
\item \textbf{Conclure par l'équation de dissolution.} Écrire la réaction avec une flèche simple $\ce{->}$ : \ce{HCl + H2O -> H3O+ + Cl-} ou \ce{NaOH -> Na+ + HO-}.
\item \textbf{Réflexe de vérification :} calculer le taux d'avancement final $\tau = \dfrac{[\ce{H3O+}]}{C} = \dfrac{10^{-\mathrm{pH}}}{C}$ (acide) ou $\tau = \dfrac{[\ce{HO-}]}{C}$ (base). Si $\tau \approx 1$ (soit $100\,\%$), l'espèce est forte.
\end{enumerate}
\end{methode}

\begin{exoresolu}[Acide chlorhydrique et soude : sont-ils forts ?]
On mesure à \SI{25}{\celsius} le pH de deux solutions : une solution $S_1$ d'acide chlorhydrique de concentration $C_1 = \SI{1,0e-2}{mol.L^{-1}}$ donne $\mathrm{pH}_1 = 2,0$ ; une solution $S_2$ d'hydroxyde de sodium de concentration $C_2 = \SI{5,0e-3}{mol.L^{-1}}$ donne $\mathrm{pH}_2 = 11,7$. Montrer que l'acide chlorhydrique est un acide fort et que la soude est une base forte.
\tcblower
\textbf{Étude de $S_1$ (acide chlorhydrique).}
Si l'acide est fort, $[\ce{H3O+}] = C_1$ donc :
\[ \mathrm{pH}_{\text{théorique}} = -\log C_1 = -\log(1,0\times 10^{-2}) = 2,0. \]
Or $\mathrm{pH}_1 = 2,0 = \mathrm{pH}_{\text{théorique}}$. De plus :
\[ \tau = \frac{10^{-\mathrm{pH}_1}}{C_1} = \frac{10^{-2,0}}{1,0\times 10^{-2}} = 1,00. \]
L'acide chlorhydrique est \textbf{totalement ionisé} : c'est un \textbf{acide fort} :
\[ \ce{HCl + H2O -> H3O+ + Cl-} \]
\textbf{Étude de $S_2$ (soude).}
Si la base est forte, $[\ce{HO-}] = C_2$, donc $[\ce{H3O+}] = \dfrac{K_e}{C_2}$ et :
\[ \mathrm{pH}_{\text{théorique}} = 14 + \log C_2 = 14 + \log(5,0\times 10^{-3}) = 14 - 2,30 = 11,7. \]
Or $\mathrm{pH}_2 = 11,7 = \mathrm{pH}_{\text{théorique}}$. La soude est \textbf{totalement dissociée} : c'est une \textbf{base forte} :
\[ \ce{NaOH -> Na+ + HO-} \]
\textbf{Conclusion :} les deux valeurs de pH coïncident avec les formules des espèces fortes ; l'acide chlorhydrique et l'hydroxyde de sodium sont respectivement un acide fort et une base forte.
\end{exoresolu}

\courssub{Montrer qu'un acide ou une base est faible}

\begin{methode}[Démontrer le caractère faible d'un acide ou d'une base]
\begin{enumerate}
\item \textbf{Calculer le pH qu'aurait la solution si l'espèce était forte} : $\mathrm{pH} = -\log C$ (acide) ou $\mathrm{pH} = 14 + \log C$ (base).
\item \textbf{Comparer au pH mesuré.} Pour un acide : si $\mathrm{pH}_{\text{mesuré}} > -\log C$, alors $[\ce{H3O+}] < C$ : l'acide est \cle{partiellement ionisé}, donc \cle{faible}. Pour une base : si $\mathrm{pH}_{\text{mesuré}} < 14 + \log C$, la base est faible.
\item \textbf{Quantifier avec le coefficient d'ionisation} $\alpha = \tau = \dfrac{[\ce{H3O+}]}{C} < 1$ (acide) ou $\alpha = \dfrac{[\ce{HO-}]}{C} < 1$ (base).
\item \textbf{Écrire l'équilibre} avec une double flèche : \ce{AH + H2O <=> A- + H3O+}.
\end{enumerate}
\end{methode}

\begin{exoresolu}[L'acide éthanoïque du vinaigre est-il faible ?]
Une solution d'acide éthanoïque \ce{CH3COOH} de concentration $C = \SI{1,0e-2}{mol.L^{-1}}$ a un pH de $3,4$ à \SI{25}{\celsius}. Montrer que l'acide éthanoïque est un acide faible et calculer son coefficient d'ionisation $\alpha$.
\tcblower
\textbf{Hypothèse d'un acide fort :} on aurait $\mathrm{pH} = -\log C = -\log(1,0\times 10^{-2}) = 2,0$.
\textbf{Comparaison :} $\mathrm{pH}_{\text{mesuré}} = 3,4 > 2,0$. La concentration réelle en ions oxonium est :
\[ [\ce{H3O+}] = 10^{-\mathrm{pH}} = 10^{-3,4} = \SI{4,0e-4}{mol.L^{-1}} < C = \SI{1,0e-2}{mol.L^{-1}}. \]
L'acide n'est donc \textbf{pas totalement ionisé} : c'est un \textbf{acide faible}, qui réagit selon l'équilibre :
\[ \ce{CH3COOH + H2O <=> CH3COO- + H3O+} \]
\textbf{Coefficient d'ionisation :}
\[ \alpha = \frac{[\ce{H3O+}]}{C} = \frac{4,0\times 10^{-4}}{1,0\times 10^{-2}} = 4,0\times 10^{-2} = 0,040. \]
\textbf{Conclusion :} seulement $4,0\,\%$ des molécules d'acide éthanoïque sont ionisées : l'acide éthanoïque est bien un acide faible ($\alpha \ll 1$).
\end{exoresolu}

\courssub{Étudier un équilibre avec le tableau d'avancement}

\begin{methode}[Exploiter l'avancement d'une réaction acido-basique]
\begin{enumerate}
\item \textbf{Écrire l'équation} de la réaction de l'acide \ce{AH} avec l'eau : \ce{AH + H2O <=> A- + H3O+}.
\item \textbf{Dresser le tableau d'avancement} en quantités de matière (ou en concentrations) : état initial ($n_0 = C\,V$), état intermédiaire ($x$), état final ($x_f$).
\item \textbf{Exploiter la mesure de pH} : $[\ce{H3O+}]_f = 10^{-\mathrm{pH}}$, d'où $x_f = [\ce{H3O+}]_f \times V$ (attention : négliger les \ce{H30+} de l'autoprotolyse de l'eau seulement si pH $< 6,5$).
\item \textbf{Calculer l'avancement maximal} $x_{\max} = n_0$ (réaction supposée totale) puis le \cle{taux d'avancement final} $\tau = \dfrac{x_f}{x_{\max}}$.
\item \textbf{Conclure} : $\tau < 1$ $\Rightarrow$ équilibre (réaction limitée) ; $\tau = 1$ $\Rightarrow$ réaction totale.
\end{enumerate}
\end{methode}

\begin{exoresolu}[Taux d'avancement de l'ionisation de l'acide méthanoïque]
On prépare $V = \SI{500}{mL}$ d'une solution d'acide méthanoïque \ce{HCOOH} de concentration $C = \SI{2,0e-2}{mol.L^{-1}}$. Le pH mesuré vaut $2,7$. À l'aide d'un tableau d'avancement, déterminer l'avancement final $x_f$, l'avancement maximal $x_{\max}$ et le taux d'avancement final $\tau$. Conclure.
\tcblower
\textbf{Équation :} \ce{HCOOH + H2O <=> HCOO- + H3O+}.
\textbf{Quantité initiale :} $n_0 = C\,V = 2,0\times 10^{-2} \times 0,500 = \SI{1,0e-2}{mol}$.
\textbf{Tableau d'avancement :}
\begin{center}
\begin{tabularx}{\linewidth}{|l|X|X|X|X|}
\hline
\rowcolor{popblueL} État & \ce{HCOOH} & \ce{H2O} & \ce{HCOO-} & \ce{H3O+} \\
\hline
Initial ($x=0$) & $n_0$ & excès & $0$ & $0$ \\
\hline
Intermédiaire ($x$) & $n_0 - x$ & excès & $x$ & $x$ \\
\hline
Final ($x_f$) & $n_0 - x_f$ & excès & $x_f$ & $x_f$ \\
\hline
\end{tabularx}
\end{center}
\textbf{Avancement final :} $[\ce{H3O+}]_f = 10^{-2,7} = \SI{2,0e-3}{mol.L^{-1}}$, donc :
\[ x_f = [\ce{H3O+}]_f \times V = 2,0\times 10^{-3} \times 0,500 = \SI{1,0e-3}{mol}. \]
\textbf{Avancement maximal :} $x_{\max} = n_0 = \SI{1,0e-2}{mol}$.
\textbf{Taux d'avancement final :}
\[ \tau = \frac{x_f}{x_{\max}} = \frac{1,0\times 10^{-3}}{1,0\times 10^{-2}} = 0,10 \quad \text{soit } 10\,\%. \]
\textbf{Conclusion :} $\tau = 0,10 < 1$ : la réaction de l'acide méthanoïque avec l'eau est \textbf{limitée} ; il s'agit d'un équilibre chimique et l'acide méthanoïque est un acide faible.
\end{exoresolu}

\courssub{Identifier les couples acide/base et calculer Ka, Kb, pKa}

\begin{methode}[Constantes d'acidité et de basicité]
\begin{enumerate}
\item \textbf{Identifier les couples} dans un équilibre : repérer les deux espèces qui ne diffèrent que par un proton \ce{H+} (ex. \ce{CH3COOH}/\ce{CH3COO-} et \ce{H3O+}/\ce{H2O}).
\item \textbf{Écrire la constante d'acidité} du couple \ce{AH}/\ce{A-} à partir de l'équilibre \ce{AH + H2O <=> A- + H3O+} :
\[ K_a = \frac{[\ce{A-}]\,[\ce{H3O+}]}{[\ce{AH}]}, \qquad \mathrm{p}K_a = -\log K_a. \]
\item \textbf{Constante de basicité} du couple \ce{AH}/\ce{A-} vu côté base (réaction de l'anion \ce{A-} avec l'eau) : \ce{A- + H2O <=> AH + HO-} :
\[ K_b = \frac{[\ce{AH}]\,[\ce{HO-}]}{[\ce{A-}]}. \]
\item \textbf{Relations à connaître} (à \SI{25}{\celsius}) : $K_a \times K_b = K_e = 10^{-14}$ et $\mathrm{p}K_a + \mathrm{p}K_b = 14$.
\item \textbf{Exploiter une mesure de pH} : pour un acide faible seul en solution, $[\ce{A-}] = [\ce{H3O+}] = 10^{-\mathrm{pH}}$ et $[\ce{AH}] = C - 10^{-\mathrm{pH}}$ (conservation de la matière).
\end{enumerate}
\end{methode}

\begin{exoresolu}[Constante d'acidité du couple de l'acide éthanoïque]
Une solution d'acide éthanoïque de concentration $C = \SI{1,0e-2}{mol.L^{-1}}$ a un pH de $3,4$. Identifier les couples mis en jeu, puis calculer $K_a$ et $\mathrm{p}K_a$ du couple \ce{CH3COOH}/\ce{CH3COO-}. En déduire $K_b$ de l'ion éthanoate.
\tcblower
\textbf{Couples mis en jeu :} \ce{CH3COOH}/\ce{CH3COO-} (l'acide cède un proton) et \ce{H3O+}/\ce{H2O} (l'eau capte le proton) :
\[ \ce{CH3COOH + H2O <=> CH3COO- + H3O+} \]
\textbf{Concentrations à l'équilibre :}
\begin{align*}
[\ce{H3O+}] &= [\ce{CH3COO-}] = 10^{-\mathrm{pH}} = 10^{-3,4} = \SI{4,0e-4}{mol.L^{-1}},\\
[\ce{CH3COOH}] &= C - [\ce{H3O+}] = 1,0\times 10^{-2} - 4,0\times 10^{-4} = \SI{9,6e-3}{mol.L^{-1}}.
\end{align*}
\textbf{Calcul de $K_a$ :}
\[ K_a = \frac{[\ce{CH3COO-}]\,[\ce{H3O+}]}{[\ce{CH3COOH}]} = \frac{(4,0\times 10^{-4})^2}{9,6\times 10^{-3}} = \frac{1,6\times 10^{-7}}{9,6\times 10^{-3}} = 1,7\times 10^{-5}. \]
\textbf{Calcul du $\mathrm{p}K_a$ :} $\mathrm{p}K_a = -\log(1,7\times 10^{-5}) = 4,8$.
\textbf{Constante de basicité de \ce{CH3COO-} :}
\[ K_b = \frac{K_e}{K_a} = \frac{1,0\times 10^{-14}}{1,7\times 10^{-5}} = 5,9\times 10^{-10} \quad (\mathrm{p}K_b = 14 - 4,8 = 9,2). \]
\textbf{Conclusion :} le couple \ce{CH3COOH}/\ce{CH3COO-} a pour constantes $K_a = 1,7\times 10^{-5}$ et $\mathrm{p}K_a = 4,8$ ; l'ion éthanoate est une base très faible ($K_b \approx 6\times 10^{-10}$).
\end{exoresolu}

\courssub{Étude de la base faible : l'ion éthanoate (hydrolyse)}

\begin{exemplebox}[pH d'une solution d'éthanoate de sodium]
On prépare une solution d'éthanoate de sodium \ce{CH3COONa} de concentration $C = \SI{1,0e-2}{mol.L^{-1}}$. L'ion éthanoate \ce{CH3COO-} est la base conjuguée de l'acide éthanoïque ($K_a = 1,7\times 10^{-5}$, $\mathrm{p}K_a = 4,8$). Prédire le pH de cette solution.
\tcblower
\textbf{Équation d'hydrolyse :} \ce{CH3COO- + H2O <=> CH3COOH + HO-}.
\textbf{Constante de basicité :} $K_b = \dfrac{K_e}{K_a} = \dfrac{10^{-14}}{1,7\times 10^{-5}} = 5,9\times 10^{-10}$.
\textbf{Tableau d'avancement (concentrations) :}
\begin{center}
\begin{tabularx}{\linewidth}{|l|X|X|X|X|}
\hline
\rowcolor{popblueL} État & \ce{CH3COO-} & \ce{H2O} & \ce{CH3COOH} & \ce{HO-} \\
\hline
Initial & $C$ & excès & $0$ & $0$ \\
\hline
Final & $C - x$ & excès & $x$ & $x$ \\
\hline
\end{tabularx}
\end{center}
Comme $K_b$ est très petit, $x \ll C$ : $C - x \approx C$.
\[ K_b = \frac{x^2}{C} \implies x = \sqrt{K_b C} = \sqrt{5,9\times 10^{-10} \times 1,0\times 10^{-2}} = \sqrt{5,9\times 10^{-12}} = 2,4\times 10^{-6}\ \text{mol.L}^{-1}. \]
$[\ce{HO-}] = x = \SI{2,4e-6}{mol.L^{-1}}$. $\mathrm{pOH} = -\log(2,4\times 10^{-6}) = 5,6$.
$\mathrm{pH} = 14 - \mathrm{pOH} = 8,4$.
\textbf{Conclusion :} une solution d'éthanoate de sodium est \textbf{basique} (pH > 7), ce qui caractérise une \textbf{base faible}.
\end{exemplebox}

\courssub{Classer les couples et prévoir le sens d'une réaction}

\begin{methode}[Classification des couples et prévision des réactions]
\begin{enumerate}
\item \textbf{Classer les acides} : plus $K_a$ est grand (donc plus $\mathrm{p}K_a$ est petit), plus l'acide est fort. Cas limites : acide fort ($K_a$ très grand, ex. \ce{H3O+}/\ce{H2O}, $\mathrm{p}K_a \approx 0$) ; acide inerte (\ce{H2O}/\ce{HO-}, $\mathrm{p}K_a = 14$).
\item \textbf{Classer les bases} : plus le $\mathrm{p}K_a$ du couple acide conjugué est grand, plus la base est forte.
\item \textbf{Prévoir le sens d'une réaction} entre l'acide $\ce{A1H}$ et la base $\ce{A2-}$ : écrire \ce{A1H + A2- <=> A1- + A2H} et calculer
\[ K = \frac{K_{a1}}{K_{a2}} = 10^{\mathrm{p}K_{a2} - \mathrm{p}K_{a1}}. \]
\item \textbf{Conclure} : si $K > 10^{4}$ la réaction est quasi totale ; si $K < 10^{-4}$ elle n'a pas lieu (c'est l'inverse qui se produit) ; entre les deux, c'est un équilibre. Réflexe : la réaction favorisée met en jeu \textbf{l'acide le plus fort} (plus petit $\mathrm{p}K_a$) et \textbf{la base la plus forte} (couple de plus grand $\mathrm{p}K_a$).
\end{enumerate}
\end{methode}

\begin{exoresolu}[Réaction entre l'acide éthanoïque et l'ammoniac]
On donne $\mathrm{p}K_a(\ce{CH3COOH}/\ce{CH3COO-}) = 4,8$ et $\mathrm{p}K_a(\ce{NH4+}/\ce{NH3}) = 9,2$. Classer ces deux couples par force d'acide croissante, écrire l'équation de la réaction entre l'acide éthanoïque et l'ammoniac, calculer sa constante $K$ et conclure sur son avancement.
\tcblower
\textbf{Classement :} plus le $\mathrm{p}K_a$ est petit, plus l'acide est fort :
\[ \ce{CH3COOH}\ (\mathrm{p}K_a = 4,8) \;>\; \ce{NH4+}\ (\mathrm{p}K_a = 9,2) \quad \text{(force d'acide)}. \]
Conséquence : \ce{NH3} est une base plus forte que \ce{CH3COO-}.
\textbf{Équation de la réaction :} l'acide le plus fort (\ce{CH3COOH}) réagit avec la base la plus forte (\ce{NH3}) :
\[ \ce{CH3COOH + NH3 <=> CH3COO- + NH4+} \]
\textbf{Constante d'équilibre :}
\[ K = \frac{[\ce{CH3COO-}]\,[\ce{NH4+}]}{[\ce{CH3COOH}]\,[\ce{NH3}]} = \frac{K_a(\ce{CH3COOH})}{K_a(\ce{NH4+})} = 10^{\mathrm{p}K_{a2} - \mathrm{p}K_{a1}} = 10^{9,2 - 4,8} = 10^{4,4} \approx 2,5\times 10^{4}. \]
\textbf{Conclusion :} $K \approx 2,5\times 10^{4} > 10^{4}$ : la réaction est \textbf{quasi totale} dans le sens direct. C'est bien la réaction entre l'acide le plus fort et la base la plus forte qui est favorisée.
\end{exoresolu}

\courssub{Utiliser la relation pH–pKa et les domaines de prédominance}

\begin{methode}[Relation de Henderson-Hasselbalch et diagramme de prédominance]
\begin{enumerate}
\item \textbf{Établir la relation} à partir de $K_a = \dfrac{[\ce{A-}][\ce{H3O+}]}{[\ce{AH}]}$ en passant au logarithme :
\[ \mathrm{pH} = \mathrm{p}K_a + \log\frac{[\ce{A-}]}{[\ce{AH}]}. \]
\item \textbf{Comparer pH et $\mathrm{p}K_a$} pour connaître l'espèce prédominante :
\begin{itemize}
\item si $\mathrm{pH} < \mathrm{p}K_a - 1$ : la forme \cle{acide} \ce{AH} prédomine ($[\ce{AH}] > 10[\ce{A-}]$) ;
\item si $\mathrm{pH} > \mathrm{p}K_a + 1$ : la forme \cle{basique} \ce{A-} prédomine ($[\ce{A-}] > 10[\ce{AH}]$) ;
\item si $\mathrm{pH} = \mathrm{p}K_a$ : $[\ce{AH}] = [\ce{A-}]$ (zone tampon optimale).
\end{itemize}
\item \textbf{Calculer un rapport de concentrations} : $\dfrac{[\ce{A-}]}{[\ce{AH}]} = 10^{\mathrm{pH} - \mathrm{p}K_a}$.
\item \textbf{Tracer le diagramme} : axe gradué en pH (0 à 14), frontière à $\mathrm{pH} = \mathrm{p}K_a$.
\end{enumerate}
\end{methode}

\begin{exoresolu}[Quelle forme du couple prédomine dans le jus de bissap ?]
Le jus de bissap (infusion de fleurs d'hibiscus) a un pH voisin de $3,0$. On considère le couple \ce{CH3COOH}/\ce{CH3COO-} ($\mathrm{p}K_a = 4,8$). Quelle espèce prédomine dans ce jus ? Calculer le rapport $\dfrac{[\ce{CH3COO-}]}{[\ce{CH3COOH}]}$. Tracer le diagramme de prédominance du couple.
\tcblower
\textbf{Comparaison pH / $\mathrm{p}K_a$ :} $\mathrm{pH} = 3,0 < \mathrm{p}K_a - 1 = 3,8$ : c'est la forme \textbf{acide} \ce{CH3COOH} qui prédomine.
\textbf{Calcul du rapport :}
\[ \frac{[\ce{CH3COO-}]}{[\ce{CH3COOH}]} = 10^{\mathrm{pH} - \mathrm{p}K_a} = 10^{3,0 - 4,8} = 10^{-1,8} \approx 1,6\times 10^{-2}. \]
Il y a donc environ $63$ fois plus de \ce{CH3COOH} que de \ce{CH3COO-} ($1/0{,}016 \approx 63$).
\textbf{Diagramme de prédominance :}
\begin{popfigure}
\begin{center}
\begin{tikzpicture}[scale=1.0]
% Axe pH 0-14 sur 14 cm
\draw[-{Stealth},thick,popdark] (0,0) -- (14.5,0) node[below left=2pt and -20pt]{\small pH};
\foreach \x/\lab in {0/0, 2/2, 4.8/4,8, 7/7, 10/10, 14/14}
  \draw[popdark] (\x,0.08) -- (\x,-0.08) node[below]{\small \lab};
% Zones de prédominance
\fill[popblue!35] (0,0.35) rectangle (4.8,0.95);
\fill[poporange!35] (4.8,0.35) rectangle (14,0.95);
\node at (2.4,0.65) {\small \textbf{\ce{CH3COOH} prédomine}};
\node at (9.4,0.65) {\small \textbf{\ce{CH3COO-} prédomine}};
% Frontière pKa
\draw[very thick,popink] (4.8,-0.15) -- (4.8,1.1);
\node[above,popink] at (4.8,1.1) {\small $\mathrm{pH}=\mathrm{p}K_a$};
% Point bissap
\draw[-{Stealth},thick,poppurple] (3.0,-0.55) -- (3.0,-0.15);
\node[below,poppurple] at (3.0,-0.55) {\small bissap : pH $= 3,0$};
\end{tikzpicture}
\end{center}
\legende{Diagramme de prédominance du couple \ce{CH3COOH}/\ce{CH3COO-} (échelle pH 0--14).}
\end{popfigure}
\textbf{Conclusion :} à pH $= 3,0$, l'acide éthanoïque est présent quasi exclusivement sous sa forme acide \ce{CH3COOH}.
\end{exoresolu}

\courssub{Déterminer expérimentalement un Ka et le degré d'acidité d'un vinaigre}

\begin{methode}[Exploiter un dosage pH-métrique d'un acide faible]
\begin{enumerate}
\item \textbf{Repérer la demi-équivalence} : à $V = \dfrac{V_E}{2}$, la moitié de l'acide a été transformée en base conjuguée : $[\ce{AH}] = [\ce{A-}]$, donc $\mathrm{pH} = \mathrm{p}K_a$. Lecture directe du $\mathrm{p}K_a$ sur la courbe $\mathrm{pH} = f(V_b)$.
\item \textbf{Déterminer $V_E$} par la méthode des tangentes (ou du maximum de la courbe dérivée $\Delta\mathrm{pH}/\Delta V$).
\item \textbf{Calculer la concentration} de l'acide : à l'équivalence, $C_A V_A = C_B V_E$, d'où $C_A = \dfrac{C_B V_E}{V_A}$.
\item \textbf{Degré d'acidité d'un vinaigre} : c'est la masse d'acide éthanoïque (en g) contenue dans $100$ g de vinaigre (assimilé à $100$ mL, $\rho \approx 1$) :
\[ d = \frac{C_A \times M(\ce{CH3COOH})}{10} \quad \text{avec } M(\ce{CH3COOH}) = \SI{60}{g.mol^{-1}}, \]
$C_A$ étant la concentration du vinaigre commercial (penser au facteur de dilution si le vinaigre a été dilué avant dosage !).
\end{enumerate}
\end{methode}

\begin{exoresolu}[Degré d'acidité d'un vinaigre du marché de Mokolo]
Un vinaigre commercial est dilué $10$ fois (solution $S$). On dose $V_A = \SI{20,0}{mL}$ de $S$ par de la soude à $C_B = \SI{0,100}{mol.L^{-1}}$. L'équivalence est atteinte pour $V_E = \SI{15,0}{mL}$ et le pH à la demi-équivalence vaut $4,8$. Données : $M(\ce{CH3COOH}) = \SI{60}{g.mol^{-1}}$.
\begin{enumerate}
\item Déterminer le $\mathrm{p}K_a$ du couple de l'acide éthanoïque.
\item Calculer la concentration $C_A$ de $S$, puis celle du vinaigre commercial.
\item En déduire le degré d'acidité du vinaigre.
\end{enumerate}
\tcblower
\textbf{1. Détermination du $\mathrm{p}K_a$.} À la demi-équivalence, $[\ce{CH3COOH}] = [\ce{CH3COO-}]$, donc :
\[ \mathrm{pH}\left(\tfrac{V_E}{2}\right) = \mathrm{p}K_a = 4,8 \quad \text{soit} \quad K_a = 10^{-4,8} = 1,6\times 10^{-5}. \]
\textbf{2. Concentrations.} À l'équivalence du dosage \ce{CH3COOH + HO- -> CH3COO- + H2O} :
\[ C_A V_A = C_B V_E \implies C_A = \frac{C_B V_E}{V_A} = \frac{0,100 \times 15,0}{20,0} = \SI{7,5e-2}{mol.L^{-1}}. \]
Le vinaigre commercial est $10$ fois plus concentré (dilution par 10) :
\[ C_0 = 10 \times C_A = \SI{0,75}{mol.L^{-1}}. \]
\textbf{3. Degré d'acidité.} Concentration massique du vinaigre :
\[ C_m = C_0 \times M = 0,75 \times 60 = \SI{45}{g.L^{-1}}. \]
Le degré est la masse d'acide dans $100$ mL de vinaigre :
\[ d = \frac{C_m}{10} = \frac{45}{10} = 4,5. \]
\textbf{Conclusion :} le vinaigre a un degré d'acidité de $4,5^{\circ}$ (soit $4,5$ g d'acide éthanoïque pour $100$ g de vinaigre), et le couple étudié a pour $\mathrm{p}K_a = 4,8$, valeur conforme aux tables.
\end{exoresolu}

\begin{attention}[Les 5 erreurs qui coûtent des points au Bac]
\begin{enumerate}
\item \textbf{Confondre les formules de pH.} $\mathrm{pH} = -\log C$ n'est valable que pour un acide \textbf{fort} ; $\mathrm{pH} = 14 + \log C$ que pour une base \textbf{forte}. Appliquer ces formules à un acide faible (vinaigre, jus de bissap) est une erreur rédhibitoire : pour un acide faible, $\mathrm{pH} > -\log C$.
\item \textbf{Oublier le facteur de dilution.} Dans le dosage du vinaigre, la solution dosée est diluée $10$ fois : la concentration du vinaigre commercial est $C_0 = 10\,C_A$. Oublier ce facteur donne un degré d'acidité $10$ fois trop faible.
\item \textbf{Écrire une flèche simple pour un équilibre.} La réaction d'un acide faible avec l'eau s'écrit obligatoirement avec $\ce{<=>}$ ; réserver $\ce{->}$ aux acides et bases forts et aux réactions quasi totales ($K > 10^{4}$).
\item \textbf{Inverser le sens de la classification.} Plus le $\mathrm{p}K_a$ est \textbf{petit}, plus l'acide est \textbf{fort} (et plus sa base conjuguée est faible). Beaucoup d'élèves affirment le contraire : vérifier avec un exemple connu (\ce{HCl} fort, $\mathrm{p}K_a$ très petit).
\item \textbf{Mal lire la demi-équivalence.} Le $\mathrm{p}K_a$ se lit au pH pour $V = V_E/2$, \textbf{pas} à l'équivalence (où le pH vaut souvent $8$ à $9$ pour un acide faible dosé par une base forte). Bien distinguer les deux points remarquables de la courbe.
\end{enumerate}
\end{attention}

\newpage

\coursec{Exercices}

\coursec{Force des acides et des bases, couples acide/base}

\courssub{Introduction : une question de force}
Dans l'industrie agroalimentaire camerounaise, le contrôle de l'acidité est crucial : le \emph{vin de palme} fermente pour donner le \emph{vinaigre de palme} (acide éthanoïque), le \emph{bili-bili} (bière de mil) a un pH acide, les savonneries artisanales de l'Ouest utilisent la soude (hydroxyde de sodium) pour la saponification de l'huile de palme. Pourquoi l'acide chlorhydrique « brûle » plus que le vinaigre à même concentration ? Pourquoi l'ammoniaque ménager nettoie-t-il sans être aussi corrosif que la soude ? La réponse réside dans la \cle{force} des acides et des bases, c'est-à-dire leur capacité à céder ou capter un proton \ce{H+} (ion oxonium \ce{H3O+} en milieu aqueux). Cette leçon pose les bases quantitatives pour prévoir et piloter les réactions acido-basiques.

\courssub{1. Acides forts et bases fortes : une ionisation totale}

\begin{savaistu}[Contexte industriel : la soude et l'acide chlorhydrique]
À la \emph{SONARA} (Société Nationale de Raffinage), l'acide chlorhydrique est utilisé pour le décapage des métaux. Dans les savonneries artisanales de Bafoussam ou Foumban, l'hydroxyde de sodium (soude caustique) transforme l'huile de palme en savon. Ces deux réactifs sont des « forces » de la chimie : ils réagissent totalement avec l'eau.
\end{savaistu}

\begin{definition}[Acide fort et base forte]
Un \textbf{acide fort} est un acide qui réagit \textbf{totalement} avec l'eau : il cède la totalité de ses protons. Son équation s'écrit avec une flèche simple \ce{->}.
\[ \ce{AH + H2O -> A- + H3O+} \]
Un \textbf{base forte} est une base qui réagit \textbf{totalement} avec l'eau : elle capte totalement les protons de l'eau (ou libère totalement des ions \ce{HO-}).
\[ \ce{B + H2O -> BH+ + HO-} \quad \text{ou} \quad \ce{HO- -> HO-} \text{ (sel soluble)} \]
\end{definition}

\begin{propriete}[pH des solutions d'acide fort et de base forte]
À \SI{25}{\celsius}, pour une solution d'acide fort de concentration molaire $C$ (en \SI{}{mol.L^{-1}}) :
\[ [\ce{H3O+}] = C \quad \Rightarrow \quad \mathrm{pH} = -\log C \]
Pour une solution de base forte de concentration $C$ (donnant $n$ ions \ce{HO-} par formule) :
\[ [\ce{HO-}] = nC \quad \Rightarrow \quad \mathrm{pOH} = -\log(nC) \quad \Rightarrow \quad \mathrm{pH} = 14 + \log(nC) \]
\end{propriete}

\begin{methode}[Vérifier qu'un acide (ou une base) est fort]
\begin{enumerate}
\item Mesurer le pH de la solution de concentration connue $C$.
\item Calculer le pH théorique attendu pour un acide fort : $\mathrm{pH}_{\text{th}} = -\log C$ (ou pour une base forte : $\mathrm{pH}_{\text{th}} = 14 + \log C$).
\item Comparer : si $\mathrm{pH}_{\text{mesuré}} \approx \mathrm{pH}_{\text{th}}$, l'espèce est forte (ionisation totale).
\end{enumerate}
\end{methode}

\begin{exemplebox}[Acide chlorhydrique et hydroxyde de sodium]
\begin{itemize}
\item Solution d'\ce{HCl} $C = \SI{1,0e-3}{mol.L^{-1}}$ : $\mathrm{pH}_{\text{th}} = -\log(10^{-3}) = 3,0$. Mesure = 3,0 $\Rightarrow$ \ce{HCl} acide fort.
\item Solution de \ce{NaOH} $C = \SI{5,0e-3}{mol.L^{-1}}$ : $[\ce{HO-}] = 5,0\times10^{-3}$, $\mathrm{pOH} = 2,30$, $\mathrm{pH}_{\text{th}} = 11,70$. Mesure = 11,7 $\Rightarrow$ \ce{NaOH} base forte.
\end{itemize}
\end{exemplebox}

\begin{attention}[Piège : concentration très faible]
Pour $C < \SI{1e-6}{mol.L^{-1}}$, l'autoprotolyse de l'eau ($K_e = 10^{-14}$) perturbe le calcul : $[\ce{H3O+}] \approx C + 10^{-7}$. Le pH ne descend pas en dessous de 7 pour un acide, ni au-dessus de 7 pour une base.
\end{attention}

\courssub{2. Acides faibles et bases faibles : un équilibre limité}

\begin{savaistu}[Le vinaigre de palme et l'ammoniaque ménager]
Le vinaigre artisanal (issu du vin de palme) contient de l'\textbf{acide éthanoïque} \ce{CH3COOH}. L'ammoniaque ménager contient de l'\textbf{ammoniac} \ce{NH3}. Contrairement à l'acide chlorhydrique ou à la soude, ils ne réagissent pas totalement avec l'eau : on observe un \textbf{équilibre chimique}.
\end{savaistu}

\begin{definition}[Acide faible et base faible]
Un \textbf{acide faible} (noté \ce{AH}) réagit \textbf{partiellement} avec l'eau. La réaction est réversible, on utilise la flèche d'équilibre \ce{<=>} :
\[ \ce{AH + H2O <=> A- + H3O+} \]
Le couple acide/base est \ce{AH}/\ce{A-}. \ce{AH} est la forme acide, \ce{A-} la forme basique (base conjuguée).
Une \textbf{base faible} (notée \ce{B}) réagit partiellement avec l'eau :
\[ \ce{B + H2O <=> BH+ + HO-} \]
Le couple est \ce{BH+}/\ce{B}. \ce{BH+} est la forme acide (acide conjugué), \ce{B} la forme basique.
\end{definition}

\begin{propriete}[Coefficient d'ionisation (ou taux d'avancement final) $\alpha = \tau$]
Pour un acide faible \ce{AH} de concentration initiale $C$, d'avancement final $x_f$ :
\[ \alpha = \tau = \frac{x_f}{C} = \frac{[\ce{A-}]_f}{C} = \frac{[\ce{H3O+}]_f}{C} \quad (\text{si on néglige l'autoprotolyse de l'eau}) \]
Pour un acide faible, $\alpha < 1$ (souvent $\alpha < 5\%$). Plus $\alpha$ est petit, plus l'acide est faible.
\end{propriete}

\begin{methode}[Montrer qu'un acide est faible à partir du pH]
\begin{enumerate}
\item Calculer $-\log C$ (pH théorique si l'acide était fort).
\item Mesurer le pH réel.
\item Si $\mathrm{pH}_{\text{mesuré}} > -\log C$ (pour un acide), alors $[\ce{H3O+}] < C$ : l'ionisation est partielle, l'acide est \textbf{faible}.
\item Pour une base : si $\mathrm{pH}_{\text{mesuré}} < 14 + \log C$, la base est \textbf{faible}.
\end{enumerate}
\end{methode}

\begin{exemplebox}[Acide éthanoïque $C = \SI{1,0e-2}{mol.L^{-1}}$]
$-\log C = 2,0$. pH mesuré = 3,4.
$3,4 > 2,0 \Rightarrow [\ce{H3O+}] = 10^{-3,4} = \SI{4,0e-4}{mol.L^{-1}} < C$.
L'acide éthanoïque est un \textbf{acide faible}. $\alpha = 4,0\times10^{-4} / 10^{-2} = 4\%$.
\end{exemplebox}

\courssub{3. Étude quantitative d'un équilibre : avancement et tableau}

\begin{definition}[Avancement d'une réaction ($x$)]
L'avancement $x$ (en \SI{}{mol}) quantifie l'évolution d'une réaction. Pour la réaction :
\[ \ce{AH + H2O <=> A- + H3O+} \]
Si $x$ moles par litre ont réagi : $[\ce{AH}] = C - x$, $[\ce{A-}] = x$, $[\ce{H3O+}] = x$ (eau en excès, concentration constante).
L'avancement final à l'équilibre est noté $x_f$. Le taux d'avancement final est $\tau = x_f / C_{\text{lim}}$.
\end{definition}

\begin{methode}[Construction du tableau d'avancement (1 L de solution)]
\begin{center}
\begin{tabularx}{\linewidth}{|l|X|X|X|}
\hline
\rowcolor{popblueL} \textbf{Espèce} & \ce{AH} & \ce{A-} & \ce{H3O+} \\
\hline
\textbf{État initial} & $C$ & $0$ & $0$ (négligé) \\
\hline
\textbf{État final} & $C - x_f$ & $x_f$ & $x_f$ \\
\hline
\end{tabularx}
\end{center}
On utilise les données expérimentales (pH $\Rightarrow$ $[\ce{H3O+}]_f = x_f$) pour remplir la ligne finale.
\end{methode}

\courssub{4. Constante d'acidité $K_a$, constante de basicité $K_b$ et $pK_a$}

\begin{propriete}[Constante d'acidité $K_a$]
Pour le couple \ce{AH}/\ce{A-} à l'équilibre :
\[ K_a = \frac{[\ce{A-}][\ce{H3O+}]}{[\ce{AH}]} \]
$K_a$ est sans dimension (activités). Elle ne dépend que de la température (valeurs usuelles à \SI{25}{\celsius}).
Plus $K_a$ est grand, plus l'équilibre est déplacé vers la droite, plus l'acide \ce{AH} est \textbf{fort}.
\end{propriete}

\begin{propriete}[Constante de basicité $K_b$ et relation $K_a \times K_b = K_e$]
Pour la base conjuguée \ce{A-} : $\ce{A- + H2O <=> AH + HO-}$.
\[ K_b = \frac{[\ce{AH}][\ce{HO-}]}{[\ce{A-}]} \]
En multipliant $K_a$ et $K_b$ : $K_a \times K_b = [\ce{H3O+}][\ce{HO-}] = K_e = 10^{-14}$ (à \SI{25}{\celsius}).
En échelle logarithmique : $\mathbf{pK_a + pK_b = 14}$ (avec $pK = -\log K$).
\end{propriete}

\begin{aretenir}[Échelle de force : $pK_a$]
\begin{itemize}
\item $pK_a = -\log K_a$. Plus $pK_a$ est \textbf{petit}, plus $K_a$ est grand, plus l'acide est \textbf{fort}.
\item \textbf{Acides forts} : $pK_a < 0$ (ex. \ce{HCl}, $pK_a \approx -7$).
\item \textbf{Acides faibles} : $0 < pK_a < 14$ (ex. \ce{CH3COOH} 4,8 ; \ce{HF} 3,2 ; \ce{NH4+} 9,2).
\item \textbf{Bases fortes} : base conjuguée d'acide très faible ($pK_a > 14$, ex. \ce{HO-}, $pK_a(\ce{H2O}) = 14$).
\item \textbf{Bases faibles} : base conjuguée d'acide faible ($pK_a$ de l'acide conjugué entre 0 et 14).
\end{itemize}
\end{aretenir}

\begin{attention}[Notation $pK_a(\ce{H2O}/\ce{HO-})$]
Certains manuels donnent $pK_a(\ce{H2O}/\ce{HO-}) = 14$. Rigoureusement, c'est le $pK_e$ (produit ionique de l'eau). Pour le couple \ce{H2O}/\ce{HO-}, l'acide est \ce{H2O}, sa constante d'acidité est $K_a = K_e / [\ce{H2O}] \approx 10^{-15,7}$ ($pK_a \approx 15,7$). Dans le cadre du programme TI, on utilise la convention $pK_a(\ce{H2O}/\ce{HO-}) = 14$ pour calculer la constante d'équilibre d'une réaction de dosage : $K = K_a(\text{acide dosé}) / K_e$.
\end{attention}

\courssub{5. Classification des couples et prévision des réactions (Règle du $\gamma$)}

\begin{propriete}[Règle du $\gamma$ (prévoir le sens d'une réaction acido-basique)]
On met en présence un acide \ce{AH_1} et une base \ce{B_2} (forme basique du couple 2). La réaction possible est :
\[ \ce{AH_1 + B_2 <=> A1- + BH2+} \]
On compare les $pK_a$ des deux acides mis en jeu : \ce{AH_1} ($pK_{a1}$) et \ce{BH2+} ($pK_{a2}$).
\begin{itemize}
\item Si $pK_{a1} < pK_{a2}$ : l'acide \ce{AH_1} est \textbf{plus fort} que \ce{BH2+}. L'équilibre se fait \textbf{vers la droite} (réaction quasi totale si $\Delta pK_a > 4$).
\item Si $pK_{a1} > pK_{a2}$ : l'équilibre se fait \textbf{vers la gauche}.
\item Si $pK_{a1} = pK_{a2}$ : équilibre au milieu.
\end{itemize}
La constante d'équilibre de cette réaction est $K = 10^{pK_{a2} - pK_{a1}} = \frac{K_{a1}}{K_{a2}}$.
\end{propriete}

\begin{methode}[Appliquer la règle du $\gamma$]
\begin{enumerate}
\item Identifier les deux couples acide/base présents.
\item Écrire l'équation de la réaction entre l'acide du couple 1 et la base du couple 2.
\item Noter $pK_{a1}$ (acide réactif) et $pK_{a2}$ (acide conjugué de la base réactive).
\item Comparer : si $pK_{a1} < pK_{a2}$, la réaction est favorable (vers la droite).
\item Calculer $K = 10^{pK_{a2} - pK_{a1}}$. Si $K > 10^4$, réaction quasi totale.
\end{enumerate}
\end{methode}

\begin{exemplebox}[Réaction \ce{HF} + \ce{NH3}]
Couples : \ce{HF}/\ce{F-} ($pK_a = 3,2$) et \ce{NH4+}/\ce{NH3} ($pK_a = 9,2$).
Réaction : $\ce{HF + NH3 <=> F- + NH4+}$.
$pK_{a1}(\ce{HF}) = 3,2 < pK_{a2}(\ce{NH4+}) = 9,2$.
$\Rightarrow$ Réaction favorable vers la droite. $K = 10^{9,2 - 3,2} = 10^6$. Quasi totale.
\end{exemplebox}

\begin{popfigure}
\begin{tikzpicture}[scale=0.9, every node/.style={font=\small}]
\draw[popline, thick, ->] (0,0) -- (14,0) node[right] {$pK_a$};
\foreach \x/\label in {0/0, 3.2/3,2, 4.8/4,8, 7/7, 9.2/9,2, 14/14} {
    \draw (\x,0.1) -- (\x,-0.1) node[below] {\label};
}
% Couples
\draw[popblue, thick] (3.2,0.5) -- (3.2,1.5) node[above, popblue] {\ce{HF}/\ce{F-}};
\draw[poporange, thick] (4.8,0.5) -- (4.8,1.5) node[above, poporange] {\ce{CH3COOH}/\ce{CH3COO-}};
\draw[popgreen, thick] (9.2,0.5) -- (9.2,1.5) node[above, popgreen] {\ce{NH4+}/\ce{NH3}};
% Zones
\node[popblue!30, draw=popblue, rounded corners, fit={(0,0.5) (3.2,1.5)}] {};
\node[poporange!30, draw=poporange, rounded corners, fit={(3.2,0.5) (4.8,1.5)}] {};
\node[popgreen!30, draw=popgreen, rounded corners, fit={(4.8,0.5) (9.2,1.5)}] {};
\node[poppurple!30, draw=poppurple, rounded corners, fit={(9.2,0.5) (14,1.5)}] {};
% Force
\node[below=0.5cm] at (1.6,-0.5) {\textbf{Acides forts} ($pK_a < 0$)};
\node[below=0.5cm] at (7,-0.5) {\textbf{Acides faibles}};
\node[below=0.5cm] at (11.6,-0.5) {\textbf{Acides très faibles}};
\node[above=2.2cm, align=center] at (7,0) {Force acide \textbf{croissante} $\leftarrow$ \quad \quad Force basique \textbf{croissante} $\rightarrow$};
\end{tikzpicture}
\legende{Axe des $pK_a$ : classification des couples acide/base. Plus on va à gauche, plus l'acide est fort ; plus on va à droite, plus la base conjuguée est forte.}
\end{popfigure}

\courssub{6. Relation $\mathrm{pH}$--$pK_a$ et domaines de prédominance}

\begin{propriete}[Équation de Henderson-Hasselbalch]
À partir de $K_a = \frac{[\ce{A-}][\ce{H3O+}]}{[\ce{AH}]}$, on prend le $-\log$ :
\[ \mathrm{pH} = pK_a + \log \frac{[\ce{A-}]}{[\ce{AH}]} \]
Cette relation lie le pH de la solution au rapport des concentrations des formes basique et acide du couple.
\end{propriete}

\begin{aretenir}[Domaines de prédominance]
\begin{itemize}
\item Si $\mathrm{pH} < pK_a$ : $\log \frac{[\ce{A-}]}{[\ce{AH}]} < 0 \Rightarrow [\ce{A-}] < [\ce{AH}]$. \textbf{La forme acide \ce{AH} prédomine}.
\item Si $\mathrm{pH} > pK_a$ : $\log \frac{[\ce{A-}]}{[\ce{AH}]} > 0 \Rightarrow [\ce{A-}] > [\ce{AH}]$. \textbf{La forme basique \ce{A-} prédomine}.
\item Si $\mathrm{pH} = pK_a$ : $[\ce{A-}] = [\ce{AH}]$. Les deux formes sont en \textbf{concentrations égales} (demi-équivalence).
\end{itemize}
\end{aretenir}

\begin{popfigure}
\begin{tikzpicture}[scale=1, every node/.style={font=\small}]
\begin{axis}[
    width=\linewidth, height=6cm,
    axis lines=middle,
    xlabel={pH}, ylabel={Fraction molaire},
    xmin=0, xmax=14, ymin=0, ymax=1.1,
    xtick={0,4.8,7,9.2,14}, xticklabels={0, $pK_a$, 7, $pK_a'$, 14},
    ytick={0,0.5,1},
    legend style={at={(0.5,1.15)}, anchor=north, legend columns=2},
    domain=0:14, samples=200,
]
% pKa = 4.8 example
\addplot[popblue, thick] {1/(1+10^(4.8-x))}; \addlegendentry{$\alpha_{\ce{A-}}$ (Forme basique)};
\addplot[poporange, thick] {1/(1+10^(x-4.8))}; \addlegendentry{$\alpha_{\ce{AH}}$ (Forme acide)};
\draw[dashed, popline] (4.8,0) -- (4.8,1);
\draw[dashed, popline] (7,0) -- (7,1);
\end{axis}
\end{tikzpicture}
\legende{Diagramme de prédominance pour un couple de $pK_a = 4,8$ (ex. acide éthanoïque). Au pH = $pK_a$, les deux formes sont à 50\%.}
\end{popfigure}

\courssub{7. Applications : détermination expérimentale de $K_a$ et degré d'acidité du vinaigre}

\begin{experience}[Détermination de $K_a$ d'un acide faible par pH-métrie]
\textbf{Objectif :} Déterminer $K_a$ d'un acide faible \ce{AH} (ex. acide méthanoïque du bissap, acide éthanoïque du vinaigre).
\textbf{Matériel :} pH-mètre étalonné (étalons pH 4, 7, 10), bécher, solution mère d'acide \ce{AH} de concentration connue $C$, eau distillée, agitation magnétique.
\textbf{Protocole :}
\begin{enumerate}
\item Préparer une solution de concentration $C$ connue (ex. \SI{5,0e-2}{mol.L^{-1}}).
\item Mesurer le pH à \SI{25}{\celsius} après stabilisation.
\item Calculer $[\ce{H3O+}] = 10^{-\mathrm{pH}}$.
\item Dresser le tableau d'avancement, en déduire $[\ce{A-}] = [\ce{H3O+}]$ et $[\ce{AH}] = C - [\ce{H3O+}]$.
\item Calculer $K_a = \frac{[\ce{H3O+}]^2}{C - [\ce{H3O+}]}$ et $pK_a = -\log K_a$.
\end{enumerate}
\textbf{Interprétation :} Comparer le $pK_a$ trouvé aux tables de données pour identifier l'acide.
\end{experience}

\begin{methode}[Calcul du degré d'acidité d'un vinaigre (dosage pH-métrique)]
Le degré d'acidité = masse d'acide éthanoïque (en g) pour \SI{100}{g} de vinaigre.
\begin{enumerate}
\item \textbf{Dilution} : Diluer le vinaigre (trop concentré) $f$ fois (ex. $f=10$) $\rightarrow$ solution $S$ ($C_A$).
\item \textbf{Dosage} : Dosage de $V_A$ de $S$ par \ce{NaOH} $C_B$ (pH-métrique). Relever $V_E$ à l'équivalence (point d'inflexion de la courbe $\mathrm{pH} = f(V_B)$).
\item \textbf{Calculs} :
    \begin{align*}
    n_{\ce{CH3COOH}} &= n_{\ce{HO-}} = C_B \times V_E \\
    C_A &= \frac{n_{\ce{CH3COOH}}}{V_A} \\
    C_0 (\text{vinaigre pur}) &= f \times C_A \\
    \text{Masse volumique } \rho &\approx \SI{1,0}{g.mL^{-1}} \Rightarrow \text{masse de 1 L} = \SI{1000}{g} \\
    \text{Masse acide dans 1 L} &= C_0 \times M_{\ce{CH3COOH}} \\
    \text{Degré} &= \frac{\text{Masse acide (g/L)}}{10} \quad (\text{car } \SI{1000}{g} \rightarrow \SI{100}{g} \text{ divise par 10})
    \end{align*}
\item \textbf{Choix indicateur} : pH à l'équivalence > 7 (hydrolyse \ce{CH3COO-}) $\Rightarrow$ indicateur à transition basique (phénolphtaléine 8,2--10,0).
\end{enumerate}
\end{methode}

\begin{attention}[Erreurs fréquentes au Bac]
\begin{itemize}
\item Confondre $C_A$ (solution diluée) et $C_0$ (vinaigre pur) : ne pas oublier le facteur de dilution $f$.
\item Oublier que le degré est en \textbf{g/100g}, pas g/L. Avec $\rho=1$, degré = $C_0 \times M / 10$.
\item Indicateur : BBT (6,0--7,6) change \textbf{avant} l'équivalence (erreur systématique négative). Phénolphtaléine (8,2--10,0) est le bon choix.
\item Demi-équivalence : $\mathrm{pH} = pK_a$ \textbf{uniquement} si on néglige la dilution (ou si on utilise les quantités de matière : $n_{\ce{AH}} = n_{\ce{A-}}$).
\end{itemize}
\end{attention}

% ============================================================
% EXERCICES
% ============================================================

\courssub{Je vérifie mes connaissances}

\begin{exercice}[QCM : force des acides et des bases]
Pour chaque question, une seule réponse est exacte. Recopie la lettre correspondante.
\begin{enumerate}
\item Une solution d'acide chlorhydrique de concentration $C = \SI{1,0e-3}{mol.L^{-1}}$ a un pH égal à :
\begin{enumerate}
\item[a)] $1{,}0$ \quad b) $2{,}0$ \quad c) $3{,}0$ \quad d) $11{,}0$
\end{enumerate}
\item Une solution d'acide éthanoïque de concentration $C = \SI{1,0e-2}{mol.L^{-1}}$ a un pH de $3{,}4$. On peut affirmer que :
\begin{enumerate}
\item[a)] l'acide est fort car le pH est acide ;
\item[b)] l'acide est faible car $\mathrm{pH} > -\log C$ ;
\item[c)] l'acide est faible car $\mathrm{pH} < -\log C$ ;
\item[d)] l'acide est totalement ionisé.
\end{enumerate}
\item Le couple \ce{NH4+}/\ce{NH3} a un $pK_a = 9{,}2$. Dans une solution de $\mathrm{pH} = 11{,}0$, l'espèce prédominante est :
\begin{enumerate}
\item[a)] \ce{NH4+} \quad b) \ce{NH3} \quad c) les deux en quantités égales \quad d) \ce{H3O+}
\end{enumerate}
\item Plus le $K_a$ d'un couple acide/base est grand :
\begin{enumerate}
\item[a)] plus l'acide est faible ;
\item[b)] plus l'acide est fort ;
\item[c)] plus la base conjuguée est forte ;
\item[d)] plus le $pK_a$ est grand.
\end{enumerate}
\end{enumerate}
\end{exercice}

\begin{exercice}[Vrai ou faux, en justifiant]
Pour chaque affirmation, réponds par \textbf{vrai} ou \textbf{faux} et justifie ta réponse par un calcul ou un argument.
\begin{enumerate}
\item Une solution d'hydroxyde de sodium de concentration $C = \SI{5,0e-3}{mol.L^{-1}}$ a un pH de $11{,}7$ à \SI{25}{\celsius}.
\item Le taux d'avancement final de la réaction d'un acide fort avec l'eau est égal à $1$.
\item Deux acides faibles de même concentration ont le même pH.
\item La constante d'acidité du couple \ce{CH3COOH}/\ce{CH3COO-} s'écrit $K_a = \dfrac{[\ce{CH3COO-}][\ce{H3O+}]}{[\ce{CH3COOH}]}$.
\item À la demi-équivalence du dosage d'un acide faible par la soude, le pH est égal au $pK_a$ du couple.
\end{enumerate}
\end{exercice}

\begin{exercice}[Texte à trous]
Recopie et complète le texte suivant avec les mots ou expressions de la liste : \emph{totale ; partielle ; taux d'avancement ; $pK_a$ ; base conjuguée ; prédomine ; ionisation ; $-\log C$.}

Un acide fort réagit avec l'eau de manière \ldots\ldots ; le pH de sa solution de concentration $C$ vaut alors \ldots\ldots. Un acide faible ne s'ionise que de façon \ldots\ldots : on caractérise cette réaction limitée par son \ldots\ldots, inférieur à $1$. La force d'un couple acide/base est repérée par la valeur de son \ldots\ldots : plus celui-ci est petit, plus l'acide est fort et plus sa \ldots\ldots est faible. Dans une solution dont le pH est supérieur au $pK_a$ d'un couple, c'est la forme basique qui \ldots\ldots.
\end{exercice}

\begin{exercice}[Définitions et couples]
\begin{enumerate}
\item Définis : acide fort, base faible, constante d'acidité $K_a$, coefficient d'ionisation $\alpha$.
\item Identifie les deux couples acide/base mis en jeu dans l'équilibre suivant, en précisant pour chacun la forme acide et la forme basique :
\[ \ce{HCOOH + NH3 <=> HCOO- + NH4+} \]
\item Écris l'équation de la réaction de l'ion ammonium \ce{NH4+} avec l'eau et donne l'expression de la constante d'acidité du couple correspondant.
\end{enumerate}
\end{exercice}

\courssub{Je m'entraîne}

\begin{exercice}[Classification de trois solutions]
On mesure, à \SI{25}{\celsius}, le pH de trois solutions aqueuses de même concentration $C = \SI{1,0e-2}{mol.L^{-1}}$ :
\begin{center}
\begin{tabularx}{\linewidth}{|l|X|X|X|}
\hline
\rowcolor{popblueL} \textbf{Solution} & \ce{HCl} & \ce{CH3COOH} & \ce{NH3} \\
\hline
\textbf{pH mesuré} & $2{,}0$ & $3{,}4$ & $10{,}6$ \\
\hline
\end{tabularx}
\end{center}
\begin{enumerate}
\item Calculer le pH théorique d'un acide fort et celui d'une base forte de concentration $C$.
\item Comparer, pour chaque solution, le pH mesuré à la valeur théorique correspondante.
\item Classer chacune des trois espèces en acide fort, acide faible ou base faible, en justifiant.
\end{enumerate}
\end{exercice}

\begin{exercice}[Coefficient d'ionisation de l'acide méthanoïque]
Une solution d'acide méthanoïque \ce{HCOOH} de concentration $C = \SI{5,0e-2}{mol.L^{-1}}$ a un pH de $2{,}6$ à \SI{25}{\celsius}.
\begin{enumerate}
\item Écrire l'équation de la réaction de l'acide méthanoïque avec l'eau.
\item Dresser le tableau d'avancement de cette réaction pour un litre de solution.
\item Calculer la concentration en ions oxonium $[\ce{H3O+}]$, puis l'avancement final $x_f$.
\item En déduire le taux d'avancement final $\tau$ et le coefficient d'ionisation $\alpha$ de l'acide.
\item Conclure sur la force de l'acide méthanoïque.
\end{enumerate}
\end{exercice}

\begin{exercice}[Classement de couples et prévision d'une réaction]
On donne les $pK_a$ de trois couples à \SI{25}{\celsius} : \ce{HF}/\ce{F-} : $3{,}2$ ; \ce{CH3COOH}/\ce{CH3COO-} : $4{,}8$ ; \ce{NH4+}/\ce{NH3} : $9{,}2$.
\begin{enumerate}
\item Classer les trois acides par force croissante, puis les trois bases conjuguées par force croissante. Justifier.
\item Placer ces couples sur un axe de $pK_a$ gradué de $0$ à $14$.
\item Prévoir, à l'aide de la règle du gamma, si la réaction entre \ce{HF} et \ce{NH3} est possible.
\item Écrire l'équation de cette réaction et calculer sa constante d'équilibre $K$. La réaction est-elle quasi totale ?
\end{enumerate}
\end{exercice}

\begin{exercice}[Relation pH--$pK_a$ et ammoniaque ménager]
\begin{enumerate}
\item À partir de l'expression de la constante d'acidité d'un couple \ce{AH}/\ce{A-}, établir la relation :
\[ \mathrm{pH} = pK_a + \log \dfrac{[\ce{A-}]}{[\ce{AH}]} \]
\item L'ammoniaque ménager, utilisée pour le nettoyage, contient de l'ammoniac \ce{NH3} en solution. Son pH vaut $11{,}0$. On donne $pK_a(\ce{NH4+}/\ce{NH3}) = 9{,}2$.
\begin{enumerate}
\item Déterminer l'espèce prédominante du couple dans cette solution.
\item Calculer la valeur du rapport $\dfrac{[\ce{NH3}]}{[\ce{NH4+}]}$.
\end{enumerate}
\item À quel pH les deux espèces du couple seraient-elles présentes en concentrations égales ?
\end{enumerate}
\end{exercice}

\courssub{Je me prépare au Bac}

\begin{exercice}[Type Bac — Degré d'acidité d'un vinaigre de vin de palme]
Un élève de Terminale TI souhaite vérifier le degré d'acidité d'un vinaigre artisanal obtenu par fermentation du vin de palme. Le vinaigre commercial est trop concentré : il prépare d'abord une solution $S$ en diluant $10$ fois le vinaigre. Il dose ensuite un volume $V_A = \SI{10,0}{mL}$ de la solution $S$ par une solution d'hydroxyde de sodium de concentration $C_B = \SI{0,10}{mol.L^{-1}}$, en suivant l'évolution du pH. L'équivalence est obtenue pour un volume versé $V_E = \SI{12,0}{mL}$.

\textbf{Données :} $pK_a(\ce{CH3COOH}/\ce{CH3COO-}) = 4{,}8$ ; masse molaire de l'acide éthanoïque $M = \SI{60}{g.mol^{-1}}$ ; le degré d'acidité d'un vinaigre est la masse d'acide éthanoïque, en grammes, contenue dans $\SI{100}{g}$ de vinaigre ; on prendra la masse volumique du vinaigre égale à $\rho = \SI{1,0}{g.mL^{-1}}$.
\begin{enumerate}
\item Écrire l'équation de la réaction de dosage et calculer sa constante d'équilibre (on donne $pK_e = 14$). Conclure sur son aptitude au dosage.
\item Faire le schéma légendé du dispositif de dosage pH-métrique.
\item Définir l'équivalence et calculer la concentration $C_A$ de la solution diluée $S$.
\item En déduire la concentration $C_0$ du vinaigre pur, puis son degré d'acidité. Commenter le résultat sachant qu'un vinaigre commercial affiche en général $6$ degrés.
\item Quelle est la valeur du pH à la demi-équivalence ? Justifier.
\item Pourquoi le pH à l'équivalence est-il supérieur à $7$ ? Nommer l'espèce responsable et préciser son caractère.
\item Choisir, en justifiant, un indicateur coloré adapté parmi : hélianthine ($3{,}1$--$4{,}4$), BBT ($6{,}0$--$7{,}6$), phénolphtaléine ($8{,}2$--$10{,}0$).
\end{enumerate}
\end{exercice}

\begin{popfigure}
\begin{tikzpicture}[scale=0.7, every node/.style={font=\small}]
% Support
\draw[popline, thick] (0,0) -- (10,0);
% Burette
\draw[popblue, thick] (8,0) -- (8,8) -- (9,8) -- (9,0);
\draw[popblue!50] (8.5,7) rectangle (8.9,7.2); % Robinet
\node[popblue, above] at (8.5,8) {Burette \ce{NaOH}};
\draw[->, popblue] (8.5,1) -- (5.5,1); % Goutte
% Bécher
\draw[popdark, thick] (2,0) -- (2,5) -- (5,5) -- (5,0);
\draw[popmuted] (2.2,1) -- (4.8,1); % Niveau liquide
% Agitation
\draw[poporange] (3.5,0.5) circle (0.3); % Barreau
\draw[poporange, decorate, decoration={snake, amplitude=1pt}] (3.5,0.2) -- (3.5,-0.5); % Aimant
% pH-metre
\draw[poppurple, thick] (1,5) -- (1,7) -- (3,7) -- (3,5); % Sonde
\node[poppurple, left] at (0.5,6) {pH-mètre};
\node[poppurple, left] at (0.5,5.5) {Électrode};
% Légendes
\node[below] at (3.5,-1) {Solution $S$ ($V_A$) + barreau aimanté};
\node[right] at (9.5,4) {Volume $V_B$};
\node[above] at (3.5,5.5) {Mesure pH};
\end{tikzpicture}
\legende{Dispositif de dosage pH-métrique : la courbe $\mathrm{pH} = f(V_B)$ permet de repérer l'équivalence (point d'inflexion, pente maximale).}
\end{popfigure}

\begin{exercice}[Type Bac — Identification d'un acide faible par détermination de son $K_a$]
Au laboratoire d'une savonnerie artisanale de la région de l'Ouest, on dispose d'une solution aqueuse d'un acide faible noté \ce{AH}, de concentration $C = \SI{1,0e-1}{mol.L^{-1}}$. La mesure du pH donne $2{,}9$ à \SI{25}{\celsius}. On néglige les ions \ce{H3O+} provenant de l'autoprotolyse de l'eau.

\textbf{Table de données :} $pK_a$ de quelques couples : acide méthanoïque \ce{HCOOH}/\ce{HCOO-} : $3{,}8$ ; acide éthanoïque \ce{CH3COOH}/\ce{CH3COO-} : $4{,}8$ ; acide fluorhydrique \ce{HF}/\ce{F-} : $3{,}2$ ; acide benzoïque \ce{C6H5COOH}/\ce{C6H5COO-} : $4{,}2$.
\begin{enumerate}
\item Écrire l'équation de la réaction de l'acide \ce{AH} avec l'eau et dresser le tableau d'avancement pour un litre de solution.
\item Montrer que cet acide est faible en comparant le pH mesuré à $-\log C$.
\item Calculer $[\ce{H3O+}]$, puis, en utilisant l'électroneutralité et la conservation de la matière, exprimer et calculer $[\ce{A-}]$ et $[\ce{AH}]$.
\item En déduire la constante d'acidité $K_a$ du couple, puis son $pK_a$.
\item Identifier l'acide \ce{AH} à l'aide de la table de données.
\item Calculer le coefficient d'ionisation $\alpha$ de cet acide dans la solution étudiée.
\item Comment évoluerait $\alpha$ si l'on diluait la solution ? Justifier qualitativement.
\end{enumerate}
\end{exercice}

\begin{exercice}[Type Bac — Étude comparative d'un acide fort et d'une base faible]
Pour préparer des solutions de nettoyage dans un atelier, un technicien dispose :
\begin{itemize}
\item d'une solution $S_1$ d'acide chlorhydrique de concentration $C_1 = \SI{2,0e-2}{mol.L^{-1}}$ ;
\item d'une solution $S_2$ d'ammoniac \ce{NH3} de concentration $C_2 = \SI{2,0e-2}{mol.L^{-1}}$.
\end{itemize}
Les mesures de pH à \SI{25}{\celsius} donnent : $\mathrm{pH}_1 = 1{,}7$ pour $S_1$ et $\mathrm{pH}_2 = 10{,}8$ pour $S_2$.

\textbf{Données :} $pK_a(\ce{NH4+}/\ce{NH3}) = 9{,}2$ ; $pK_e = 14$ à \SI{25}{\celsius}.
\begin{enumerate}
\item Montrer, par le calcul, que l'acide chlorhydrique est un acide fort. Écrire l'équation de sa réaction avec l'eau.
\item Calculer le pH qu'aurait une base forte de concentration $C_2$. En déduire que l'ammoniac est une base faible et écrire l'équation de sa réaction avec l'eau.
\item Calculer, pour la solution $S_2$, les concentrations $[\ce{HO-}]$, $[\ce{NH4+}]$ et $[\ce{NH3}]$.
\item En déduire le taux d'avancement final $\tau$ de la réaction de l'ammoniac avec l'eau.
\item Calculer la constante d'acidité $K_a$ du couple \ce{NH4+}/\ce{NH3} à partir des concentrations de la question 3, puis le $pK_a$. Comparer à la valeur donnée.
\item On mélange un volume $V_1 = \SI{20,0}{mL}$ de $S_1$ avec un volume $V_2 = \SI{20,0}{mL}$ de $S_2$.
\begin{enumerate}
\item Écrire l'équation de la réaction qui se produit et calculer sa constante d'équilibre.
\item Déterminer la composition du mélange à l'état final (quantités de matière des espèces du couple).
\item En déduire le pH du mélange obtenu.
\end{enumerate}
\end{enumerate}
\end{exercice}



\newpage

\coursec{Corrigés des exercices}

\coursec{Corrigés des exercices d'application — Force des acides et des bases, couples acide/base}

\begin{corrige}[Exercice 1 — QCM : force des acides et des bases]
\begin{enumerate}
\item \textbf{Réponse c)} $\mathrm{pH} = 3{,}0$.
\ce{HCl} est un \cle{acide fort} : ionisation totale, donc $[\ce{H3O+}] = C = \SI{1,0e-3}{mol.L^{-1}}$ et $\mathrm{pH} = -\log(1{,}0\times10^{-3}) = 3{,}0$. La réponse d) ($11{,}0$) correspondrait à une \cle{base forte} de même concentration.

\item \textbf{Réponse b)} l'acide est faible car $\mathrm{pH} > -\log C$.
Si l'acide était fort : $\mathrm{pH} = -\log(1{,}0\times10^{-2}) = 2{,}0$. Or le pH mesuré ($3{,}4$) est supérieur : $[\ce{H3O+}] = 10^{-3{,}4} = \SI{4,0e-4}{mol.L^{-1}} < C$. L'ionisation est partielle : l'acide éthanoïque est un \cle{acide faible}.

\item \textbf{Réponse b)} \ce{NH3}.
$\mathrm{pH} = 11{,}0 > pK_a = 9{,}2$ : on est dans le \cle{domaine de prédominance} de la forme basique \ce{NH3}.

\item \textbf{Réponse b)} plus l'acide est fort.
$K_a$ grand signifie que l'équilibre $\ce{AH + H2O <=> A- + H3O+}$ est déplacé vers la droite : l'acide est fort. Conséquences : $pK_a = -\log K_a$ est petit, et la \cle{base conjuguée} est faible (réponses c) et d) fausses).
\end{enumerate}
\end{corrige}

\begin{corrige}[Exercice 2 — Vrai ou faux]
\begin{enumerate}
\item \textbf{Vrai.} \ce{NaOH} est une \cle{base forte} : $[\ce{HO-}] = C = \SI{5,0e-3}{mol.L^{-1}}$.
$\mathrm{pOH} = -\log(5{,}0\times10^{-3}) = 2{,}30$ ; $\mathrm{pH} = 14 - 2{,}30 = 11{,}70$ à \SI{25}{\celsius}.

\item \textbf{Vrai.} Par définition, un acide fort réagit totalement avec l'eau : l'avancement final atteint l'avancement maximal, $x_f = x_{\max} = C \times V$, donc $\tau = \dfrac{x_f}{x_{\max}} = 1$ (soit $100\,\%$).

\item \textbf{Faux.} Le pH d'une solution d'acide faible dépend de $C$ \emph{et} de $K_a$ : pour $\alpha \ll 1$, $[\ce{H3O+}] \approx \sqrt{K_a \cdot C}$. Deux acides de $K_a$ différents donnent des pH différents. Exemple à $C = \SI{1,0e-2}{mol.L^{-1}}$ : \ce{HF} ($pK_a = 3{,}2$) donne $\mathrm{pH} \approx 2{,}6$, tandis que \ce{CH3COOH} ($pK_a = 4{,}8$) donne $\mathrm{pH} \approx 3{,}4$.

\item \textbf{Vrai.} Pour l'équilibre $\ce{CH3COOH + H2O <=> CH3COO- + H3O+}$, la \cle{loi d'action de masse} s'écrit :
\[ K_a = \frac{[\ce{CH3COO-}][\ce{H3O+}]}{[\ce{CH3COOH}]} \]
(l'eau, solvant, n'apparaît pas dans l'expression).

\item \textbf{Vrai.} À la \cle{demi-équivalence}, la moitié de l'acide a été transformée : $n(\ce{AH}) = n(\ce{A-})$, donc $\dfrac{[\ce{A-}]}{[\ce{AH}]} = 1$ et $\mathrm{pH} = pK_a + \log 1 = pK_a$.
\end{enumerate}
\end{corrige}

\begin{corrige}[Exercice 3 — Texte à trous]
Un \cle{acide fort} réagit avec l'eau de manière \textbf{totale} ; le pH de sa solution de concentration $C$ vaut alors \textbf{$-\log C$}. Un \cle{acide faible} ne s'ionise que de façon \textbf{partielle} : on caractérise cette réaction limitée par son \cle{taux d'avancement}, inférieur à $1$. La force d'un \cle{couple acide/base} est repérée par la valeur de son \textbf{$pK_a$} : plus celui-ci est petit, plus l'acide est fort et plus sa \cle{base conjuguée} est faible. Dans une solution dont le pH est supérieur au $pK_a$ d'un couple, c'est la forme basique qui \textbf{prédomine}.

\medskip
\textbf{Point méthode :} le mot \emph{ionisation} était un leurre : c'est bien le \emph{taux d'avancement} (ou \cle{coefficient d'ionisation} $\alpha$) qui quantifie le caractère limité de la réaction.
\end{corrige}

\begin{corrige}[Exercice 4 — Définitions et couples]
\begin{enumerate}
\item \textbf{Acide fort :} espèce chimique qui cède totalement son proton à l'eau ; sa réaction avec l'eau est totale : $\ce{AH + H2O -> A- + H3O+}$ ($\tau = 1$).

\textbf{Base faible :} espèce chimique qui capte partiellement les protons de l'eau ; sa réaction avec l'eau est un équilibre limité : $\ce{B + H2O <=> BH+ + HO-}$ ($\tau < 1$).

\textbf{Constante d'acidité $K_a$ :} constante d'équilibre associée à la réaction d'un acide \ce{AH} avec l'eau : $K_a = \dfrac{[\ce{A-}][\ce{H3O+}]}{[\ce{AH}]}$. Elle ne dépend que de la température.

\textbf{Coefficient d'ionisation $\alpha$ :} fraction de l'acide (ou de la base) qui a réagi avec l'eau à l'équilibre : $\alpha = \dfrac{x_f}{C} = \dfrac{[\ce{A-}]_f}{C}$.

\item Dans $\ce{HCOOH + NH3 <=> HCOO- + NH4+}$ :
\begin{itemize}
\item Couple 1 : \ce{HCOOH}/\ce{HCOO-} — forme acide : \ce{HCOOH} (cède un proton), forme basique : \ce{HCOO-}.
\item Couple 2 : \ce{NH4+}/\ce{NH3} — forme acide : \ce{NH4+}, forme basique : \ce{NH3} (capte un proton).
\end{itemize}

\item L'ion ammonium est un acide faible :
\[ \ce{NH4+ + H2O <=> NH3 + H3O+} \]
\[ K_a = \frac{[\ce{NH3}][\ce{H3O+}]}{[\ce{NH4+}]} \]
\end{enumerate}
\end{corrige}

\begin{corrige}[Exercice 5 — Classement de trois solutions]
\begin{enumerate}
\item \textbf{pH théoriques pour $C = \SI{1,0e-2}{mol.L^{-1}}$ :}
\begin{itemize}
\item Acide fort : $\mathrm{pH} = -\log C = -\log(1{,}0\times10^{-2}) = 2{,}0$.
\item Base forte : $[\ce{HO-}] = C$ ; $\mathrm{pOH} = 2{,}0$ ; $\mathrm{pH} = 14 - 2{,}0 = 12{,}0$.
\end{itemize}

\item \textbf{Comparaisons :}
\begin{center}
\begin{tabularx}{\linewidth}{|l|X|X|X|}
\hline
\rowcolor{popblueL} \textbf{Solution} & \textbf{pH mesuré} & \textbf{pH théorique} & \textbf{Conclusion} \\
\hline
\ce{HCl} & $2{,}0$ & $2{,}0$ (acide fort) & mesuré $=$ théorique \\
\hline
\ce{CH3COOH} & $3{,}4$ & $2{,}0$ (acide fort) & mesuré $>$ théorique \\
\hline
\ce{NH3} & $10{,}6$ & $12{,}0$ (base forte) & mesuré $<$ théorique \\
\hline
\end{tabularx}
\end{center}

\item \textbf{Classement :}
\begin{itemize}
\item \ce{HCl} : $\mathrm{pH}_{\text{mesuré}} = -\log C$ $\Rightarrow$ ionisation totale $\Rightarrow$ \textbf{acide fort}.
\item \ce{CH3COOH} : $\mathrm{pH}_{\text{mesuré}} = 3{,}4 > -\log C = 2{,}0$ $\Rightarrow$ $[\ce{H3O+}] < C$, ionisation partielle $\Rightarrow$ \textbf{acide faible}.
\item \ce{NH3} : $\mathrm{pH}_{\text{mesuré}} = 10{,}6 < 14 + \log C = 12{,}0$ $\Rightarrow$ $[\ce{HO-}] < C$, réaction partielle avec l'eau $\Rightarrow$ \textbf{base faible}.
\end{itemize}
\end{enumerate}
\end{corrige}

\begin{corrige}[Exercice 6 — Coefficient d'ionisation de l'acide méthanoïque]
\begin{enumerate}
\item Équation de la réaction avec l'eau :
\[ \ce{HCOOH + H2O <=> HCOO- + H3O+} \]

\item \cle{Tableau d'avancement} pour \SI{1}{L} de solution (concentrations en \SI{}{mol.L^{-1}}) :
\begin{center}
\begin{tabularx}{\linewidth}{|l|X|X|X|}
\hline
\rowcolor{popblueL} \textbf{Espèce} & \ce{HCOOH} & \ce{HCOO-} & \ce{H3O+} \\
\hline
\textbf{État initial} & $5{,}0\times10^{-2}$ & $0$ & $0$ (négligé) \\
\hline
\textbf{État final} & $5{,}0\times10^{-2} - x_f$ & $x_f$ & $x_f$ \\
\hline
\end{tabularx}
\end{center}

\item Concentration en ions oxonium :
\[ [\ce{H3O+}] = 10^{-\mathrm{pH}} = 10^{-2{,}6} = \SI{2,5e-3}{mol.L^{-1}} \]
D'après le tableau d'avancement : $x_f = [\ce{H3O+}]_f = \SI{2,5e-3}{mol.L^{-1}}$ (pour 1 L, $x_f = \SI{2,5e-3}{mol}$).

\item Taux d'avancement final :
\[ \tau = \alpha = \frac{x_f}{C} = \frac{2{,}5\times10^{-3}}{5{,}0\times10^{-2}} = 5{,}0\times10^{-2} = 5{,}0\,\% \]

\item \textbf{Conclusion :} $\tau = 5{,}0\,\% \ll 100\,\%$ : la réaction est très limitée. L'acide méthanoïque est un \textbf{acide faible}. Vérification : $-\log C = 1{,}3 < \mathrm{pH} = 2{,}6$, ce qui confirme l'ionisation partielle.
\end{enumerate}
\end{corrige}

\begin{corrige}[Exercice 7 — Classement de couples et prévision d'une réaction]
\begin{enumerate}
\item \textbf{Acides par force croissante} (plus $pK_a$ est grand, plus l'acide est faible) :
\[ \ce{NH4+}\ (9{,}2) \;<\; \ce{CH3COOH}\ (4{,}8) \;<\; \ce{HF}\ (3{,}2) \]
\textbf{Bases conjuguées par force croissante} (plus l'acide est fort, plus sa base conjuguée est faible ; ordre inverse) :
\[ \ce{F-} \;<\; \ce{CH3COO-} \;<\; \ce{NH3} \]

\item Axe des $pK_a$ :
\begin{popfigure}
\begin{tikzpicture}[scale=0.9, every node/.style={font=\small}]
\draw[popline, thick, ->] (0,0) -- (14.5,0) node[right] {$pK_a$};
\foreach \x/\l in {0/0, 2/2, 4/4, 6/6, 8/8, 10/10, 12/12, 14/14} {
    \draw (\x,0.15) -- (\x,-0.15) node[below] {\l};
}
% Couples
\draw[popblue, very thick] (3.2,0.2) -- (3.2,1.2) node[above, popblue] {\ce{HF}/\ce{F-} (3,2)};
\draw[poporange, very thick] (4.8,-0.2) -- (4.8,-1.2) node[below, poporange] {\ce{CH3COOH}/\ce{CH3COO-} (4,8)};
\draw[popgreen, very thick] (9.2,0.2) -- (9.2,1.2) node[above, popgreen] {\ce{NH4+}/\ce{NH3} (9,2)};
% Flèches indicatives
\node[popmuted] at (1.5,1.8) {Acide plus fort $\leftarrow$};
\draw[popmuted, ->] (3.2,1.8) -- (0.5,1.8);
\node[popmuted] at (12.5,1.8) {$\rightarrow$ Base conjuguée plus forte};
\draw[popmuted, ->] (10.5,1.8) -- (13.5,1.8);
\end{tikzpicture}
\legende{Positionnement des trois couples sur l'axe des $pK_a$.}
\end{popfigure}

\item \textbf{Règle du gamma :} la réaction envisagée met en jeu l'acide \ce{HF} ($pK_{a1} = 3{,}2$) et la base \ce{NH3}, dont l'acide conjugué est \ce{NH4+} ($pK_{a2} = 9{,}2$). Comme $pK_{a1} < pK_{a2}$, l'acide le plus fort (\ce{HF}) cède son proton à la base la plus forte (\ce{NH3}) : la réaction est \textbf{possible et favorisée vers la droite}.

\item Équation et constante d'équilibre :
\[ \ce{HF + NH3 <=> F- + NH4+} \]
\[ K = 10^{pK_{a2} - pK_{a1}} = 10^{9{,}2 - 3{,}2} = 10^{6} \]
$K = 10^{6} > 10^{4}$ : la réaction est \textbf{quasi totale}.
\end{enumerate}
\end{corrige}

\begin{corrige}[Exercice 8 — Relation pH--$pK_a$ et ammoniaque ménager]
\begin{enumerate}
\item \textbf{Établissement de la relation.} Pour le couple \ce{AH}/\ce{A-} :
\[ K_a = \frac{[\ce{A-}][\ce{H3O+}]}{[\ce{AH}]} \quad \Rightarrow \quad [\ce{H3O+}] = K_a \times \frac{[\ce{AH}]}{[\ce{A-}]} \]
On applique $-\log$ aux deux membres :
\[ -\log[\ce{H3O+}] = -\log K_a - \log\frac{[\ce{AH}]}{[\ce{A-}]} \]
Or $-\log\dfrac{[\ce{AH}]}{[\ce{A-}]} = +\log\dfrac{[\ce{A-}]}{[\ce{AH}]}$, d'où la \cle{relation de Henderson-Hasselbalch} :
\[ \boxed{\mathrm{pH} = pK_a + \log\frac{[\ce{A-}]}{[\ce{AH}]}} \]

\item \textbf{a)} $\mathrm{pH} = 11{,}0 > pK_a = 9{,}2$ : la solution se situe dans le \cle{domaine de prédominance} de la forme basique. L'espèce prédominante est donc \textbf{\ce{NH3}} (ammoniac), ce qui est cohérent avec l'odeur piquante de l'ammoniaque ménager.

\textbf{b)} Application numérique :
\[ \log\frac{[\ce{NH3}]}{[\ce{NH4+}]} = \mathrm{pH} - pK_a = 11{,}0 - 9{,}2 = 1{,}8 \]
\[ \frac{[\ce{NH3}]}{[\ce{NH4+}]} = 10^{1{,}8} \approx 63 \]
Il y a environ $63$ fois plus d'ammoniac que d'ions ammonium : la prédominance de \ce{NH3} est écrasante.

\item Les deux espèces sont en concentrations égales lorsque $\dfrac{[\ce{NH3}]}{[\ce{NH4+}]} = 1$, c'est-à-dire $\log 1 = 0$, donc :
\[ \mathrm{pH} = pK_a = 9{,}2 \]
\end{enumerate}
\end{corrige}

\begin{corrige}[Exercice 9 — Type Bac : degré d'acidité d'un vinaigre de vin de palme]
\textbf{Barème indicatif (sur 10 points) :} Q1 : 1,5 pt ; Q2 : 1 pt ; Q3 : 1,5 pt ; Q4 : 2,5 pts ; Q5 : 1 pt ; Q6 : 1 pt ; Q7 : 1,5 pt.

\begin{enumerate}
\item \textbf{Équation de la réaction de dosage :}
\[ \ce{CH3COOH + HO- -> CH3COO- + H2O} \]
Constante d'équilibre (acide 1 : \ce{CH3COOH}, $pK_{a1} = 4{,}8$ ; acide 2 : \ce{H2O}, $pK_{a2} = pK_e = 14$) :
\[ K = 10^{pK_{a2} - pK_{a1}} = 10^{14 - 4{,}8} = 10^{9{,}2} \approx 1{,}6\times10^{9} \]
$K \gg 10^{4}$ : la réaction est \textbf{quasi totale}. Étant de plus rapide et unique, elle est \textbf{apte au dosage}.

\item \textbf{Schéma du dispositif :} burette graduée contenant la soude ($C_B$) au-dessus d'un bécher contenant $V_A$ de solution $S$ et un barreau aimanté sur agitateur magnétique ; la sonde du pH-mètre (préalablement étalonné) plonge dans le bécher.

\item \textbf{Équivalence :} état du mélange où les réactifs ont été introduits dans les proportions stœchiométriques de la réaction de dosage : $n(\ce{CH3COOH})_{\text{initial}} = n(\ce{HO-})_{\text{versé}}$.
\[ C_A V_A = C_B V_E \quad \Rightarrow \quad C_A = \frac{C_B V_E}{V_A} = \frac{0{,}10 \times 12{,}0}{10{,}0} = \SI{0,12}{mol.L^{-1}} \]

\item \textbf{Concentration du vinaigre pur} (dilution 10 fois) :
\[ C_0 = 10 \times C_A = \SI{1,2}{mol.L^{-1}} \]
Masse d'acide éthanoïque dans \SI{1}{L} de vinaigre :
\[ m = C_0 \times M = 1{,}2 \times 60 = \SI{72}{g} \]
Masse d'\SI{1}{L} de vinaigre : $m_V = \rho \times V = 1{,}0 \times 1000 = \SI{1000}{g}$.
Le \cle{degré d'acidité} est la masse d'acide pour \SI{100}{g} de vinaigre :
\[ \text{Degré} = \frac{72}{1000} \times 100 = \boxed{7{,}2^{\circ}} \]
\textbf{Commentaire :} ce vinaigre artisanal de vin de palme ($7{,}2^{\circ}$) est plus acide qu'un vinaigre commercial standard ($6^{\circ}$) : sa fermentation a été bien conduite, voire légèrement prolongée.

\item À la \textbf{demi-équivalence} ($V_B = V_E/2 = \SI{6,0}{mL}$), la moitié de l'acide a été transformée : $[\ce{CH3COOH}] = [\ce{CH3COO-}]$, donc :
\[ \mathrm{pH} = pK_a + \log 1 = pK_a = \boxed{4{,}8} \]

\item À l'équivalence, la solution contient l'ion \textbf{éthanoate \ce{CH3COO-}}, base conjuguée d'un acide faible : c'est une \textbf{base faible} qui réagit avec l'eau :
\[ \ce{CH3COO- + H2O <=> CH3COOH + HO-} \]
La production d'ions \ce{HO-} rend la solution basique : $\mathrm{pH}_E > 7$.

\item \textbf{Choix de l'indicateur :} la zone de virage doit encadrer le pH à l'équivalence ($\mathrm{pH}_E \approx 8{,}7$ ici).
\begin{itemize}
\item Hélianthine ($3{,}1$--$4{,}4$) : vire bien avant l'équivalence $\Rightarrow$ rejetée.
\item BBT ($6{,}0$--$7{,}6$) : vire juste avant l'équivalence $\Rightarrow$ erreur systématique, rejeté.
\item \textbf{Phénolphtaléine ($8{,}2$--$10{,}0$) : adaptée}, sa zone de virage contient le saut de pH à l'équivalence.
\end{itemize}
\end{enumerate}

\textbf{Points de vigilance :}
\begin{itemize}
\item Ne pas oublier le facteur de dilution $f = 10$ entre $C_A$ et $C_0$ : oublier donne $0{,}72^{\circ}$, absurde pour un vinaigre.
\item Le degré s'exprime en g d'acide pour \SI{100}{g} de vinaigre, pas en \SI{}{g.L^{-1}} : avec $\rho = \SI{1,0}{g.mL^{-1}}$, degré $= \dfrac{C_0 \times M}{10}$.
\item Les volumes doivent être dans la même unité dans $C_A V_A = C_B V_E$ (ici mL partout : inutile de convertir).
\end{itemize}
\end{corrige}

\begin{corrige}[Exercice 10 — Type Bac : identification d'un acide faible par détermination de son $K_a$]
\textbf{Barème indicatif (sur 10 points) :} Q1 : 1,5 pt ; Q2 : 1 pt ; Q3 : 2 pts ; Q4 : 1,5 pt ; Q5 : 1 pt ; Q6 : 1 pt ; Q7 : 2 pts.

\begin{enumerate}
\item Équation de la réaction :
\[ \ce{AH + H2O <=> A- + H3O+} \]
Tableau d'avancement pour \SI{1}{L} (en \SI{}{mol.L^{-1}}) :
\begin{center}
\begin{tabularx}{\linewidth}{|l|X|X|X|}
\hline
\rowcolor{popblueL} \textbf{Espèce} & \ce{AH} & \ce{A-} & \ce{H3O+} \\
\hline
\textbf{État initial} & $1{,}0\times10^{-1}$ & $0$ & $0$ (négligé) \\
\hline
\textbf{État final} & $1{,}0\times10^{-1} - x_f$ & $x_f$ & $x_f$ \\
\hline
\end{tabularx}
\end{center}

\item Si l'acide était fort : $\mathrm{pH} = -\log C = -\log(1{,}0\times10^{-1}) = 1{,}0$.
Or $\mathrm{pH}_{\text{mesuré}} = 2{,}9 > 1{,}0$, donc $[\ce{H3O+}] = 10^{-2{,}9} < C$ : l'ionisation est partielle, \textbf{\ce{AH} est un acide faible}.

\item Concentrations à l'équilibre :
\[ [\ce{H3O+}] = 10^{-2{,}9} = \SI{1,26e-3}{mol.L^{-1}} \]
\textbf{Électroneutralité :} $[\ce{H3O+}] = [\ce{A-}] + [\ce{HO-}]$. En milieu acide, $[\ce{HO-}] = \dfrac{K_e}{[\ce{H3O+}]} = 10^{-11{,}1} \ll [\ce{H3O+}]$, donc :
\[ [\ce{A-}] \approx [\ce{H3O+}] = \SI{1,26e-3}{mol.L^{-1}} \]
\textbf{Conservation de la matière :} $[\ce{AH}] + [\ce{A-}] = C$, donc :
\[ [\ce{AH}] = 1{,}0\times10^{-1} - 1{,}26\times10^{-3} = \SI{9,87e-2}{mol.L^{-1}} \]

\item Constante d'acidité :
\[ K_a = \frac{[\ce{A-}][\ce{H3O+}]}{[\ce{AH}]} = \frac{(1{,}26\times10^{-3})^2}{9{,}87\times10^{-2}} = \frac{1{,}59\times10^{-6}}{9{,}87\times10^{-2}} = \SI{1,6e-5}{} \]
\[ pK_a = -\log(1{,}6\times10^{-5}) = 4{,}8 \]

\item Comparaison à la table : $pK_a = 4{,}8$ correspond au couple \ce{CH3COOH}/\ce{CH3COO-}. L'acide \ce{AH} est l'\textbf{acide éthanoïque} \ce{CH3COOH}.

\item Coefficient d'ionisation :
\[ \alpha = \frac{x_f}{C} = \frac{1{,}26\times10^{-3}}{1{,}0\times10^{-1}} = 1{,}26\times10^{-2} \approx 1{,}3\,\% \]

\item \textbf{Effet de la dilution :} si l'on dilue la solution, $\alpha$ \textbf{augmente}. Justification : pour un acide faible peu ionisé, $K_a = \dfrac{C\alpha^2}{1-\alpha} \approx C\alpha^2$, d'où $\alpha \approx \sqrt{\dfrac{K_a}{C}}$ (\cle{loi de dilution d'Ostwald}) : quand $C$ diminue, $\alpha$ croît. La dilution déplace l'équilibre $\ce{AH + H2O <=> A- + H3O+}$ vers la droite (sens qui augmente le nombre de particules dissoutes), conformément au principe de Le Chatelier. L'acide faible se rapproche du comportement d'un acide fort à grande dilution.
\end{enumerate}

\textbf{Points de vigilance :}
\begin{itemize}
\item Ne pas écrire $[\ce{AH}] = C$ dans le calcul de $K_a$ sans le justifier : ici $\alpha = 1{,}3\,\%$, l'approximation est licite mais la valeur exacte $C - x_f$ est préférable au Bac.
\item Vérifier la cohérence : $pK_a$ calculé doit être comparé aux valeurs de la table \emph{avec les incertitudes de mesure du pH} en tête.
\end{itemize}
\end{corrige}

\begin{corrige}[Exercice 11 — Type Bac : étude comparative d'un acide fort et d'une base faible]
\textbf{Barème indicatif (sur 12 points) :} Q1 : 1,5 pt ; Q2 : 2 pts ; Q3 : 2 pts ; Q4 : 1 pt ; Q5 : 1,5 pt ; Q6 : 4 pts.

\begin{enumerate}
\item \textbf{\ce{HCl} est un acide fort.} Si l'ionisation est totale : $[\ce{H3O+}] = C_1 = \SI{2,0e-2}{mol.L^{-1}}$, soit :
\[ \mathrm{pH}_{\text{th}} = -\log(2{,}0\times10^{-2}) = 1{,}70 \]
Le pH mesuré ($1{,}7$) est égal au pH théorique : l'ionisation est totale, \ce{HCl} est un acide fort.
\[ \ce{HCl + H2O -> Cl- + H3O+} \]

\item \textbf{\ce{NH3} est une base faible.} Si \ce{NH3} était une base forte : $[\ce{HO-}] = C_2 = \SI{2,0e-2}{mol.L^{-1}}$ :
\[ \mathrm{pOH} = -\log(2{,}0\times10^{-2}) = 1{,}70 \quad \Rightarrow \quad \mathrm{pH}_{\text{th}} = 14 - 1{,}70 = 12{,}30 \]
Or $\mathrm{pH}_{\text{mesuré}} = 10{,}8 < 12{,}30$ : $[\ce{HO-}] < C_2$, la réaction avec l'eau est partielle : \ce{NH3} est une \textbf{base faible}.
\[ \ce{NH3 + H2O <=> NH4+ + HO-} \]

\item \textbf{Concentrations dans $S_2$} ($\mathrm{pH} = 10{,}8$) :
\[ [\ce{H3O+}] = 10^{-10{,}8} = \SI{1,6e-11}{mol.L^{-1}} \]
\[ [\ce{HO-}] = \frac{K_e}{[\ce{H3O+}]} = \frac{10^{-14}}{1{,}6\times10^{-11}} = \SI{6,3e-4}{mol.L^{-1}} \]
Électroneutralité : $[\ce{NH4+}] + [\ce{H3O+}] = [\ce{HO-}]$, avec $[\ce{H3O+}] \ll [\ce{HO-}]$ :
\[ [\ce{NH4+}] \approx [\ce{HO-}] = \SI{6,3e-4}{mol.L^{-1}} \]
Conservation de la matière :
\[ [\ce{NH3}] = C_2 - [\ce{NH4+}] = 2{,}0\times10^{-2} - 6{,}3\times10^{-4} = \SI{1,94e-2}{mol.L^{-1}} \]

\item Taux d'avancement final :
\[ \tau = \frac{[\ce{NH4+}]}{C_2} = \frac{6{,}3\times10^{-4}}{2{,}0\times10^{-2}} = 3{,}2\times10^{-2} \approx 3{,}2\,\% \]
$\tau \ll 100\,\%$ : confirmation quantitative du caractère faible de la base.

\item Constante d'acidité du couple \ce{NH4+}/\ce{NH3} :
\[ K_a = \frac{[\ce{NH3}][\ce{H3O+}]}{[\ce{NH4+}]} = \frac{1{,}94\times10^{-2} \times 1{,}6\times10^{-11}}{6{,}3\times10^{-4}} = \SI{4,9e-10}{} \]
\[ pK_a = -\log(4{,}9\times10^{-10}) = 9{,}3 \]
Valeur proche de la donnée ($9{,}2$) : l'écart de $0{,}1$ unité provient des arrondis sur le pH mesuré (donné au dixième). Le résultat expérimental est \textbf{cohérent}.

\item \textbf{Mélange de $S_1$ et $S_2$ :}

\textbf{a)} On met en présence l'acide le plus fort (\ce{H3O+}) et la base \ce{NH3} :
\[ \ce{H3O+ + NH3 -> NH4+ + H2O} \]
Constante d'équilibre (acide 1 : \ce{H3O+}, $pK_{a1} = 0$ ; acide 2 : \ce{NH4+}, $pK_{a2} = 9{,}2$) :
\[ K = 10^{pK_{a2} - pK_{a1}} = 10^{9{,}2 - 0} = 10^{9{,}2} \approx 1{,}6\times10^{9} \gg 10^{4} \]
La réaction est \textbf{quasi totale} : on peut la traiter comme totale pour le bilan de matière.

\textbf{b)} Quantités initiales :
\[ n(\ce{H3O+}) = C_1 V_1 = 2{,}0\times10^{-2} \times 20{,}0\times10^{-3} = \SI{4,0e-4}{mol} \]
\[ n(\ce{NH3}) = C_2 V_2 = 2{,}0\times10^{-2} \times 20{,}0\times10^{-3} = \SI{4,0e-4}{mol} \]
Les réactifs sont en proportions stœchiométriques : ils sont \textbf{totalement consommés}. Bilan final (couple \ce{NH4+}/\ce{NH3}) :
\[ n_f(\ce{NH4+}) = \SI{4,0e-4}{mol} \quad ; \quad n_f(\ce{NH3}) \approx 0 \]
(\ce{Cl-} est ion spectateur.)

\textbf{c)} Le mélange final est une solution d'ions ammonium \ce{NH4+} (acide faible) de concentration :
\[ [\ce{NH4+}] = \frac{n_f}{V_1 + V_2} = \frac{4{,}0\times10^{-4}}{40{,}0\times10^{-3}} = \SI{1,0e-2}{mol.L^{-1}} \]
Pour un acide faible peu ionisé : $[\ce{H3O+}] \approx \sqrt{K_a \times C}$ :
\[ [\ce{H3O+}] = \sqrt{10^{-9{,}2} \times 1{,}0\times10^{-2}} = \sqrt{10^{-11{,}2}} = 10^{-5{,}6} = \SI{2,5e-6}{mol.L^{-1}} \]
\[ \boxed{\mathrm{pH} = 5{,}6} \]
Le mélange est acide : c'est la solution d'un \textbf{sel d'acide fort et de base faible} (chlorure d'ammonium), dont le caractère acide est un résultat classique à connaître.
\end{enumerate}

\textbf{Points de vigilance :}
\begin{itemize}
\item Question 6 : bien raisonner sur les \textbf{quantités de matière} avant de repasser aux concentrations (le volume final est $V_1 + V_2 = \SI{40,0}{mL}$, pas \SI{20,0}{mL}).
\item Ne pas conclure « pH = 7 car les réactifs sont en proportions stœchiométriques » : c'est faux ici car le produit formé, \ce{NH4+}, est un acide faible. La neutralité ($\mathrm{pH} = 7$) ne vaut que pour le mélange équimolaire acide fort + base forte.
\item Question 5 : utiliser $[\ce{H3O+}]$ (et non $[\ce{HO-}]$) dans l'expression de $K_a$ ; les deux sont liées par $K_e$.
\end{itemize}
\end{corrige}

\newpage

\coursec{Fiche bilan}

\coursec{Force des acides et des bases, couples acide/base}

\begin{popfigure}
\begin{center}
\begin{tikzpicture}[scale=0.92, transform shape]
  % Nœud central
  \node[fill=popdark, text=white, rounded corners=4pt, align=center, font=\bfseries\small, inner sep=5pt] (C) at (0,0) {Force des acides\\ et des bases};
  % Branches
  \node[fill=popblueL, text=popdark, rounded corners=3pt, align=center, font=\scriptsize, inner sep=4pt] (N1) at (-5.6,2.6) {\textbf{Acide fort / Base forte}\\ réaction totale\\ pH $=-\log c$ \; ou \; pH $=14+\log c$};
  \node[fill=poporangeL, text=popdark, rounded corners=3pt, align=center, font=\scriptsize, inner sep=4pt] (N2) at (5.6,2.6) {\textbf{Acide faible / Base faible}\\ réaction limitée (équilibre)\\ $[\ce{H3O+}] < c$};
  \node[fill=popgreenL, text=popdark, rounded corners=3pt, align=center, font=\scriptsize, inner sep=4pt] (N3) at (-6.4,-0.2) {\textbf{Avancement}\\ $\tau = \dfrac{x_f}{x_{\max}} = \dfrac{[\ce{H3O+}]}{c}$};
  \node[fill=poppurpleL, text=popdark, rounded corners=3pt, align=center, font=\scriptsize, inner sep=4pt] (N4) at (6.4,-0.2) {\textbf{Constantes}\\ $K_A$, $pK_A=-\log K_A$\\ $K_B$, $pK_A+pK_B=14$};
  \node[fill=poppinkL, text=popdark, rounded corners=3pt, align=center, font=\scriptsize, inner sep=4pt] (N5) at (-4.6,-3.0) {\textbf{Classification}\\ petit $pK_A$ $\Rightarrow$ acide fort\\ réaction : grand $K_A$ + grand $K_A$};
  \node[fill=popgoldL, text=popdark, rounded corners=3pt, align=center, font=\scriptsize, inner sep=4pt] (N6) at (4.6,-3.0) {\textbf{pH--$pK_A$}\\ pH $=pK_A+\log\dfrac{[\text{base}]}{[\text{acide}]}$\\ domaines de prédominance};
  \node[fill=popcream, text=popdark, rounded corners=3pt, align=center, font=\scriptsize, inner sep=4pt, draw=popline] (N7) at (0,-4.4) {\textbf{Applications}\\ $K_A$ expérimental \;·\; degré du vinaigre};
  % Liaisons
  \draw[popblue, very thick] (C) -- (N1);
  \draw[poporange, very thick] (C) -- (N2);
  \draw[popgreen, very thick] (C) -- (N3);
  \draw[poppurple, very thick] (C) -- (N4);
  \draw[poppink, very thick] (C) -- (N5);
  \draw[popgold, very thick] (C) -- (N6);
  \draw[popmuted, very thick] (C) -- (N7);
\end{tikzpicture}
\end{center}
\legende{Carte mentale de la leçon : les sept idées-force à maîtriser pour le Bac TI.}
\end{popfigure}

\coursec{1. Acides forts et bases fortes}

\begin{definition}[Acide fort, base forte]
Un \cle{acide fort} est un acide qui réagit \emph{totalement} avec l'eau : la réaction de transfert de proton est totale, son avancement final est maximal ($x_f = x_{\max}$). On écrit une flèche simple $\ce{->}$.
Un \cle{base forte} est une base qui réagit \emph{totalement} avec l'eau (ou se dissocie totalement en libérant $\ce{HO-}$), son avancement final est maximal.
\end{definition}

\begin{propriete}[Équations-types des acides et bases forts]
\begin{itemize}
  \item \textbf{Acide chlorhydrique (exemple type) :} $\ce{HCl(aq) + H2O(l) -> H3O+(aq) + Cl-(aq)}$.
  \item \textbf{Acide nitrique :} $\ce{HNO3 + H2O -> H3O+ + NO3-}$.
  \item \textbf{Acide sulfurique (1er proton seulement) :} $\ce{H2SO4 + H2O -> H3O+ + HSO4-}$.
  \item \textbf{Hydroxyde de sodium (base forte type) :} $\ce{NaOH(s) ->[\text{eau}] Na+(aq) + HO-(aq)}$.
  \item \textbf{Hydroxyde de potassium :} $\ce{KOH -> K+ + HO-}$.
\end{itemize}
\end{propriete}

\begin{methode}[Montrer qu'un acide (ou une base) est fort]
\textbf{Protocole} (situation classique du Bac) :
\begin{enumerate}
  \item Préparer une solution de l'acide (ou de la base) de concentration molaire $c$ connue (ex. \SI{1,0e-2}{mol.L^{-1}}).
  \item Mesurer son pH à \SI{25}{\celsius} à l'aide d'un pH-mètre étalonné.
  \item Calculer le pH \emph{théorique} si l'espèce était forte :
        \begin{itemize}
          \item Acide : $\text{pH}_{\text{th}} = -\log c$.
          \item Base : $\text{pH}_{\text{th}} = 14 + \log c$ (à \SI{25}{\celsius}).
        \end{itemize}
  \item \textbf{Comparer} : si $\text{pH}_{\text{mesuré}} \approx \text{pH}_{\text{th}}$ (à $\pm 0,1$ près), l'espèce est forte.
\end{enumerate}
\end{methode}

\begin{exemplebox}[Preuve par le pH : HCl \SI{1,0e-2}{mol.L^{-1}}]
$c = \SI{1,0e-2}{mol.L^{-1}}$. $\text{pH}_{\text{th}} = -\log(1,0 \times 10^{-2}) = 2,00$.
Mesure : $\text{pH}_{\text{mes}} = 2,0$. $\Rightarrow$ \textbf{HCl est un acide fort}.
\end{exemplebox}

\begin{attention}[Pièges fréquents]
\begin{itemize}
  \item \textbf{Concentrations élevées ($c > \SI{e-2}{mol.L^{-1}}$) :} les interactions ioniques faussent la relation $\text{pH} = -\log c$ (on utilise les activités, pas les concentrations). Restez dans le domaine dilué.
  \item \textbf{Acide sulfurique $\ce{H2SO4}$ :} seul le \emph{premier} proton est fort ($\ce{H2SO4 -> H3O+ + HSO4-}$). Le second ($\ce{HSO4- <=> H3O+ + SO4^2-}$, $pK_{A2} \approx 1,9$) est faible ! Ne pas écrire $\ce{H2SO4 -> 2H3O+ + SO4^2-}$.
  \item \textbf{Base forte vs Base concentrée :} "Fort" concerne le \emph{degré de dissociation} ($\tau=1$), "concentré" concerne la \emph{quantité de matière par litre}. Une solution diluée de $\ce{NaOH}$ reste une base forte.
\end{itemize}
\end{attention}

\coursec{2. Acides faibles et bases faibles}

\begin{definition}[Acide faible, base faible]
Un \cle{acide faible} réagit \emph{partiellement} avec l'eau : la réaction s'arrête à un \cle{équilibre chimique}. L'avancement final $x_f$ est \emph{inférieur} à l'avancement maximal $x_{\max}$ ($x_f < x_{\max}$). On écrit une double flèche $\ce{<=>}$.
Une \cle{base faible} réagit partiellement avec l'eau (équilibre).
\end{definition}

\begin{propriete}[Exemples courants au programme]
\begin{center}
\begin{tabularx}{\linewidth}{|l|X|X|}
\hline
\rowcolor{popblueL} \textbf{Espèce} & \textbf{Équation d'équilibre avec l'eau} & \textbf{Couple conjugué} \\
\hline
Acide éthanoïque (acétique) & $\ce{CH3COOH + H2O <=> CH3COO- + H3O+}$ & $\ce{CH3COOH}/\ce{CH3COO-}$ \\
\hline
Ion ammonium & $\ce{NH4+ + H2O <=> NH3 + H3O+}$ & $\ce{NH4+}/\ce{NH3}$ \\
\hline
Ion hydrogénocarbonate & $\ce{HCO3- + H2O <=> CO3^2- + H3O+}$ & $\ce{HCO3-}/\ce{CO3^2-}$ \\
\hline
Eau (acide) & $\ce{H2O + H2O <=> HO- + H3O+}$ & $\ce{H2O}/\ce{HO-}$ \\
\hline
Ion éthanoate (base) & $\ce{CH3COO- + H2O <=> CH3COOH + HO-}$ & $\ce{CH3COOH}/\ce{CH3COO-}$ \\
\hline
Ammoniac (base) & $\ce{NH3 + H2O <=> NH4+ + HO-}$ & $\ce{NH4+}/\ce{NH3}$ \\
\hline
\end{tabularx}
\end{center}
\end{propriete}

\begin{methode}[Montrer qu'un acide (ou une base) est faible]
Même protocole que pour les forts. Si $\text{pH}_{\text{mesuré}} > -\log c$ (pour un acide de concentration $c$), alors $[\ce{H3O+}]_{\text{mes}} < c$ : la réaction n'est pas totale, l'acide est \textbf{faible}.
\end{methode}

\begin{exemplebox}[Acide éthanoïque \SI{1,0e-2}{mol.L^{-1}}]
$c = \SI{1,0e-2}{mol.L^{-1}}$. $\text{pH}_{\text{th}}(\text{fort}) = 2,00$.
Mesure : $\text{pH}_{\text{mes}} = 3,38$.
$[\ce{H3O+}] = 10^{-3,38} = \SI{4,2e-4}{mol.L^{-1}} \ll c$.
$\Rightarrow$ \textbf{Acide éthanoïque est un acide faible}.
\end{exemplebox}

\begin{savaistu}[Contexte camerounais : le vinaigre et le vin de palme]
Le \textbf{vinaigre} (issu de l'acétification du vin de palme ou du bili-bili) contient de l'acide éthanoïque (\SIrange{4}{8}{\%} en masse, soit \SIrange{0,7}{1,3}{mol.L^{-1}}). C'est un acide faible : il ne "brûle" pas l'œsophage comme l'acide chlorhydrique concentré, mais il acidifie suffisamment pour conserver les aliments (pickles, condiments). Le \textbf{vin de palme} frais est légèrement acide (pH $\approx$ 4-5) du fait de la fermentation lactique et acétique naturelle.
\end{savaistu}

\coursec{3. Étude quantitative d'un équilibre : avancement et taux d'avancement}

\begin{definition}[Avancement $x$ et avancement maximal $x_{\max}$]
Pour la réaction d'un acide faible $\ce{AH}$ avec l'eau : $\ce{AH + H2O <=> A- + H3O+}$.
\begin{itemize}
  \item $x$ : quantité de matière (en mol) d'acide $\ce{AH}$ qui a réagi à un instant $t$.
  \item $x_{\max}$ : avancement maximal \emph{théorique} si la réaction était totale. Ici, le réactif limitant est $\ce{AH}$ (l'eau est en immense excès), donc $x_{\max} = n_{\text{initial}}(\ce{AH}) = c \times V$.
\end{itemize}
\end{definition}

\begin{methode}[Tableau d'avancement (incontournable au Bac)]
\begin{center}
\begin{tabularx}{\linewidth}{|c|X|X|X|X|}
\hline
\rowcolor{popblueL} \textbf{État} & \textbf{$\ce{AH}$} & \textbf{$\ce{H2O}$} & \textbf{$\ce{A-}$} & \textbf{$\ce{H3O+}$} \\
\hline
Initial & $n_0 = cV$ & Excès & $0$ & $10^{-7}V$ (négligeable) \\
\hline
Intermédiaire & $n_0 - x$ & Excès & $x$ & $x$ \\
\hline
Final (équilibre) & $n_0 - x_f$ & Excès & $x_f$ & $x_f$ \\
\hline
\end{tabularx}
\end{center}
\emph{Concentrations à l'équilibre (volume $V$ constant) :} $[\ce{AH}] = \dfrac{n_0 - x_f}{V} = c - \dfrac{x_f}{V}$, $[\ce{A-}] = [\ce{H3O+}] = \dfrac{x_f}{V}$.
\end{methode}

\begin{definition}[Taux d'avancement (ou degré d'ionisation) $\tau$]
$\tau = \dfrac{x_f}{x_{\max}} = \dfrac{x_f}{n_0} = \dfrac{[\ce{H3O+}]_{\text{eq}}}{c} = \dfrac{10^{-\text{pH}}}{c}$.
\begin{itemize}
  \item $\tau = 1$ (ou $100\%$) $\Rightarrow$ Acide fort.
  \item $\tau < 1$ (souvent $< 5\%$) $\Rightarrow$ Acide faible.
\end{itemize}
\end{definition}

\begin{exoresolu}[Calcul de $\tau$ pour l'acide éthanoïque]
\textbf{Énoncé :} Une solution d'acide éthanoïque de concentration $c = \SI{1,0e-2}{mol.L^{-1}}$ a un pH mesuré de $3,38$ à \SI{25}{\celsius}. Calculer son taux d'avancement $\tau$.
\textbf{Solution :}
$[\ce{H3O+}] = 10^{-\text{pH}} = 10^{-3,38} = \SI{4,17e-4}{mol.L^{-1}}$.
$\tau = \dfrac{[\ce{H3O+}]}{c} = \dfrac{4,17 \times 10^{-4}}{1,0 \times 10^{-2}} = 0,0417 = \SI{4,17}{\%}$.
$\tau < 1 \Rightarrow$ réaction limitée, acide faible.
\tcblower
\textbf{Conclusion :} Seules \SI{4,2}{\%} des molécules d'acide éthanoïque ont donné un proton à l'eau.
\end{exoresolu}

\coursec{4. Constante d'acidité $K_A$, constante de basicité $K_B$ et $pK_A$}

\begin{definition}[Constante d'acidité $K_A$]
Pour le couple acide/base $\ce{AH}/\ce{A-}$ à l'équilibre : $\ce{AH + H2O <=> A- + H3O+}$.
$K_A = \dfrac{[\ce{A-}]_{\text{eq}} [\ce{H3O+}]_{\text{eq}}}{[\ce{AH}]_{\text{eq}}}$ \quad (activités approchées par concentrations molaires, eau omise car activité $\approx 1$).
$K_A$ est sans dimension (quotient d'activités). Elle ne dépend \emph{que} de la température (valeurs usuelles données à \SI{25}{\celsius}).
\end{definition}

\begin{definition}[p$K_A$]
$pK_A = -\log_{10} K_A$ \quad $\Leftrightarrow$ \quad $K_A = 10^{-pK_A}$.
Plus $K_A$ est grand (plus $pK_A$ est petit), plus l'acide $\ce{AH}$ est fort.
\end{definition}

\begin{propriete}[Constante de basicité $K_B$ de la base conjuguée $\ce{A-}$]
Réaction de la base : $\ce{A- + H2O <=> AH + HO-}$.
$K_B = \dfrac{[\ce{AH}]_{\text{eq}} [\ce{HO-}]_{\text{eq}}}{[\ce{A-}]_{\text{eq}}}$.
\textbf{Relation fondamentale :} $K_A \times K_B = K_e = 10^{-14}$ (à \SI{25}{\celsius}).
$\Rightarrow$ \textbf{$pK_A + pK_B = 14$} (à \SI{25}{\celsius}).
\end{propriete}

\begin{aretenir}[Valeurs de $pK_A$ à connaître (à \SI{25}{\celsius})]
\begin{center}
\begin{tabularx}{\linewidth}{|l|c|c|}
\hline
\rowcolor{popblueL} \textbf{Couple $\ce{AH}/\ce{A-}$} & \textbf{$pK_A$} & \textbf{Force relative} \\
\hline
$\ce{H3O+}/\ce{H2O}$ & $-1,7$ & Acide le plus fort existant en milieu aqueux \\
\hline
$\ce{HCl}/\ce{Cl-}$ & $\approx -7$ & Acide fort (réaction totale) \\
\hline
$\ce{CH3COOH}/\ce{CH3COO-}$ & $4,80$ & Acide faible (référence) \\
\hline
$\ce{NH4+}/\ce{NH3}$ & $9,2$ & Acide très faible \\
\hline
$\ce{HCO3-}/\ce{CO3^2-}$ & $10,3$ & Acide très faible \\
\hline
$\ce{H2O}/\ce{HO-}$ & $14,0$ & Acide le plus faible (base forte $\ce{HO-}$) \\
\hline
\end{tabularx}
\end{center}
\end{aretenir}

\begin{exemplebox}[Calcul de $K_A$ et $K_B$ pour le couple acide éthanoïque/éthanoate]
$pK_A = 4,80 \Rightarrow K_A = 10^{-4,80} = \SI{1,58e-5}{}$.
$pK_B = 14 - 4,80 = 9,20 \Rightarrow K_B = 10^{-9,20} = \SI{6,31e-10}{}$.
$K_A \gg K_B$ : l'acide éthanoïque est bien plus fort que l'éthanoate n'est basique.
\end{exemplebox}

\coursec{5. Classification des couples et prévision des réactions}

\begin{propriete}[Règle de classification]
\begin{itemize}
  \item Plus $K_A$ est \textbf{grand} (plus $pK_A$ est \textbf{petit}), plus l'acide $\ce{AH}$ est \textbf{fort}.
  \item Plus l'acide $\ce{AH}$ est fort, plus sa base conjuguée $\ce{A-}$ est \textbf{faible} (petit $K_B$, grand $pK_B$).
  \item \textbf{Échelle des couples :} On classe les couples par $pK_A$ croissant (ou $K_A$ décroissant).
\end{itemize}
\end{propriete}

\begin{methode}[Prévoir le sens spontané d'une réaction acido-basique entre deux couples]
On met en présence un acide $\ce{AH_1}$ (couple 1 : $\ce{AH_1}/\ce{A_1-}$, $pK_{A1}$) et une base $\ce{A_2-}$ (couple 2 : $\ce{AH_2}/\ce{A_2-}$, $pK_{A2}$).
Réaction possible : $\ce{AH_1 + A_2- <=> A_1- + AH_2}$.
\begin{enumerate}
  \item Identifier les deux couples et leurs $pK_A$.
  \item Calculer $K = \dfrac{K_{A1}}{K_{A2}} = 10^{\,pK_{A2} - pK_{A1}}$.
  \item \textbf{Règle du "gamma" (ou règle des $pK_A$) :} La réaction se fait dans le sens qui met en jeu l'\textbf{acide du couple de plus petit $pK_A$} (le plus fort) et la \textbf{base de l'autre couple}.
  \item Si $pK_{A2} - pK_{A1} > 4$ ($K > 10^4$), la réaction est \textbf{quasi-totale} (flèche $\ce{->}$).
  \item Si $|pK_{A2} - pK_{A1}| < 4$, la réaction est \textbf{limitée} (flèche $\ce{<=>}$).
\end{enumerate}
\end{methode}

\begin{exoresolu}[Prévision : $\ce{CH3COOH}$ + $\ce{NH3}$]
Couple 1 : $\ce{CH3COOH}/\ce{CH3COO-}$ ($pK_{A1} = 4,80$).
Couple 2 : $\ce{NH4+}/\ce{NH3}$ ($pK_{A2} = 9,20$).
Réaction : $\ce{CH3COOH + NH3 <=> CH3COO- + NH4+}$.
$pK_{A2} - pK_{A1} = 9,20 - 4,80 = 4,40 > 4$.
$K = 10^{4,40} \approx \SI{2,5e4}{} \gg 1$.
$\Rightarrow$ \textbf{Réaction quasi-totale} : $\ce{CH3COOH + NH3 -> CH3COO- + NH4+}$.
L'acide le plus fort ($\ce{CH3COOH}$, $pK_A$ plus petit) a réagi avec la base de l'autre couple ($\ce{NH3}$).
\end{exoresolu}

\begin{attention}[Erreur classique : inverser les couples]
Ne pas confondre $pK_A$ de l'acide et $pK_B$ de la base. La règle compare les $pK_A$ des \textbf{deux acides} mis en jeu ($\ce{AH_1}$ et $\ce{AH_2}$). Celui qui a le plus petit $pK_A$ "gagne" : il donne son proton.
\end{attention}

\coursec{6. Relation pH–pKa et domaines de prédominance}

\begin{propriete}[Relation de Henderson-Hasselbalch]
À partir de $K_A = \dfrac{[\ce{A-}][\ce{H3O+}]}{[\ce{AH}]}$, on tire :
$[\ce{H3O+}] = K_A \dfrac{[\ce{AH}]}{[\ce{A-}]}$.
En prenant $-\log$ :
\[ \text{pH} = pK_A + \log \dfrac{[\ce{A-}]}{[\ce{AH}]} = pK_A + \log \dfrac{[\text{base}]}{[\text{acide}]} \]
Cette relation est valide pour \textbf{toute solution contenant les deux formes du couple} (mélange acide/base, point d'équivalence d'un dosage, solution d'un sel d'acide faible...).
\end{propriete}

\begin{definition}[Domaines de prédominance]
\begin{itemize}
  \item Si $\text{pH} < pK_A$ : $\log \dfrac{[\text{base}]}{[\text{acide}]} < 0 \Rightarrow [\text{acide}] > [\text{base}]$. \textbf{Forme acide $\ce{AH}$ prédominante}.
  \item Si $\text{pH} > pK_A$ : $\log \dfrac{[\text{base}]}{[\text{acide}]} > 0 \Rightarrow [\text{base}] > [\text{acide}]$. \textbf{Forme basique $\ce{A-}$ prédominante}.
  \item Si $\text{pH} = pK_A$ : $[\ce{AH}] = [\ce{A-}]$. \textbf{Formes équimolaires} (frontière).
\end{itemize}
\end{definition}

\begin{popfigure}
\begin{center}
\begin{tikzpicture}[scale=1.1]
\begin{axis}[
    width=\linewidth, height=6cm,
    axis lines=middle,
    xlabel={pH}, ylabel={Fraction molaire},
    xmin=0, xmax=14, ymin=-0.05, ymax=1.05,
    xtick={0,2,4,4.8,7,9.2,10,12,14},
    xticklabels={0,2,4,4,8,7,9,2,10,12,14},
    ytick={0,0.5,1},
    legend style={at={(0.5,1.15)},anchor=north,legend columns=2},
    domain=0:14, samples=200,
    clip=false
]
% pKa = 4.8 (Acide éthanoïque)
\addplot[popblue, thick] {1/(1+10^(x-4.8))}; % AH
\addplot[poporange, thick] {1/(1+10^(4.8-x))}; % A-
% Ligne verticale pKa
\draw[popmuted, dashed, thick] (axis cs:4.8,0) -- (axis cs:4.8,1.1) node[above, font=\tiny, popmuted] {$pH = pK_A$};
% Zones
\node[popblue, font=\small\bfseries] at (axis cs:2,0.85) {$\ce{AH}$ prédomine};
\node[poporange, font=\small\bfseries] at (axis cs:10,0.85) {$\ce{A-}$ prédomine};
\legend{$\ce{AH}$ (acide), $\ce{A-}$ (base)}
\end{axis}
\end{tikzpicture}
\end{center}
\legende{Diagramme de prédominance du couple $\ce{CH3COOH}/\ce{CH3COO-}$ ($pK_A = 4,80$).}
\end{popfigure}

\begin{exemplebox}[Application : pH d'un mélange tampon]
Mélange : $V_{\ce{CH3COOH}} = \SI{100}{mL}$, $c = \SI{0,10}{mol.L^{-1}}$ + $V_{\ce{CH3COONa}} = \SI{50}{mL}$, $c = \SI{0,10}{mol.L^{-1}}$.
$n_{\ce{AH}} = \SI{10}{mmol}$, $n_{\ce{A-}} = \SI{5}{mmol}$. Volume total = \SI{150}{mL}.
$[\ce{AH}] = 10/150$, $[\ce{A-}] = 5/150$. Rapport $[\ce{A-}]/[\ce{AH}] = 0,5$.
$\text{pH} = 4,80 + \log(0,5) = 4,80 - 0,30 = 4,50$.
$\text{pH} < pK_A \Rightarrow$ forme acide prédominante (cohérent : $n_{\ce{AH}} > n_{\ce{A-}}$).
\end{exemplebox}

\coursec{7. Applications : détermination expérimentale de $K_A$ et degré d'acidité du vinaigre}

\courssub{7.1 Détermination expérimentale de $K_A$ d'un acide faible}

\begin{methode}[Protocole : Mesure de pH d'une solution de concentration connue]
\begin{enumerate}
  \item Préparer une solution de l'acide faible $\ce{AH}$ de concentration $c$ connue (ex. \SI{1,0e-2}{mol.L^{-1}}), température \SI{25}{\celsius}.
  \item Mesurer le pH à l'équilibre (pH-mètre étalonné).
  \item Calculer $[\ce{H3O+}] = 10^{-\text{pH}}$.
  \item Au tableau d'avancement : $[\ce{A-}] = [\ce{H3O+}]$, $[\ce{AH}] = c - [\ce{H3O+}]$.
  \item Calculer $K_A = \dfrac{[\ce{H3O+}]^2}{c - [\ce{H3O+}]}$.
  \item En déduire $pK_A = -\log K_A$.
\end{enumerate}
\end{methode}

\begin{attention}[Approximation $\tau < 5\%$]
Si $\tau = \dfrac{[\ce{H3O+}]}{c} < 0,05$, alors $c - [\ce{H3O+}] \approx c$.
Formule simplifiée : $K_A \approx \dfrac{[\ce{H3O+}]^2}{c} = \dfrac{10^{-2\text{pH}}}{c}$.
\textbf{Vérifiez toujours la condition $\tau < 5\%$ avant d'utiliser l'approximation.}
\end{attention}

\courssub{7.2 Détermination du degré d'acidité du vinaigre (Dosage)}

\begin{definition}[Degré d'acidité (degré acétique)]
Le \cle{degré d'acidité} (noté $d$ ou degré) d'un vinaigre est la masse d'acide éthanoïque ($\ce{CH3COOH}$, $M = \SI{60,05}{g.mol^{-1}}$) contenue dans \SI{100}{g} (ou \SI{100}{mL} si $\rho \approx 1$) de vinaigre.
$1^{\circ} = \SI{1}{g}$ d'acide éthanoïque pour \SI{100}{g} de vinaigre.
C'est un \%\text{m/m} (ou \%\text{m/V} si densité 1).
\end{definition}

\begin{methode}[Dosage du vinaigre par une solution de soude de concentration connue]
\begin{enumerate}
  \item Pipeter un volume $V_0$ de vinaigre (ex. \SI{10,0}{mL}) dans un erlenmeyer. Ajouter quelques gouttes d'indicateur (phénolphtaléïne : incolore $\to$ rose).
  \item Titrer avec une solution de $\ce{NaOH}$ de concentration $c_B$ connue (ex. \SI{0,100}{mol.L^{-1}}) jusqu'au virage (équivalence).
  \item Noter le volume versé $V_B$ (équivalence).
  \item \textbf{Calculs :}
  \begin{align*}
    n_{\ce{NaOH}} &= c_B \times V_B \quad (\text{en mol}) \\
    n_{\ce{CH3COOH}} &= n_{\ce{NaOH}} \quad (\text{stoïchiométrie 1/1}) \\
    m_{\ce{CH3COOH}} &= n_{\ce{CH3COOH}} \times M_{\ce{CH3COOH}} \quad (\text{en g}) \\
    \text{Degré } d &= \dfrac{m_{\ce{CH3COOH}}}{m_{\text{vinaigre}}} \times 100 \quad (\text{ou } \dfrac{m_{\ce{CH3COOH}}}{V_0 \times \rho} \times 100)
  \end{align*}
  Si $\rho_{\text{vinaigre}} \approx \SI{1,01}{g.mL^{-1}}$, on approxime souvent $m_{\text{vinaigre}} \approx V_0$ (en g) pour le degré.
\end{enumerate}
\end{methode}

\begin{exoresolu}[Dosage d'un vinaigre commercial]
\textbf{Énoncé :} On dose \SI{10,0}{mL} de vinaigre ($\rho = \SI{1,01}{g.mL^{-1}}$) par une solution de $\ce{NaOH}$ \SI{0,100}{mol.L^{-1}}. Volume à l'équivalence : $V_B = \SI{15,2}{mL}$. Calculer le degré d'acidité.
\textbf{Solution :}
$n_{\ce{NaOH}} = \SI{0,100}{mol.L^{-1}} \times \SI{15,2e-3}{L} = \SI{1,52e-3}{mol}$.
$n_{\ce{CH3COOH}} = \SI{1,52e-3}{mol}$.
$m_{\ce{CH3COOH}} = \SI{1,52e-3}{mol} \times \SI{60,05}{g.mol^{-1}} = \SI{0,0913}{g}$.
Masse de vinaigre dosé $m_0 = \SI{10,0}{mL} \times \SI{1,01}{g.mL^{-1}} = \SI{10,1}{g}$.
$d = \dfrac{0,0913}{10,1} \times 100 = \SI{0,904}{\%} \approx \textbf{0,90 degré}$.
\tcblower
\textbf{Interprétation :} Ce vinaigre est peu concentré (vinaigre "faible" ou dilué). Un vinaigre standard culinaire titre souvent 5 à 8 degrés.
\end{exoresolu}

\begin{savaistu}[Le vinaigre de palme au Cameroun]
Le vinaigre de palme, obtenu par acétification naturelle du vin de palme (bactéries acétiques \emph{Acetobacter} transformant l'éthanol en acide éthanoïque à l'air libre), titre traditionnellement entre \textbf{4 et 6 degrés}. Sa saveur fruitée unique en fait un ingrédient clé de la cuisine camerounaise (sauces, marinades). Le contrôle de son degré d'acidité par dosage est une application industrielle directe de cette leçon (qualité SONARA, brasseries, PME agroalimentaires).
\end{savaistu}

\begin{bilan}[L'essentiel de la leçon — Version complète pour le Bac]
\textbf{Définitions clés}
\begin{itemize}\setlength{\itemsep}{1pt}
  \item \cle{Acide fort} / \cle{Base forte} : réaction totale avec l'eau ($\tau = 1$, flèche $\ce{->}$). Ex. : $\ce{HCl}$, $\ce{NaOH}$.
  \item \cle{Acide faible} / \cle{Base faible} : réaction limitée, équilibre ($\tau < 1$, flèche $\ce{<=>}$). Ex. : $\ce{CH3COOH}$, $\ce{NH3}$.
  \item \cle{Taux d'avancement} $\tau = \dfrac{x_f}{x_{\max}} = \dfrac{[\ce{H3O+}]}{c} = \dfrac{10^{-\text{pH}}}{c}$.
  \item \cle{Constante d'acidité} $K_A = \dfrac{[\text{base}][\ce{H3O+}]}{[\text{acide}]}$ (équilibre) ; $pK_A = -\log K_A$.
  \item \cle{Couple conjugué} $\ce{AH}/\ce{A-}$ : $pK_A + pK_B = 14$ (à \SI{25}{\celsius}).
\end{itemize}

\textbf{Équations-types à connaître par cœur}
\begin{itemize}\setlength{\itemsep}{1pt}
  \item Acide fort : $\ce{HCl + H2O -> H3O+ + Cl-}$.
  \item Base forte : $\ce{NaOH -> Na+ + HO-}$.
  \item Acide faible : $\ce{CH3COOH + H2O <=> CH3COO- + H3O+}$.
  \item Base faible : $\ce{CH3COO- + H2O <=> CH3COOH + HO-}$.
  \item Réaction entre couples : $\ce{AH1 + A2- <=> A1- + AH2}$.
\end{itemize}

\textbf{Formulaire de poche}
\begin{center}
\begin{tabularx}{\linewidth}{|l|X|X|}
\hline
\rowcolor{popblueL} \textbf{Grandeur} & \textbf{Relation} & \textbf{Condition / Note} \\
\hline
pH acide fort & $\text{pH} = -\log c$ & $c \le \SI{e-2}{mol.L^{-1}}$, \SI{25}{\celsius} \\
\hline
pH base forte & $\text{pH} = 14 + \log c$ & $c \le \SI{e-2}{mol.L^{-1}}$, \SI{25}{\celsius} \\
\hline
$K_A$ & $K_A = \dfrac{[\text{base}][\ce{H3O+}]}{[\text{acide}]}$ & À l'équilibre \\
\hline
$pK_A$ & $pK_A = -\log K_A$ ; $K_A = 10^{-pK_A}$ & --- \\
\hline
Taux d'ionisation & $\tau = \dfrac{10^{-\text{pH}}}{c}$ & Acide faible \\
\hline
$K_A$ exact & $K_A = \dfrac{10^{-2\text{pH}}}{c - 10^{-\text{pH}}}$ & Sans approximation \\
\hline
$K_A$ approx. & $K_A \approx \dfrac{10^{-2\text{pH}}}{c}$ & Si $\tau < 5\%$ \\
\hline
pH mélange (Henderson) & $\text{pH} = pK_A + \log \dfrac{[\text{base}]}{[\text{acide}]}$ & Couple présent \\
\hline
Prédominance & $\text{pH} < pK_A \Rightarrow \ce{AH}$ ; $\text{pH} > pK_A \Rightarrow \ce{A-}$ & Frontière $\text{pH}=pK_A$ \\
\hline
Constante réaction & $K = 10^{\,pK_{A2} - pK_{A1}}$ & $\ce{AH1 + A2- <=> A1- + AH2}$ \\
\hline
Réaction totale ? & $K > 10^4$ ($pK_{A2}-pK_{A1} > 4$) $\Rightarrow$ flèche $\ce{->}$ & Règle du gamma \\
\hline
Degré vinaigre & $d = \dfrac{c_B V_B M_{\ce{CH3COOH}}}{V_0 \rho} \times 100$ & $M=\SI{60,05}{g.mol^{-1}}$, $\rho\approx\SI{1,01}{g.mL^{-1}}$ \\
\hline
\end{tabularx}
\end{center}

\textbf{Méthodes express}
\begin{itemize}\setlength{\itemsep}{1pt}
  \item \textbf{Fort ou faible ?} Comparer pH mesuré à $-\log c$ (acide) ou $14+\log c$ (base). Égal $\Rightarrow$ fort ; écart $\Rightarrow$ faible.
  \item \textbf{Tableau d'avancement} : Toujours écrire l'état initial, intermédiaire ($x$), final ($x_f$). $x_{\max} = n_0(\text{réactif limitant})$.
  \item \textbf{Classer des couples} : $pK_A$ croissant = force acide décroissante. Base conjuguée : force croissante.
  \item \textbf{Prévoir réaction} : Acide du plus petit $pK_A$ + Base de l'autre couple. $\Delta pK_A > 4 \Rightarrow$ totale.
  \item \textbf{$K_A$ expé.} : Mesure pH $\to [\ce{H3O+}] \to K_A = \frac{[\ce{H3O+}]^2}{c-[\ce{H3O+}]}$.
  \item \textbf{Degré vinaigre} : Dosage $\ce{NaOH}$ $\to n_{\ce{CH3COOH}} \to m \to d = \frac{m}{m_{\text{vinaigre}}}\times 100$.
\end{itemize}
\end{bilan}

\begin{aretenir}[Checklist finale avant le Bac TI]
\begin{itemize}\setlength{\itemsep}{2pt}
  \item[$\square$] Je définis sans hésiter acide fort, base forte, acide faible, base faible.
  \item[$\square$] J'écris correctement les équations avec $\ce{->}$ ou $\ce{<=>}$ et les charges.
  \item[$\square$] Je prouve le caractère fort/faible par la comparaison pH mesuré / pH théorique.
  \item[$\square$] Je construis un tableau d'avancement complet (initial, intermédiaire, final).
  \item[$\square$] Je calcule $\tau = 10^{-\text{pH}}/c$ et j'interprète ($\tau=1$ fort, $\tau<1$ faible).
  \item[$\square$] J'identifie les deux couples $\ce{AH}/\ce{A-}$ dans une équation d'équilibre.
  \item[$\square$] J'écris $K_A$, $pK_A$, $K_B$, $pK_B$ et la relation $pK_A+pK_B=14$.
  \item[$\square$] Je classe des couples par $pK_A$ croissant et j'en déduis la force relative.
  \item[$\square$] J'applique la règle du gamma ($K=10^{\Delta pK_A}$) pour prévoir le sens et la totalité.
  \item[$\square$] J'utilise $\text{pH} = pK_A + \log\frac{[\text{base}]}{[\text{acide}]}$ pour calculer un pH ou un rapport.
  \item[$\square$] Je trace et j'exploite un diagramme de prédominance (frontière $\text{pH}=pK_A$).
  \item[$\square$] Je détermine $K_A$ expérimentalement à partir d'une mesure de pH (avec/sans approx).
  \item[$\square$] Je calcule le degré d'acidité d'un vinaigre après dosage par $\ce{NaOH}$.
  \item[$\square$] Je soigne les unités (mol, L, g, \SI{}{\celsius}), les virgules décimales et les chiffres significatifs.
\end{itemize}
\end{aretenir}

\begin{exercice}[Entraînement Bac - Niveau 1 : Identification et force]
On dispose de trois solutions aqueuses \SI{1,0e-2}{mol.L^{-1}} : Solution A (pH = 2,00), Solution B (pH = 3,38), Solution C (pH = 11,95).
\begin{enumerate}
  \item Identifier la nature (acide/base) et la force (forte/faible) de chaque solution.
  \item Écrire l'équation de la réaction prépondérante pour chaque solution.
  \item Calculer le taux d'avancement $\tau$ pour la solution B.
\end{enumerate}
\end{exercice}

\begin{exercice}[Entraînement Bac - Niveau 2 : Constantes et prévision]
Le couple $\ce{NH4+}/\ce{NH3}$ a $pK_A = 9,20$. Le couple $\ce{HCO3-}/\ce{CO3^2-}$ a $pK_A = 10,30$.
\begin{enumerate}
  \item Classer les acides $\ce{NH4+}$ et $\ce{HCO3-}$ par force croissante.
  \item Classer les bases $\ce{NH3}$ et $\ce{CO3^2-}$ par force croissante.
  \item Prédire le sens et la nature (totale/limitée) de la réaction entre $\ce{NH4+}$ et $\ce{CO3^2-}$.
  \item Écrire l'équation de cette réaction avec la flèche appropriée.
\end{enumerate}
\end{exercice}

\begin{exercice}[Entraînement Bac - Niveau 3 : pH-pKa et dosage]
\begin{enumerate}
  \item On prépare un mélange contenant $\ce{CH3COOH}$ et $\ce{CH3COO-}$ tels que $[\ce{CH3COO-}] = 2[\ce{CH3COOH}]$. $pK_A = 4,80$. Calculer le pH. Quelle forme prédomine ?
  \item On dose \SI{20,0}{mL} de vinaigre ($\rho = \SI{1,01}{g.mL^{-1}}$) par une solution de $\ce{NaOH}$ \SI{0,500}{mol.L^{-1}}. L'équivalence est atteinte pour $V_B = \SI{16,4}{mL}$. Calculer le degré d'acidité de ce vinaigre.
\end{enumerate}
\end{exercice}

\begin{corrige}[Exercice Niveau 1]
\begin{enumerate}
  \item \textbf{A :} pH=2,00. pH théorique acide fort $=-\log(10^{-2})=2,00$. $\Rightarrow$ \textbf{Acide fort} (ex. $\ce{HCl}$).\\
  \textbf{B :} pH=3,38 $> 2,00$. $\Rightarrow$ \textbf{Acide faible} (ex. $\ce{CH3COOH}$).\\
  \textbf{C :} pH=11,95. pH théorique base forte $=14+\log(10^{-2})=12,00$. $\approx$ égal. $\Rightarrow$ \textbf{Base forte} (ex. $\ce{NaOH}$).
  \item A : $\ce{H3O+ + Cl-}$ (déjà ions) ou $\ce{HCl + H2O -> H3O+ + Cl-}$.\\
  B : $\ce{CH3COOH + H2O <=> CH3COO- + H3O+}$.\\
  C : $\ce{NaOH -> Na+ + HO-}$.
  \item $\tau_B = \frac{10^{-3,38}}{10^{-2}} = 10^{-1,38} = 0,0417 = \SI{4,17}{\%}$.
\end{enumerate}
\end{corrige}

\begin{corrige}[Exercice Niveau 2]
\begin{enumerate}
  \item Acides : $\ce{NH4+}$ ($pK_A=9,20$) plus fort que $\ce{HCO3-}$ ($pK_A=10,30$). Ordre force croissante : $\ce{HCO3-} < \ce{NH4+}$.
  \item Bases conjuguées : plus l'acide est fort, plus la base est faible. $\ce{NH3}$ (base de $\ce{NH4+}$) plus faible que $\ce{CO3^2-}$ (base de $\ce{HCO3-}$). Ordre force croissante : $\ce{NH3} < \ce{CO3^2-}$.
  \item Réaction : $\ce{NH4+ + CO3^2- <=> NH3 + HCO3-}$.
  Couple 1 : $\ce{NH4+}/\ce{NH3}$ ($pK_{A1}=9,20$). Couple 2 : $\ce{HCO3-}/\ce{CO3^2-}$ ($pK_{A2}=10,30$).
  $\Delta pK_A = pK_{A2} - pK_{A1} = 10,30 - 9,20 = 1,10 < 4$.
  $K = 10^{1,10} \approx 12,6$.
  $\Rightarrow$ \textbf{Réaction limitée} (équilibre), sens direct favorisé ($K>1$).
  \item $\ce{NH4+ + CO3^2- <=> NH3 + HCO3-}$ (double flèche).
\end{enumerate}
\end{corrige}

\begin{corrige}[Exercice Niveau 3]
\begin{enumerate}
  \item $\text{pH} = pK_A + \log\frac{[\ce{CH3COO-}]}{[\ce{CH3COOH}]} = 4,80 + \log 2 = 4,80 + 0,30 = \textbf{5,10}$.
  $\text{pH} > pK_A \Rightarrow$ \textbf{Forme basique $\ce{CH3COO-}$ prédomine} (cohérent : concentration double).
  \item $n_{\ce{NaOH}} = \SI{0,500}{mol.L^{-1}} \times \SI{16,4e-3}{L} = \SI{8,20e-3}{mol}$.
  $n_{\ce{CH3COOH}} = \SI{8,20e-3}{mol}$.
  $m_{\ce{CH3COOH}} = \SI{8,20e-3}{mol} \times \SI{60,05}{g.mol^{-1}} = \SI{0,492}{g}$.
  $m_{\text{vinaigre}} = \SI{20,0}{mL} \times \SI{1,01}{g.mL^{-1}} = \SI{20,2}{g}$.
  $d = \frac{0,492}{20,2} \times 100 = \textbf{2,44 degré}$.
\end{enumerate}
\end{corrige}

\findecours

\end{document}
